\DocumentMetadata{}
\documentclass[%
10pt,%
oneside,
a4paper,%
french,%
notitlepage
]%
{book}%


\usepackage{luatextra}
\usepackage{amsmath,amssymb,marvosym,stmaryrd,calrsfs,minted,setspace,array,caption,fvextra}%
\usepackage[math-style=french]{unicode-math}
\usepackage[math-style=french]{fourier-otf}%
\usepackage[svgnames]{xcolor}
\usepackage[margin=1.5cm,top=1cm,bottom=1.5cm,headheight=5mm,headsep=2mm]{geometry}
\usepackage[unicode,pdfstartview = FitH, colorlinks, linkcolor=blue]{hyperref}

\usepackage{multicol,xspace}%
\usepackage[protrusion=true,expansion=true]{microtype}
\usepackage{babel}%

%\usepackage{minitoc}
\usepackage[Lenny]{fncychap}
\usepackage{animate}
\usepackage[noexec,3d,cachedir=tkz]{luadraw}
\usetikzlibrary{shadings}
\def\version{3.1}

\hypersetup{
  pdfauthor={P. Fradin},
  pdflang={fr},
  %hidelinks,
  pdfcreator={LuaLatex}
}

\title{\textbf{Le paquet \emph{luadraw} (\version)}}

\author{Patrick Fradin}
%\date{}


\onehalfspacing
%\newminted{Lua}{bgcolor=Beige,ignorelexererrors=true,linenos,numbersep=6pt,breaklines,fontsize=\footnotesize}%
\newminted{TeX}{bgcolor=Gray!30,breaklines,fontsize=\footnotesize}%
\newminted{Lua}{bgcolor=Beige,ignorelexererrors=true,breaklines,fontsize=\footnotesize}%

%\columnseprule 0.4pt
\columnsep 25pt

\definecolor{bgcol}{RGB}{238,240,252}% couleur de fond des listings

\newenvironment*{demo}[2][]{%
\gdef\legende{#2}%
\gdef\lab{#1}%
\bgroup
\VerbatimOut{\jobname.tmp}%
}%
{%
\endVerbatimOut%
\egroup%
%\begin{center}
%\begin{tabular}{|m{10cm}|m{7cm}|}
%\hline
%\begin{minipage}{10cm}
%\inputminted[ignorelexererrors=true,breaklines,bgcolor=Beige,linenos,numbersep=6pt,frame=single,fontsize=\footnotesize]{Lua}{\jobname.tmp}%
\inputminted[ignorelexererrors=true,breaklines,bgcolor=Beige,frame=single,fontsize=\footnotesize]{Lua}{\jobname.tmp}%
%\end{minipage}
%&
\begin{minipage}{0.9\textwidth}
%\par\smallskip
\begin{center}
\captionof{figure}{\legende}\label{\lab}%
\input{\jobname.tmp}%
%\end{center}
%\par\smallskip
%\hline
%\end{tabular}
\end{center}
\end{minipage}
}%


\newcommand*{\opt}[1]{\texttt{\textcolor{violet}{#1}}}%option
\newcommand*{\val}[1]{\texttt{\textcolor{blue}{#1}}}%value
\newcommand*{\cmd}[1]{\textbf{#1}}%commande
\newcommand*{\cmdln}[1]{\par\hfil\cmd{#1}\hfil\par}%commande sur une ligne
\newcommand*{\fac}[1]{[#1]}%arguments facultatifs
\newcommand*{\code}[1]{\textbf{\texttt{#1}}}%code
\newcommand*{\codeln}[1]{\par\hfil\code{#1}\hfil\par}%code sur une ligne
%\newcommand*{\argu}[1]{\texttt{\textcolor{Maroon}{<#1>}}}%argument
\newcommand*{\argu}[1]{$\langle$\emph{#1}$\rangle$}%argument
\newcommand*{\varglob}[1]{\texttt{\textcolor{DarkGreen}{#1}}}%variable globale
\newcommand*{\true}{\val{true}\xspace}%
\newcommand*{\false}{\val{false}\xspace}%
\newcommand*{\nil}{\val{nil}\xspace}%
\newcommand*{\drawcmd}{\emph{\textbackslash draw}\xspace}%
\newcommand*{\luadrawenv}{\emph{luadraw}\xspace}%

\begin{document}
\maketitle

\hfil{\Large Dessin 2D et 3D avec Lua et TikZ.}\hfil\par
\hfil{\small \url{https://github.com/pfradin/luadraw}}\hfil\par

\vspace*{2cm} 

\begin{small}
\hfil\textbf{Résumé}\hfil\par
Le paquet \luadrawenv définit l'environnement du même nom, celui-ci permet de créer des graphiques mathématiques en utilisant le langage Lua. Ces graphiques sont dessinés au final par TikZ (et automatiquement sauvegardés), alors pourquoi les faire en Lua ? Parce que celui-ci apporte toute la puissance d'un langage de programmation simple, efficace, capable de faire des calculs, tout en utilisant les possibilités graphiques de TikZ.
\end{small}%

%\include{body-fr/luadraw-doc-saisie.tex}% saisie d'une section ou sous-section seulement

\clearpage%
\setcounter{tocdepth}{3}%
\setcounter{secnumdepth}{2}%

\begin{multicols}{2}
\tableofcontents
\end{multicols}

\listoffigures

\renewcommand{\labelitemi}{$\bullet$}
\renewcommand{\labelitemii}{--}
\renewcommand{\labelitemiii}{$*$}
%\renewcommand{\thechapter}{\Roman{chapter}}
\renewcommand{\thesection}{\Roman{section}~}
\renewcommand{\thesubsection}{\arabic{subsection})}
\renewcommand{\thesubsubsection}{\arabic{subsection}.\arabic{subsubsection}}
\renewcommand{\thefigure}{\arabic{figure}}

\chapter{Dessin 2D}

\begin{center}
\input{tkz/trigo.tkz}%
\par
\captionof{figure}{Un premier exemple : trois sous-figures dans un même graphique}
\end{center}

\section{Introduction}

\subsection{Nouveautés depuis la version 3.0}

Depuis la version 3.0, \textbf{toutes} les données relatives au paquet \emph{luadraw} sont encapsulées dans un même espace de noms, celui-ci s'appelle \textbf{luadraw} (c'est une table). Ceci a pour conséquence que les données sont accessibles uniquement avec la notation pointée, plus précisément :
\begin{enumerate}
    \item la classe \emph{graph} devient \emph{luadraw.graph},
    \item la classe \emph{cpx} (nombres complexes) devient \emph{luadraw.cpx},    
    \item la classe \emph{graph3d} devient \emph{luadraw.graph3d},
    \item la classe \emph{pt3d} (points 3D) devient \emph{luadraw.pt3d},
    \item toutes les variables globales, toutes les fonctions non graphiques sont dans l'espace de nommage \emph{luadraw}, sauf quelques exceptions propres aux nombres complexes qui sont passées dans la classe \emph{cpx}, et quelques exceptions propres aux points 3D qui sont passées dans la classe \emph{pt3d},
    \item par contre il n'y a pas de changement pour les méthodes graphiques puisqu'elles étaient déjà encapsulées dans les classes \emph{graph} et \emph{graph3d}.
\end{enumerate}

Le langage Lua n'autorise pas la factorisation de l'espace de noms, il faut donc réellement utiliser la notation pointée, mais heureusement on peut se créer des raccourcis (ou alias), par exemple :

\begin{Luacode}
local ld = luadraw -- alias sur l'espace de nommage
local cpx, pt3d = ld.cpx, ld.pt3d -- raccourcis pour les classes cpx et pt3d
local Z, M, i = cpx.Z, pt3d.M, cpx.I -- raccourcis vers les fonctions Z et M et le complexe i
local Origin, vecI, vecJ, vecK = pt3d.Origin, pt3d.vecI, pt3d.vecJ, pt3d.vecK
local g = ld.graph3d:new{}
...
\end{Luacode}

Du fait de ce changement majeur, les scripts \emph{luadraw} des versions antérieures \textbf{ne sont plus compilables directement}, ils doivent être adaptés à cette nouvelle version. Mais avec des raccourcis, les changements sont minimes.

\paragraph{Autre nouveauté} : l'environnement \emph{luadraw} fait appel à un environnement \emph{luacode} (et non plus \emph{luacode*}) ce qui permet d'introduire des macros \TeX{} dans le code \emph{Lua}, on peut donc mettre les raccourcis les plus fréquents dans une macro et l'utiliser dans ses codes \emph{luadraw}, par exemple :

\begin{TeXcode}
\documentclass{article}
\usepackage[3d]{luadraw}

\newcommand*{\shortcuts}{%
local ld = luadraw
local cpx, pt3d = ld.cpx, ld.pt3d
local Z, M, i = cpx.Z, pt3d.M, cpx.I
local Origin, vecI, vecJ, vecK = pt3d.Origin, pt3d.vecI, pt3d.vecJ, pt3d.vecK
}%

\begin{document}

\begin{luadraw}{name=test}
\shortcuts
local g = ld.graph3d:new{}
local P = ld.polyreg(0, Z(2,3), 5)
g:Dpolyline(P, true, "red,line width=1.2pt")
local Q = ld.map(pt3d.toPoint3d, P)
P = ld.pyramid(Q, M(0,0,4))
g:Dpoly(P, {color="blue", opacity=0.6})
g:Show()
\end{luadraw}
\end{document}
\end{TeXcode}

L'environnement étoilé \emph{luadraw*} quant à lui, fait appel à un environnement \emph{luacode*}, dans ce cas il n'y a aucune ingérence de \TeX{} dans le code Lua. 

\paragraph{NB} : dans toute la suite de cette documentation nous utiliserons les raccourcis suivants : 

\begin{Luacode}
local ld = luadraw -- alias sur l'espace de nommage
local cpx = ld.cpx -- raccourci pour la classe cpx
local pt3d = ld.pt3d -- raccourci pour la classe pt3d
\end{Luacode}



\subsection{Prérequis}

\begin{itemize}
\item Dans le préambule, il faut déclarer le package \luadrawenv : \verb|\usepackage[options globales]{luadraw}|
\item La compilation se fait avec LuaLatex \textbf{exclusivement}.
\item Les couleurs dans l'environnement \emph{luadraw} sont des chaînes de caractères qui doivent correspondre à des couleurs connues de TikZ. Il est fortement conseillé d'utiliser le package \emph{xcolor} avec l'option \emph{svgnames}.
\end{itemize}

Quelque soient les options globales choisies, ce paquet charge le module \emph{luadraw\_graph2d.lua} qui définit la classe \emph{ld.graph}, et fournit l'environnement \emph{luadraw} qui permet de faire des graphiques en Lua. Pour pouvoir faire des graphiques en 3D il faut charger en plus la classe \emph{ld.graph3d} grâce à l'option globale \opt{3d} :\par
\hfil\verb|\usepackage[3d]{luadraw}|.\hfil

\paragraph{Options globales du paquet} : \opt{noexec}, \opt{3d} et \opt{cachedir=}.

\begin{itemize}
    \item \opt{noexec} : lorsque cette option globale est mentionnée la valeur par défaut de l'option \emph{exec} pour l'environnement \emph{luadraw} sera false (et non plus true).
    \item \opt{3d} : lorsque cette option globale est mentionnée, le module \emph{luadraw\_graph3d.lua} est également chargé. Celui-ci définit en plus la classe \emph{ld.graph3d} (qui s'appuie sur la classe \emph{ld.graph}) pour des dessins en 3D. 
    \item \opt{cachedir=<dossier>} : par défaut les fichiers créés sont enregistrés dans le dossier \emph{\_luadraw} qui est un sous-dossier du dossier courant (contenant le document maître). Ce dossier peut être changé avec l'option \opt{cachedir}, par exemple \opt{cachedir=\{tikz\}}.
\end{itemize}

\noindent\textbf{NB} : dans ce chapitre nous ne parlerons pas de l'option \opt{3d}. Celle-ci fait l'objet du chapitre suivant. Nous ne parlerons donc que de la version 2D.

Lorsqu'un graphique est terminé il est exporté au format TikZ, donc ce paquet charge également le paquet TikZ ainsi que les librairies :
\begin{itemize}
\item\emph{patterns}
\item\emph{plotmarks}
\item\emph{arrows.meta}
\item\emph{decorations.markings}
\end{itemize}

Les graphiques sont créés dans un environnement \luadrawenv, celui-ci appelle \emph{luacode}, c'est donc du \textbf{langage Lua} qu'il faut utiliser dans cet environnement :

\begin{TeXcode}
\begin{luadraw}{ name = <filename>, exec = true/false, auto = true/false }
local ld = luadraw
-- création d'un nouveau graphique en lui donnant un nom local (cette ligne est un commentaire)
local g = ld.graph:new{ window={x1,x2,y1,y2 [,xscale,yscale]}, margin={top,right,bottom,left},
                     size={width,height,ratio}, bg="color", border=true/false }
-- construction du graphique g
    instructions graphiques en langage Lua ...
-- affichage du graphique g et sauvegarde dans le fichier <filename>.tkz
g:Show()
-- ou bien sauvegarde uniquement dans le fichier <filename>.tkz
g:Save()
\end{luadraw}
\end{TeXcode}

\paragraph{Sauvegarde du fichier \emph{*.tkz}} : le graphique est exporté au format TikZ dans un fichier (avec l'extension \emph{tkz}), par défaut celui-ci est sauvegardé dans le dossier \emph{\_luadraw} qui est un sous-dossier du dossier courant (contenant le document maître), mais il est possible d'imposer un chemin vers un autre sous-dossier avec l'option globale \opt{cachedir=\ldots}.

\subsection{Options de l'environnement \luadrawenv}

Celles-ci sont :
\begin{itemize}
    \item \opt{name=\ldots{}} : permet de donner un nom au fichier TikZ produit, on donne un nom sans extension (celle-ci sera automatiquement ajoutée, c'est \emph{.tkz}). Si cette option est omise, alors il y a un nom par défaut, qui est le nom du fichier maître suivi d'un numéro.
    
    \item \opt{exec=true/false} : permet d'exécuter ou non le code Lua compris dans l'environnement. Par défaut cette option vaut true, \textbf{SAUF} si l'option globale \opt{noexec} a été mentionnée dans le préambule avec la déclaration du paquet. Lorsqu'un graphique complexe qui demande beaucoup de calculs est au point, il peut être intéressant de lui ajouter l'option \opt{exec=false}, cela évitera les recalculs de ce même graphique pour les compilations à venir.
    
    \item \opt{auto=true/false} : permet d'inclure ou non automatiquement le fichier TikZ en lieu et place de l'environnement \luadrawenv lorsque l'option \opt{exec} est à \false. Par défaut l'option \opt{auto} vaut \true.
\end{itemize}


\subsection{La classe ld.cpx (complexes)}

Rappel : nous utilisons l'alias \code{local cpx = ld.cpx}.

Elle est automatiquement chargée par le module \emph{luadraw\_graph2d} et donc au chargement du paquet \luadrawenv. Cette classe permet de manipuler les nombres complexes et de faire les calculs habituels. On crée un complexe avec la fonction \cmd{cpx.Z(a, b)} pour \(a+i\times b\), ou bien avec la fonction \cmd{cpx.Zp(r, theta)} pour \(r\times e^{i\theta}\) en coordonnées polaires.

\begin{itemize}
    \item Exemple : \code{local z = cpx.Z(a,b)} va créer le complexe correspondant à \(a+i\times b\) dans la variable $z$. On accède alors aux parties réelle et imaginaire de $z$ comme ceci : \code{z.re} et \code{z.im}. On peut bien sûr créer un raccourci vers cette fonction avec l'instruction \code{local Z = cpx.Z} en début de code.
    
    \item \textbf{Attention} : un nombre réel $x$ n'est pas considéré comme complexe par Lua. Cependant, les fonctions proposées pour les constructions graphiques font la vérification et la conversion réel vers complexe. On peut néanmoins, utiliser $Z(x,0)$ à la place de $x$.
    
    \item Les opérateurs habituels ont été surchargés ce qui permet l'utilisation des symboles habituels, à savoir : +, x, -, /, ainsi que le test d'égalité avec = (ce test d'égalité est fait à \varglob{ld.epsilon} près où \varglob{ld.epsilon} et une variable globale qui vaut $1e-16$ par défaut). Lorsqu'un calcul échoue le résultat renvoyé en principe doit être égal à \nil.
    
    \item À cela s'ajoutent les fonctions suivantes :
  \begin{itemize}
  \item le module : \cmd{cpx.abs(z)},
  \item le module au carré: \cmd{cpx.abs2(z)},
  \item la normalisation : \cmd{cpx.normalize(z)} (renvoie \emph{nil} si $z$ est nul),
  \item la norme 1 : \cmd{cpx.N1(z)},
  \item l'argument principal : \cmd{cpx.arg(z)},
  \item le conjugué : \cmd{cpx.bar(z)},
  \item l'égalité parfaite : \cmd{cpx.equal(z1,z2)},
  \item l'exponentielle complexe : \cmd{cpx.exp(z)},
  \item le cosinus et sinus hyperboliques complexes : \cmd{cpx.cosh(z)} et \cmd{cpx.sinh(z)},
  \item la fonction puissance $a$ : cmd{cpx.pow(z,a)} (le calcul utilise l'argument principal de $z$),
  \item le produit scalaire : \cmd{cpx.dot(z1, z2)}, où les complexes représentent des affixes de vecteurs,
  \item le déterminant : \cmd{cpx.det(z1, z2)},
  \item l'angle orienté (en radians) entre deux vecteurs non nuls : \cmd{cpx.angle(z1, z2)}
  \item l'arrondi : \cmd{cpx.round(z, nb decimales)},
  \item la fonction : \cmd{cpx.isNul(z)} teste si les parties réelle et imaginaire de \emph{z} sont en valeur absolue inférieures à la variable globale \varglob{ld.epsilon} qui vaut $1e-16$ par défaut,
  \item la fonction \cmd{cpx.isComplex(u)} teste si \argu{u} est un complexe ou non et renvoie un booléen (attention : si \argu{u} est réel, la fonction renvoie \false),
  \item la fonction \cmd{cpx.toComplex(x)} teste si \argu{x} est un réel, si c'est le cas, elle renvoie le nombre complexe \code{cpx.Z(x,0)}, si \argu{x} est un nombre complexe la fonction renvoie \argu{x}, et dans les autres cas, elle renvoie \nil,
  \item la fonction \cmd{cpx.isobar(L)} où \argu{L} est une liste de nombres complexes, renvoie l'isobarycentre des points de \argu{L}.
  \end{itemize}
\end{itemize}

On dispose également de la constante \varglob{cpx.I} qui représente l'imaginaire pur $\imath$. 

Exemple :

\begin{Luacode}
local cpx = luadraw.cpx
local i = cpx.I
local A = 2+3*i
\end{Luacode}

Le symbole de multiplication est obligatoire.

\subsection{Afficher une variable dans le terminal}

L'instruction \cmd{ld.whatis(variable \fac{, msg})} affiche dans le terminal lors de la compilation, le type de la \argu{variable} ainsi que son contenu. Les types reconnus sont, les types prédéfinis plus : \emph{complex number}, \emph{list of (complex) numbers}, \emph{list of lists of (complex) numbers}.  L'argument \emph{msg} est une chaîne optionnelle (vide par défaut) qui est affichée avec le type pour repérer la variable dans le terminal.

\subsection{Création d'un graphe}

Comme cela a été vu plus haut, la création se fait dans un environnement \luadrawenv, cette création se fait en nommant le graphique :

\begin{Luacode}
local ld = luadraw
local g = ld.graph:new{ window = {x1,x2,y1,y2 [,xscale,yscale]}, margin = {left,right,top,bottom}, 
                     size = {width,height,ratio}, bg = "color", border =true/false, 
                     bbox = true/false, pictureoptions = "" }
\end{Luacode}

La classe \emph{ld.graph} est définie dans le paquet \luadrawenv. On instancie cette classe en invoquant son constructeur et en donnant un nom (ici c'est \emph{g}), on le fait en local de sorte que le graphique \emph{g} ainsi créé, n'existera plus une fois sorti de l'environnement (sinon \emph{g} resterait en mémoire jusqu'à la fin du document). Voici les options du constructeur \cmd{ld.graph:new{}}:

\begin{itemize}
 \item \opt{window=\{x1,x2,y1,y2 \fac{,xscale,yscale}\}}. Cette option (facultative) définit le pavé de $\mathbf R^2$ dans lequel se fait au graphique (c'est $[x_1;x_2]\times[y_1;y_2]$), ainsi que les échelles sur les axes, qui sont les arguments \argu{xscale} et \argu{yscale}. Ces deux derniers paramètres sont facultatifs et valent $1$ par défaut, ils représentent l'échelle (cm par unité) sur les axes. Par défaut on a \opt{window=\{-5,5,-5,5,1,1\}}.
 
\item \opt{margin=\{left,right,top,bottom\}}. Cette option (facultative) définit des marges autour du graphique en cm. Par défaut on a \opt{margin=\{0.5,0.5,0.5,0.5\}}.

\item \opt{size=\{width,height \fac{,ratio}\}}. Cette option (facultative) permet d'imposer une taille (en cm, marges incluses) pour le graphique, l'argument \argu{ratio} est facultatif et correspond au rapport d'échelle souhaité (\argu{xscale}/\argu{yscale}), un ratio de $1$ donnera un repère orthonormé, et si le ratio n'est pas précisé alors le ratio par défaut est conservé (celui défini par l'option \opt{window}). Un ratio égal à $0$ détermine les échelles de telle sorte que le graphique ait exactement la taille demandée. L'utilisation de ce paramètre va modifier les valeurs de \argu{xscale} et \argu{yscale}. Par défaut la taille est de $11\times11$ (en cm) avec les marges ($10\times10$ sans les marges).

\item \opt{bg="color"}. Cette option (facultative) permet de définir une couleur de fond pour le graphique, cette couleur est une chaîne de caractères représentant une couleur pour TikZ. Par défaut cette chaîne est vide ce qui signifie que le fond ne sera pas peint.

\item \opt{border=true/false}. Cette option (facultative) indique si un cadre doit être dessiné ou non autour du graphique (en incluant les marges). Par défaut ce paramètre vaut \false.

\item \opt{bbox=true/false}. Cette option (facultative) indique si une boundingbox doit être ajoutée au graphique de telle sorte que celui-ci ait la taille souhaitée, tout ce qui en sort est clippé par TikZ. Par défaut ce paramètre vaut \true. Avec la valeur \false il n'y a pas de boundingbox ajoutée, mais tout ce qui sort de la fenêtre 2D, sauf les path, est clippé par \luadrawenv, la taille du graphique peut être plus petite que celle demandée.

\item \opt{pictureoptions=""}. Cette option (facultative) est une chaîne de caractères destinée à contenir des options qui seront passées à l'environnement \emph{tikzpicture} comme ceci:
\begin{TeXcode}
\begin{tikzpicture}[line join=round <,pictureoptions>]
\end{TeXcode}
\end{itemize}


\paragraph{Construction du graphique.}

\begin{itemize}
    \item L'objet instancié (\emph{g} ici dans l'exemple) possède un certain nombre de méthodes permettant de faire du dessin (segments, droites, courbes,\dots). Les instructions de dessins ne sont pas directement envoyées à \TeX, elles sont enregistrées sous forme de chaînes dans une table qui est une propriété de l'objet \emph{g}. C'est la méthode \cmd{g:Show()} qui va envoyer ces instructions à \TeX\ tout en les sauvegardant dans un fichier texte\footnote{Ce fichier contiendra un environnement \emph{tikzpicture}.}. La méthode \cmd{g:Save()} enregistre le graphique dans le fichier désigné par l'option (de l'environnement) \opt{name} mais sans envoyer les instructions à \TeX.
    
    \item On peut faire une sauvegarde du graphique en cours dans un autre fichier avec la méthode: \cmdln{g:Savetofile(<nom de fichier avec extension>)}.
    
    \item On peut réinitialiser un graphique en cours, c'est-à-dire supprimer tous les éléments déjà créés, avec la méthode: \cmdln{g:Cleargraph()}.
    
    \item Le paquet \luadrawenv fournit aussi un certain nombre de fonctions mathématiques, ainsi que des fonctions permettant des calculs sur les listes (tables) de complexes, des transformations géométriques, \ldots{} etc. Ces fonctions sont dans l'espace de noms \luadrawenv.
\end{itemize}


\paragraph{Système de coordonnées. Repérage}

\begin{itemize}
\item L'objet instancié (\emph{g} ici dans l'exemple) possède :
    \begin{enumerate}
        \item Une vue originelle : c'est le pavé de $\mathbf R^2$ défini par l'option \opt{window} à la création. Celui-ci \textbf{ne doit pas être modifié} par la suite.
        
        \item Une vue courante : c'est un pavé de $\mathbf R^2$ qui doit être inclus dans la vue originelle, ce qui sort de ce pavé sera clippé. Par défaut la vue courante est la vue originelle. Pour retrouver la vue courante on peut utiliser la méthode \cmd{g:Getview()} qui renvoie une table \code{\{x1,x2,y1,y2\}}, celle-ci représente le pavé $[x_1;x_2]\times [y_1;y_2]$.
        
        \item Une matrice de transformation : celle-ci est initialisée à la matrice identité. Lors d'une instruction de dessin les points sont automatiquement transformés par cette matrice avant d'être envoyés à TikZ.
        
        \item Un système de coordonnées (repère cartésien) lié à la vue courante, c'est le repère de l'utilisateur. Par défaut c'est le repère canonique de $\mathbf R^2$, mais il est possible d'en changer. Admettons que la vue courante soit le pavé $[-5;5]\times[-5;5]$, il est possible par exemple, de décider que ce pavé représente l'intervalle $[-1;12]$ pour les abscisses et l'intervalle $[0;8]$ pour les ordonnées, la méthode qui fait ce changement va modifier la matrice de transformation du graphe, de telle sorte que pour l'utilisateur tout se passe comme s'il était dans le pavé $[-1;12]\times [0,8]$. On peut retrouver les intervalles du repère de l'utilisateur avec les méthodes:\par
        \hfil\cmd{g:Xinf()}, \cmd{g:Xsup()}, \cmd{g:Yinf()} et \cmd{g:Ysup()}.\hfil
    \end{enumerate}
\item On utilise les nombres complexes pour représenter les points ou les vecteurs dans le repère cartésien de l'utilisateur.

\item Dans l'export TikZ les coordonnées seront différentes car le coin inférieur gauche (hors marges) aura pour coordonnées $(0,0)$, et le coin supérieur droit (hors marges) aura des coordonnées correspondant à la taille (hors marges) du graphique, et avec $1$ cm par unité sur les deux axes. Ce qui fait que normalement, TikZ ne devrait manipuler que de \og petits\fg\ nombres.

\item La conversion se fait automatiquement avec la méthode \cmd{g:strCoord(x, y)} qui renvoie une chaîne de la forme \emph{(a,b)}, où \emph{a} et \emph{b} sont les coordonnées pour TikZ, ou bien avec la méthode \cmd{g:Coord(z)} qui renvoie aussi une chaîne de la forme \emph{(a,b)} représentant les coordonnées TikZ du point d'affixe $z$ dans le repère de l'utilisateur.
\end{itemize}

\subsection{Peut-on utiliser directement du TikZ dans l'environnement \luadrawenv ?}

Supposons que l'on soit en train de créer un graphique nommé \emph{g} dans un environnement \luadrawenv. Il est possible d'écrire une instruction TikZ lors de cette création, mais pas en utilisant \code{tex.sprint("<instruction TikZ>")}, car alors cette instruction ne ferait pas partie du graphique \emph{g}. Il faut pour cela utiliser la méthode \cmd{g:Writeln("<instruction TikZ>")}, avec la contrainte que \textbf{les antislash doivent être doublés}, et sans oublier que les coordonnées graphiques d'un point dans \emph{g} ne sont pas les mêmes pour TikZ. Par exemple : 

\begin{Luacode}
g:Writeln("\\draw"..g:Coord(Z(1,-1)).." node[red] {Texte};")
\end{Luacode}

Ou encore pour changer des styles :
\begin{Luacode}
g:Writeln("\\tikzset{every node/.style={fill=white}}")
\end{Luacode}

Dans une présentation beamer, cela peut aussi être utilisé pour insérer des pauses dans un graphique :
\begin{Luacode}
g:Writeln("\\pause")
\end{Luacode}  
%
\section{Méthodes graphiques}

On peut créer des lignes polygonales, des courbes, des chemins, des points, des labels.


\subsection{Lignes polygonales}

\subsubsection{Dpolyline}

\textbf{Une ligne polygonale est une liste (table) de composantes connexes, et une composante connexe est une liste (table) de complexes qui représentent les affixes des points}. Par exemple l'instruction :
\begin{Luacode}
local Z = luadraw.cpx.Z -- alias
local L = { {Z(-4,0), Z(0,2), Z(1,3)}, {Z(0,0), Z(4,-2), Z(1,-1)} }
\end{Luacode}
crée une ligne polygonale à deux composantes dans une variable \emph{L}.

\paragraph{Dessin d'une ligne polygonale.}

C'est la méthode \cmdln{g:Dpolyline(L \fac{, close, draw\_options, clip})} où \emph{g} désigne le graphique en cours de création. L'argument \argu{L} est une ligne polygonale, \argu{close} un argument facultatif qui vaut \true ou \false indiquant si la ligne doit être refermée ou non (\false par défaut), \argu{draw\_options} est une chaîne de caractères qui sera passée directement à l'instruction \drawcmd dans l'export (vide par défaut). L'argument \argu{clip} doit contenir ou bien \nil (valeur par défaut) ou bien une table de la forme \code{\{x1,x2,y1,y2\}},  dans le premier cas la ligne est clippée par la fenêtre 2D courante \textbf{après} sa transformation par la matrice 2D du graphe, dans le second cas la ligne est clippée par la fenêtre $[x_1;x_2]\times[y_1;y_2]$ \textbf{avant} d'être transformée par la matrice du graphe.

\subsubsection{Options de dessin d'une ligne polygonale}

On peut passer des options de dessin directement à l'instruction \drawcmd dans l'export, mais elles auront un effet local uniquement. Il est possible de modifier ces options de manière globale avec la méthode \cmdln{g:Lineoptions(style \fac{, color, width, arrows})} lorsqu'un des arguments est égal à \nil, c'est sa valeur par défaut qui est utilisée.

\begin{itemize}
\item \argu{color} est une chaîne de caractères représentant une couleur connue de TikZ (\val{"black"} par défaut),
\item \argu{style} est une chaîne de caractères représentant le type de ligne à dessiner, ce style peut être :
 \begin{itemize}
  \item \val{"noline"} : trait non dessiné,
  \item \val{"solid"} : trait plein, valeur par défaut,
  \item \val{"dotted"} : trait pointillé,
  \item \val{"dashed"} : tirets,
  \item style personnalisé : l'argument \argu{style} peut être une chaîne de la forme (exemple) \code{"\{2.5pt\}\{2pt\}"} ce qui signifie : un trait de 2.5pt suivi d'un espace de 2pt, le nombre de valeurs peut être supérieur à 2, ex : \code{"\{2.5pt\}\{2pt\}\{1pt\}\{2pt\}"} (succession de on, off).
  \end{itemize}
\item \argu{width} est un nombre représentant l'épaisseur de ligne exprimée en dixième de points, par exemple $8$ pour une épaisseur réelle de 0.8pt (valeur de $4$ par défaut),
\item \argu{arrows} est une chaîne qui précise le type de flèche qui sera dessiné, cela peut être :
 \begin{itemize}
  \item \val{"-"} qui signifie pas de flèche (valeur par défaut),
  \item \val{"->"} qui signifie une flèche à la fin,
  \item \val{"<-"} qui signifie une flèche au début,
  \item \val{"<->"} qui signifie une flèche à chaque bout.\\
  \textbf{ATTENTION} : la flèche n'est pas dessinée lorsque l'argument \argu{close} vaut \true.
  \end{itemize}
\end{itemize}


On peut modifier les options individuellement avec les méthodes : 
\begin{itemize}
    \item \cmd{g:Linecolor(color)}, 
    \item \cmd{g:Linestyle(style)}, 
    \item \cmd{g:Linewidth(width)}, 
    \item \cmd{g:Arrows(arrows)}, 
    \item plus les méthodes suivantes : 
        \begin{itemize}
            \item \cmd{g:Lineopacity(opacity)} qui règle l'opacité du tracé de la ligne, l'argument \argu{opacity} doit être une
valeur entre $0$ (totalement transparent) et $1$ (totalement opaque), par défaut la valeur est de $1$.
            \item \cmd{g:Linecap(style)}, pour jouer sur les extrémités de la ligne, l'argument \argu{style} est une chaîne qui peut valoir :
            \begin{itemize}
                \item \val{"butt"} (bout droit au point d'arrêt, valeur par défaut),
                \item \val{"round"} (bout arrondi en demi-cercle),
                \item \val{"square"} (bout \og arrondi\fg\ en carré).
            \end{itemize}
          
          \item \cmd{g:Linejoin(style)}, pour jouer sur la jointure entre segments, l'argument \argu{style} est une chaîne qui peut valoir :
            \begin{itemize}
                \item \val{"miter"} (coin pointu, valeur par défaut),
                \item \val{"round"} (ou coin arrondi),
                \item \val{"bevel"} (coin coupé).
            \end{itemize}
        \end{itemize}
\end{itemize}


\paragraph{Options de remplissage d'une ligne polygonale.}

C'est la méthode \cmdln{g:Filloptions(style \fac{, color, opacity, evenodd})} (qui utilise la librairie \emph{patterns} de TikZ, celle-ci est chargée avec le paquet). Lorsqu'un des arguments vaut \nil, c'est sa valeur par défaut qui est utilisée :

\begin{itemize}
  \item \argu{color} est une chaîne de caractères représentant une couleur connue de TikZ (\val{"black"} par défaut).
  \item \argu{style} est une chaîne de caractères représentant le type de remplissage, ce style peut être :
      \begin{itemize}
      \item \val{"none"} : pas de remplissage, c'est la valeur par défaut,
      \item \val{"full"} : remplissage plein,
      \item \val{"bdiag"} : hachures descendantes de gauche à droite,
      \item \val{"fdiag"} : hachures montantes de gauche à droite,
      \item \val{"horizontal"} : hachures horizontales,
      \item \val{"`vertical"} : hachures verticales,
      \item \val{"hvcross"} : hachures horizontales et verticales,
      \item \val{"diagcross"} : diagonales descendantes et montantes,
      \item \val{"gradient"} : dans ce cas le remplissage se fait avec un gradient défini avec la méthode \cmd{g:Gradstyle(chaîne)}, ce style est passé tel quel à l'instruction \drawcmd. Par défaut la chaîne définissant le style de gradient est : \codeln{"left color = white,right color = red"}
      \item tout autre style connu de la librairie \emph{patterns} est également possible.
      \end{itemize}
  \item \argu{opacity} : nombre entre $0$ et $1$ ($1$ par défaut).
  \item \argu{evenodd} : booléen indiquant si l'option \emph{even odd rule} doit être utilisée pour le remplissage (\false par défaut).
\end{itemize}

On peut modifier certaines options individuellement avec les méthodes :
\begin{itemize}
    \item \cmd{g:Fillopacity(opacity)}, 
    \item \cmd{g:Filleo(evenodd)}.
\end{itemize}

\begin{demo}{Champ de vecteurs, courbe intégrale de $y'= 1-xy^2$}
\begin{luadraw}{name=champ}
local ld = luadraw
local cpx = ld.cpx
local Z = cpx.Z
local g = ld.graph:new{window={-5,5,-5,5},bg="Cyan!30",size={10,10}}
local f = function(x,y) -- éq. diff. y'= 1-x*y^2=f(x,y)
    return 1-x*y^2
end
local A = Z(-1,1) -- A = -1+i
local deltaX, deltaY, long = 0.5, 0.5, 0.4
local champ = function(f)
    local vecteurs, v = {}
    for y = g:Yinf(), g:Ysup(), deltaY do
        for x = g:Xinf(), g:Xsup(), deltaX do
            v = Z(1,f(x,y)) -- coordonnées 1 et f(x,y)
            v = v/cpx.abs(v)*long -- normalisation de v
            table.insert(vecteurs, {Z(x,y), Z(x,y)+v} )
        end
    end
    return vecteurs -- vecteurs est une ligne polygonale
end

g:Daxes( {0,1,1}, {labelpos={"none","none"}, arrows="->"} )
g:Dpolyline( champ(f), "->,blue")
g:Dodesolve(f, A.re, A.im, {t={-2.75,5},draw_options="red,line width=0.8pt"})
g:Dlabeldot("$A$", A, {pos="S"})
g:Show()
\end{luadraw}
\end{demo}
\label{champ}


\subsection{Segments et droites}

\textbf{Un segment est une liste (table) de deux complexes représentant les extrémités. Une droite est une liste (table) de deux complexes, le premier représente un point de la droite, et le second un vecteur directeur de la droite (celui-ci doit être non nul).}

\subsubsection{Dangle}
\begin{itemize}
    \item La méthode \cmd{g:Dangle(B, A, C \fac{, r, draw\_options})} dessine l'angle \(BAC\) avec un parallélogramme (deux côtés seulement sont dessinés), l'argument facultatif \argu{r} précise la longueur d'un côté ($0.25$ par défaut). L'argument \argu{draw\_options} est une chaîne (vide par défaut) qui sera passée telle quelle à l'instruction  \drawcmd.
    \item La fonction \cmd{ld.angleD(B, A, C, r)} renvoie la liste des points de cet angle.
\end{itemize}

\newcommand*{\arguScale}{%
L'argument \argu{scale} (qui vaut $1$ par défaut) est soit un nombre (pourcentage) permettant de faire varier la longueur du segment affiché (à partir de son milieu), soit une table de deux nombres (pourcentages) \{scaleA, scaleB\} permettant de faire varier la longueur du segment affiché à chacune de ses extrémités.
}%

\subsubsection{Dbissec}
\begin{itemize}
    \item La méthode \cmd{g:Dbissec(B, A, C \fac{, interior, draw\_options, scale})} dessine une bissectrice de l'angle géométrique BAC, intérieure si l'argument facultatif \argu{interior} vaut \true (valeur par défaut), extérieure sinon. L'argument \argu{draw\_options} est une chaîne (vide par défaut) qui sera passée telle quelle à l'instruction \drawcmd. \arguScale
  \item La fonction \cmd{ld.bissec(B, A, C, interior)} renvoie cette bissectrice sous forme d'une liste $\{A,u\}$ où $u$ est un vecteur directeur de la droite.
\end{itemize}

\subsubsection{Dhline}
La méthode \cmd{g:Dhline(d \fac{, draw\_options, scale})} dessine une demi-droite, l'argument \argu{d} doit être une liste de deux nombres complexes $\{A,B\}$, c'est la demi-droite $[A,B)$ qui est dessinée.

Variante : \cmd{g:Dhline(A, B \fac{, draw\_options, scale})}. L'argument \argu{draw\_options} est une chaîne (vide par défaut) qui sera passée telle quelle à l'instruction \drawcmd. \arguScale


\subsubsection{Dline}
La méthode \cmd{g:Dline(d \fac{, draw\_options, scale})} trace la droite \argu{d}, celle-ci est une liste du type $\{A,u\}$ où $A$ représente un point de la droite (un nombre complexe) et $u$ un vecteur directeur (un nombre complexe non nul). 

Variante : la méthode \cmd{g:Dline(A, B \fac{, draw\_options, scale})} trace la droite passant par les points \argu{A} et \argu{B} (deux nombres complexes). L'argument \argu{draw\_options} est une chaîne (vide par défaut) qui sera passée telle quelle à l'instruction \drawcmd. \arguScale

\subsubsection{DlineEq}
\begin{itemize}
    \item La méthode \cmd{g:DlineEq(a, b, c \fac{, draw\_options, scale})} dessine la droite d'équation \(ax+by+c=0\). L'argument \argu{draw\_options} est une chaîne (vide par défaut) qui sera passée telle quelle à l'instruction \drawcmd. \arguScale
    \item La fonction \cmd{ld.lineEq(a, b, c)} renvoie la droite d'équation \(ax+by+c=0\) sous la forme d'une liste $\{A,u\}$ où $A$ est un point de la droite et $u$ un vecteur directeur.
\end{itemize}

\subsubsection{Dmarkarc}
La méthode \cmd{g:Dmarkarc(b, a, c, r, n \fac{,long,espace,draw\_options})} permet de marquer l'arc de cercle de centre \argu{a}, de rayon \argu{r}, allant de \argu{b} à \argu{c}, avec \argu{n} petits segments. Par défaut la longueur (argument \argu{long}) est de $0.25$, et l'espacement (argument \argu{espace}) et de $0.0625$. L'argument \argu{draw\_options} est une chaîne (vide par défaut) qui sera passée telle quelle à l'instruction \drawcmd.

\subsubsection{Dmarkseg}
La méthode \cmd{g:Dmarkseg(a, b, n, \fac{long, espace, angle, draw\_options})} permet de marquer le segment défini par \argu{a} et \argu{b} avec \argu{n} petits segments penchés de \argu{angle} degrés (45° par défaut). Par défaut la longueur (argument \argu{long}) est de $0.25$, et l'espacement (argument \argu{espace}) et de $0.125$. L'argument \argu{draw\_options} est une chaîne (vide par défaut) qui sera passée telle quelle à l'instruction \drawcmd.

\subsubsection{Dmed}
\begin{itemize}
    \item La méthode \cmd{g:Dmed(A, B \fac{, draw\_options, scale})} trace la médiatrice du segment $[A;B]$.

  Variante : \cmd{g:Dmed(seg \fac{, draw\_options, scale})} où segment est une liste de deux points représentant le segment. L'argument \argu{draw\_options} est une chaîne (vide par défaut) qui sera passée telle quelle à l'instruction \drawcmd. \arguScale
  \item La fonction \cmd{ld.med(A, B)} (ou \cmd{ld.med(seg)}) renvoie la médiatrice du segment $[A;B]$ sous la forme d'une liste $\{C,u\}$ où $C$ est un point de la droite et $u$ un vecteur directeur.
\end{itemize}

\subsubsection{Dparallel}
\begin{itemize}
    \item La méthode \cmd{g:Dparallel(d, A \fac{, draw\_options, scale})} trace la parallèle à \argu{d} passant par \argu{A}. L'argument \argu{d} peut être soit une droite (une liste constituée d'un point et un vecteur directeur) soit un vecteur non nul. L'argument \argu{draw\_options} est une chaîne (vide par défaut) qui sera passée telle quelle à l'instruction \drawcmd. \arguScale
    \item La fonction \cmd{ld.parallel(d, A)} renvoie la parallèle à \argu{d} passant par \argu{A} sous la forme d'une liste $\{B,u\}$ où $B$ est un point de la droite et $u$ un vecteur directeur.
\end{itemize}

\subsubsection{Dperp}
\begin{itemize}
    \item La méthode \cmd{g:Dperp(d, A \fac{, draw\_options, scale})} trace la perpendiculaire à \argu{d} passant par \argu{A}. L'argument \argu{d} peut être soit une droite (une liste constituée d'un point et un vecteur directeur) soit un vecteur non nul. L'argument \argu{draw\_options} est une chaîne (vide par défaut) qui sera passée  telle quelle à l'instruction \drawcmd. \arguScale
    \item La fonction \cmd{ld.perp(d, A)} renvoie la perpendiculaire à \argu{d} passant par \argu{A} sous la forme d'une liste $\{B,u\}$ où $B$ est un point de la droite et $u$ un vecteur directeur.
\end{itemize}

\subsubsection{Dseg}
La méthode \cmd{g:Dseg(seg \fac{, scale, draw\_options})} dessine le segment défini par l'argument \argu{seg} qui doit être une liste de deux complexes. L'argument \argu{draw\_options} est une chaîne (vide par défaut) qui sera passée telle quelle à l'instruction \drawcmd. \arguScale

\subsubsection{Dtangent}
\begin{itemize}
    \item La méthode \cmd{g:Dtangent(p, t0 \fac{, long, draw\_options})} dessine la tangente à la courbe paramétrée par \argu{p}: \(t \mapsto p(t)\) (à valeurs complexes), au point de paramètre \argu{t0}. Si l'argument \argu{long} vaut \nil (valeur par défaut) alors c'est la droite entière qui est dessinée, sinon c'est un segment de longueur \argu{long}. L'argument \argu{draw\_options} est une chaîne (vide par défaut) qui sera passée telle quelle à l'instruction \drawcmd.
    \item La fonction \cmd{ld.tangent(p, t0 \fac{, long})} renvoie soit la droite, soit un segment (suivant que \argu{long} vaut \nil ou pas).
\end{itemize}

\subsubsection{DtangentC}
\begin{itemize}
    \item La méthode \cmd{g:DtangentC(f, x0 \fac{, long, draw\_options})} dessine la tangente à la courbe cartésienne d'équation \(y=f(x)\), au point d'abscisse \argu{x0}. Si l'argument \argu{long} vaut \nil (valeur par défaut) alors c'est la droite entière qui est dessinée, sinon c'est un segment de longueur \argu{long}. L'argument \argu{draw\_options} est une chaîne (vide par défaut) qui sera passée telle quelle à l'instruction \drawcmd.
    \item La fonction \cmd{ld.tangentC(f, x0 \fac{, long})} renvoie soit la droite, soit un segment (suivant que \argu{long} vaut \nil ou pas).
\end{itemize}

\subsubsection{DtangentI}
\begin{itemize}
    \item La méthode \cmd{g:DtangentI(f, x0, y0 \fac{, long, draw\_options})} dessine la tangente à la courbe implicite d'équation \(f(x,y)=0\), au point \argu{(x0,y0)} supposé sur la courbe. Si l'argument \argu{long} vaut \nil (valeur par défaut) alors c'est la droite entière qui est dessinée, sinon c'est un segment de longueur \argu{long}. L'argument \argu{draw\_options} est une chaîne (vide par défaut) qui sera passée telle quelle à l'instruction \drawcmd.
    \item La fonction \cmd{ld.tangentI(f, x0, y0 \fac{, long})} renvoie soit la droite, soit un segment (suivant que \argu{long} vaut \nil ou pas).
\end{itemize}

\subsubsection{Dtangent\_from}
\begin{itemize}
    \item La méthode \cmd{g:Dtangent\_from(A, p, t1, t2 \fac{, dp, draw\_options, out, scale})} dessine la ou les tangentes à la courbe paramétrée par \argu{p}: \(t \mapsto p(t)\) (à valeurs complexes) issue(s) du point \argu{A} (nombre complexe). Les arguments \argu{t1} et \argu{t2} représentent les bornes de l'intervalle de recherche. L'argument optionnel \argu{dp} est une fonction représentant la dérivée de la fonction \argu{p}, par défaut cet argument vaut \nil et la dérivée de \argu{p} est remplacée par une approximation. L'argument optionnel \argu{draw\_options} est une chaîne (vide par défaut) qui sera passée telle quelle à l'instruction \drawcmd. L'argument optionnel \argu{out}, s'il est utilisé, doit être le nom d'une variable, celle-ci doit être une table, elle contiendra à la fin de l'exécution les points de la courbe pour lesquels la tangente passe par \argu{A}. \arguScale
    \item La fonction \cmd{ld.tangent\_from(A, p, t1, t2 \fac{, dp})} renvoie la liste les points de la courbe paramétrée par \argu{p}, sur l'intervalle \([t_1;t_2]\) pour lesquels la tangente passe par \argu{A} (nombre complexe). S'il n'y en a pas, la fonction renvoie une liste vide.
\end{itemize}

\begin{demo}{Tangentes issues d'un point}
\begin{luadraw}{name=tangent_from}
local ld = luadraw
local cpx = ld.cpx
local i, cos, sin, pi = cpx.I, math.cos, math.sin, math.pi
local g = ld.graph:new{window={-5,5,-5,5}, size={10,10}}
local p1 = function(t) return t+i*(t^2/4-1) end -- parabole
local p2 = function(t) return 1/2+i+ld.rotate(2*cos(t)+i*sin(t),15) end -- ellipse
local p3 = function(t) return math.sinh(t)+i*math.cosh(t)-i*2 end -- branche d'hyperbole
local p4 = function(t) return 2*cos(3*t)*cpx.exp(i*t) end -- autre
local P = 0.5-1.25*i
local draw = function(p,t1,t2)
    local axis_style = { arrows='-Stealth', legend={'$x$','$y$'} }
    local S = {}
    g:Daxes({0, 1, 1}, axis_style)
    g:Dparametric(p,{t={t1,t2}, draw_options="line width=0.8pt, Crimson"})
    g:Dtangent_from(P,p,t1,t2,"blue",S)
    g:Ddots(S,"Crimson"); g:Ddots(P); g:Dlabel("$A$",P,{pos="S"})
end
g:Saveattr(); g:Viewport(-5,-0.25,0.25,5); g:Coordsystem(-3,4,-3,4); draw(p1,-4,4); g:Restoreattr()
g:Saveattr(); g:Viewport(0.25,5,0.25,5); g:Coordsystem(-3,3,-3,3); draw(p2,-pi,pi); g:Restoreattr()
g:Saveattr(); g:Viewport(-5,-0.25,-5,-0.25); g:Coordsystem(-3,4,-3,3); draw(p3,-4,4); g:Restoreattr()
g:Saveattr(); g:Viewport(0.25,5,-5,-0.25); g:Coordsystem(-3,3,-3,3); draw(p4,-pi,pi); g:Restoreattr()
g:Show()
\end{luadraw}
\end{demo}

\subsubsection{Dnormal}
\begin{itemize}
    \item La méthode \cmd{g:Dnormal(p, t0 \fac{, long, draw\_options})} dessine la normale à la courbe paramétrée par \argu{p}: \(t \mapsto p(t)\) (à valeurs complexes), au point de paramètre \argu{t0}. Si l'argument \argu{long} vaut \nil (valeur par défaut) alors c'est la droite entière qui est dessinée, sinon c'est un segment de longueur \argu{long}. L'argument \argu{draw\_options} est une chaîne (vide par défaut) qui sera passée telle quelle à l'instruction \drawcmd.
    \item La fonction \cmd{ld.normal(p, t0 \fac{, long})} renvoie soit la droite, soit un segment (suivant que \emph{long} vaut \nil ou pas).
\end{itemize}

\subsubsection{DnormalC}
\begin{itemize}
    \item La méthode \cmd{g:DnormalC(f, x0 \fac{, long, draw\_options})} dessine la normale à la courbe cartésienne d'équation \(y=f(x)\), au point d'abscisse \argu{x0}. Si l'argument \argu{long} vaut \nil (valeur par défaut) alors c'est la droite entière qui est dessinée, sinon c'est un segment de longueur \argu{long}. L'argument \argu{draw\_options} est une chaîne (vide par défaut) qui sera passée telle quelle à l'instruction \drawcmd.
    \item La fonction \cmd{ld.normalC(f, x0 \fac{, long})} renvoie soit la droite, soit un segment (suivant que \argu{long} vaut \nil ou pas).
\end{itemize}

\subsubsection{DnormalI}
\begin{itemize}
    \item La méthode \cmd{g:DnormalI(f, x0, y0 \fac{, long, draw\_options})} dessine la normale à la courbe implicite d'équation \(f(x,y)=0\), au point \argu{(x0,y0)} supposé sur la courbe. Si l'argument \argu{long} vaut \nil (valeur par défaut) alors c'est la droite entière qui est dessinée, sinon c'est un segment de longueur \argu{long}. L'argument \argu{draw\_options} est une chaîne (vide par défaut) qui sera passée telle quelle à l'instruction \drawcmd.
    \item La fonction \cmd{ld.normalI(f, x0, y0 \fac{, long})} renvoie soit la droite, soit un segment (suivant que \argu{long} vaut \nil ou pas).
\end{itemize}
  

\begin{demo}{Symétrique de l'orthocentre}
\begin{luadraw}{name=orthocentre}
local ld = luadraw
local cpx = ld.cpx
local g = ld.graph:new{window={-5,5,-5,5},bg="cyan!30",size={10,10}}
local i = cpx.I
local A, B, C = 4*i, -2-2*i, 3.5
local h1, h2 = ld.perp({B,C-B},A), ld.perp({A,B-A},C) -- hauteurs
local A1, F = ld.proj(A,{B,C-B}), ld.proj(C,{A,B-A}) -- projetés
local H = ld.interD(h1,h2) -- orthocentre
local A2 = 2*A1-H -- symétrique de H par rapport à BC
g:Dpolyline({A,B,C},true, "draw=none,fill=Maroon,fill opacity=0.3") -- fond du triangle
g:Linewidth(6); g:Filloptions("full", "blue", 0.2)
g:Dangle(C,A1,A,0.25); g:Dangle(B,F,C,0.25) -- angles droits
g:Linecolor("black"); g:Filloptions("full","cyan",0.5)
g:Darc(H,C,A2,1); g:Darc(B,A,A1,1) -- arcs
g:Filloptions("none","black",1) -- on rétablit l'opacité à 1
g:Dmarkarc(H,C,A1,1,2); g:Dmarkarc(A1,C,A2,1,2) -- marques
g:Dmarkarc(B,A,H,1,2)
g:Linewidth(8); g:Linecolor("black")
g:Dseg({A,B},1.25); g:Dseg({C,B},1.25); g:Dseg({A,C},1.25) -- côtés
g:Linecolor("red"); g:Dcircle(A,B,C) -- cercle
g:Linecolor("blue"); g:Dline(h1); g:Dline(h2) -- hauteurs
g:Dseg({A2,C}); g:Linecolor("red"); g:Dseg({H,A2}) -- segments
g:Dmarkseg(H,A1,2); g:Dmarkseg(A1,A2,2) -- marques
g:Labelcolor("blue") -- pour les labels
g:Dlabel("$A$",A, {pos="NW",dist=0.1}, "$B$",B, {pos="SW"}, "$A_2$",A2,{pos="E"},
"$C$", C, {pos="S" }, "$H$", H, {pos="NE"}, "$A_1$", A1, {pos="SW"})
g:Linecolor("black"); g:Filloptions("full"); g:Ddots({A,B,C,H,A1,A2}) -- dessin des points
g:Show()
\end{luadraw}
\end{demo}

\subsection{Figures géométriques}

\subsubsection{Darc}
\begin{itemize}
    \item La méthode \cmd{g:Darc(B, A, C, r \fac{, sens, draw\_options})} dessine un arc de cercle de centre \argu{A} (complexe), de rayon \argu{r}, allant de \argu{B} (complexe) vers \argu{C} (complexe) dans le sens trigonométrique si l'argument \argu{sens} vaut $1$ (valeur par défaut), le sens inverse sinon. L'argument \argu{draw\_options} est une chaîne (vide par défaut) qui sera passée telle quelle à l'instruction \drawcmd.
    \item La fonction \cmd{ld.arc(B, A, C, r, sens)} renvoie la liste des points de cet arc (ligne polygonale). 
    \item La fonction \cmd{ld.arcb(B, A, C, r, sens)} renvoie cet arc sous forme d'un chemin (voir \emph{Dpath} page \pageref{Dpath}) utilisant des courbes de Bézier.
\end{itemize}

\subsubsection{Dcircle}
\begin{itemize}
    \item La méthode \cmd{g:Dcircle(c, r \fac{, d, draw\_options})} trace un cercle. Lorsque l'argument \argu{d} vaut \nil, c'est le cercle de centre \argu{c} (complexe) et de rayon \argu{r}, lorsque \argu{d} est précisé (complexe) alors c'est le cercle passant par les points d'affixe \argu{c}, \argu{r} et \argu{d}. L'argument \argu{draw\_options} est une chaîne (vide par défaut) qui sera passée telle quelle à l'instruction \drawcmd. Autre syntaxe possible :  \cmd{g:Dcircle(C \fac{, draw\_options})} où \argu{C}$=\{c,r \fac{,d}\}$.
  \item La fonction \cmd{ld.circle(c, r \fac{, d})} renvoie la liste des points de ce cercle (ligne polygonale). 
  \item La fonction \cmd{ld.circle(\{c, r \fac{, d}\}, nbdots)} renvoie la liste des points de ce cercle (ligne polygonale) avec \argu{nbdots} points. 
  \item La fonction \cmd{ld.circleb(c, r \fac{, d})} renvoie ce cercle sous forme d'un chemin (voir \emph{Dpath} page \pageref{Dpath}) utilisant des courbes de Bézier.
\end{itemize}

\subsubsection{Dellipse}
\begin{itemize}
    \item  La méthode \cmd{g:Dellipse(c, rx, ry \fac{, inclin, draw\_options})} dessine l'ellipse de centre \argu{c} (complexe), les arguments \argu{rx} et \argu{ry} précisent les deux rayons (sur $x$ et sur $y$), l'argument facultatif \argu{inclin} est un angle en degrés qui indique l'inclinaison de l'ellipse par rapport à l'axe \(Ox\) (angle nul par défaut). L'argument \argu{draw\_options} est une chaîne (vide par défaut) qui sera passée telle quelle à l'instruction \drawcmd.
  \item La fonction \cmd{ld.ellipse(c, rx, ry \fac{, inclin})} renvoie la liste des points de cette ellipse (ligne polygonale). 
  \item La fonction \cmd{ld.ellipse(\{c, rx, ry \fac{, inclin}\}, nbdots)} renvoie la liste des points de cette ellipse (ligne polygonale) avec \argu{nbdots} points. 
  \item La fonction \cmd{ld.ellipseb(c, rx, ry \fac{, inclin})} renvoie cette ellipse sous forme d'un chemin (voir \emph{Dpath} page \pageref{Dpath}) utilisant des courbes de Bézier.
\end{itemize}

\subsubsection{Dellipticarc}
\begin{itemize}
    \item La méthode \cmd{g:Dellipticarc(B, A, C, rx, ry \fac{, sens, inclin, draw\_options})} dessine un arc d'ellipse de centre \argu{A} (complexe) de rayons \argu{rx} et \argu{ry}, faisant un angle égal à \argu{inclin} par rapport à l'axe \(Ox\) (angle nul par défaut), allant de \argu{B} (complexe) vers \argu{C} (complexe) dans le sens trigonométrique si l'argument \argu{sens} vaut $1$ (valeur par défaut), le sens inverse sinon. L'argument \argu{draw\_options} est une chaîne (vide par défaut) qui sera passée telle quelle à l'instruction \drawcmd.
    \item La fonction \cmd{ld.ellipticarc(B, A, C, rx, ry \fac{, sens, inclin})} renvoie la liste des points de cet arc (ligne polygonale). 
    \item La fonction \cmd{ld.ellipticarcb(B, A, C, rx, ry \fac{, sens, inclin})} renvoie cet arc sous forme d'un chemin (voir \emph{Dpath} page \pageref{Dpath}) utilisant des courbes de Bézier.
\end{itemize}

\subsubsection{Dparallelogram}
\begin{itemize}
    \item La méthode \cmd{g:Dparallelogram(a, u, v \fac{, draw\_options})} où \argu{a}, \argu{u}, \argu{v} désigne trois nombres complexes désignant un sommet et deux vecteurs, dessine le parallélogramme de sommets $\{a,a+u,a+u+v,a+v\}$. L'argument \argu{draw\_options} est une chaîne (vide par défaut) qui sera passée telle quelle à l'instruction \drawcmd.
    \item La fonction \cmd{ld.parallelogram(a, u, v)} renvoie la liste des sommets de ce parallélogramme.
\end{itemize}

\subsubsection{Dpolyreg}
\begin{itemize}
    \item La méthode \cmd{g:Dpolyreg(sommet1, sommet2, nbcotes, sens \fac{, draw\_options)}} ou \par \cmd{g:Dpolyreg(centre, sommet, nbcotes \fac{, draw\_options})} dessine un polygone régulier. L'argument \argu{draw\_options} est une chaîne (vide par défaut) qui sera passée telle quelle à l'instruction \drawcmd.
    \item La fonction \cmd{ld.polyreg(sommet1, sommet2, nbcotes, sens)} et la fonction \cmd{ld.polyreg(centre, sommet, nbcotes)}, renvoient la liste des sommets de ce polygone régulier.
\end{itemize}

\subsubsection{Drectangle}
\begin{itemize}
    \item La méthode \cmd{g:Drectangle(a, b, c \fac{, draw\_options})} dessine le rectangle ayant comme sommets consécutifs \argu{a} et \argu{b} et dont le côté opposé passe par \argu{c}. L'argument \argu{draw\_options} est une chaîne (vide par défaut) qui sera passée telle quelle à l'instruction \drawcmd. \par    
    La fonction \cmd{ld.rectangle(a, b, c)} renvoie la liste des sommets de ce rectangle.
    
    \item La méthode \cmd{g:Drectangle(a, b \fac{, draw\_options})} dessine le rectangle dont les côtés sont parallèles aux axes et ayant comme sommets opposés \argu{a} et \argu{b}. L'argument \argu{draw\_options} est une chaîne (vide par défaut) qui sera passée telle quelle à l'instruction \drawcmd. \par
    La fonction \cmd{ld.rectangle(a, b)} renvoie la liste des sommets de ce rectangle.
\end{itemize}

\subsubsection{Dsequence}
\begin{itemize}
    \item La méthode \cmd{g:Dsequence(f, u0, n \fac{, draw\_options})} fait le dessin des "escaliers" de la suite récurrente définie par son premier terme \argu{u0} et la relation \(u_{k+1}=f(u_k)\). L'argument \argu{f} doit être une fonction d'une variable réelle et à valeurs réelles, l'argument \argu{n} est le nombre de termes calculés. L'argument \argu{draw\_options} est une chaîne (vide par défaut) qui sera passée telle quelle à l'instruction \drawcmd.
    \item La fonction \cmd{ld.sequence(f, u0, n)} renvoie la liste des points constituant ces "escaliers".
\end{itemize}

\begin{demo}{Suite $u_{n+1}=\cos(u_n)$}
\begin{luadraw}{name=sequence}
local ld = luadraw
local cpx = ld.cpx
local g = ld.graph:new{window={-0.1,1.7,-0.1,1.1},size={10,10,0}}
local i, pi, cos = cpx.I, math.pi, math.cos
local f = function(x) return cos(x)-x end
local ell = ld.solve(f,0,pi/2)[1]
local L = ld.sequence(cos,0.2,5) -- u_{n+1}=cos(u_n), u_0=0.2
local seg, z = {}, L[1]
for k = 2, #L do
    table.insert(seg,{z,L[k]})
    z = L[k]
end -- seg est la liste des segments de l'escalier
local styleA = "\\tikzset{->-/.style={decoration={markings, mark="
local styleB = "at position #1 with {\\arrow{Stealth}}}, postaction={decorate}}}"
g:Writeln(styleA..styleB)
g:Daxes({0,1,1}, {arrows="-Stealth"})
g:DlineEq(1,-1,0,"line width=0.8pt,ForestGreen")
g:Dcartesian(cos, {x={0,pi/2},draw_options="line width=1.2pt,Crimson"})
g:Dpolyline(seg,false,"->-=0.65,blue")
g:Dlabel("$u_0$",0.2,{pos="S",node_options="blue"})
g:Dseg({ell, ell*(1+i)},1,"dashed,gray")
g:Dlabel("$\\ell\\approx"..ld.round(ell,3).."$", ell,{pos="S"})
g:Ddots(ell*(1+i)); g:Labelcolor("Crimson")
g:Dlabel("${\\mathcal C}_{\\cos}$",1+i*cos(1),{pos="E"})
g:Labelcolor("ForestGreen"); g:Labelangle(g:Arg(1+i)*180/pi)
g:Dlabel("$y=x$",0.4+i*0.4,{pos="S",dist=0.1})
g:Show()
\end{luadraw}
\end{demo}

La méthode \cmd{g:Arg(z)} calcule et renvoie l'argument \textit{réel} du complexe $z$, c'est à dire son argument (en radians) à l'export dans le repère de TikZ (il faut pour cela appliquer la matrice de transformation du graphe à $z$, puis faire le changement de repère vers celui de TikZ). Si le repère du graphe est orthonormé et si la matrice de transformation est l'identité alors le résultat est identique à celui de \code{cpx.arg(z)} (ce n'est pas le cas dans l'exemple ci-dessus).

De même, la méthode \cmd{g:Abs(z)} calcule et renvoie le module \textit{réel} du complexe $z$, c'est à dire son module à l'export dans le repère de TikZ, c'est donc une longueur en centimètres. Si le repère du graphe est orthonormé avec 1cm par unité sur chaque axe, et si la matrice de transformation est une isométrie alors le résultat est identique à celui de \code{cpx.abs(z)}.


\subsubsection{Dsquare}
\begin{itemize}
    \item La méthode \cmd{g:Dsquare(a, b \fac{, sens, draw\_options})} dessine le carré ayant comme sommets consécutifs \argu{a} et \argu{b}, dans le sens trigonométrique lorsque \argu{sens} vaut $1$ (valeur par défaut). L'argument \argu{draw\_options} est une chaîne (vide par défaut) qui sera passée telle quelle à l'instruction \drawcmd.
  \item La fonction \cmd{ld.square(a, b, sens)} renvoie la liste des sommets de ce carré.
\end{itemize}

\subsubsection{Dwedge}
La méthode \cmd{g:Dwedge(B, A, C, r, sens \fac{, draw\_options})} dessine un secteur angulaire de centre \argu{A} (complexe), de rayon \argu{r}, allant de \argu{B} (complexe) vers \argu{C} (complexe) dans le sens trigonométrique si l'argument \argu{sens} vaut $1$, le sens inverse sinon. L'argument \argu{draw\_options} est une chaîne (vide par défaut) qui sera passée telle quelle à l'instruction \drawcmd.


\subsection{Courbes}

\subsubsection{Paramétriques : Dparametric}

\begin{itemize}
\item La fonction \cmd{ld.parametric(p, t1, t2 \fac{, nbdots, discont, nbdiv})} fait le calcul des points de la courbe et renvoie une ligne polygonale (une liste de listes de complexes, pas de dessin).
  \begin{itemize}
    \item L'argument \argu{p} est le paramétrage, ce doit être une fonction d'une variable réelle $t$ et à valeurs complexes, par exemple :
    \mintinline{Lua}{local p = function(t) return cpx.exp(t*cpx.I) end}
    \item  Les arguments \argu{t1} et \argu{t2} sont obligatoires avec \(t_1 < t_2\), ils forment les bornes de l'intervalle pour le paramètre.
    \item L'argument \argu{nbdots} est facultatif, c'est le nombre de points (minimal) à calculer, il vaut $40$ par défaut.
    \item L'argument \argu{discont} est un booléen facultatif qui indique s'il y a des discontinuités ou non, c'est \false par défaut.
    \item L'argument \argu{nbdiv} est un entier positif qui vaut $5$ par défaut et indique le nombre de fois que l'intervalle entre deux valeurs consécutives du paramètre peut être coupé en deux (dichotomie) lorsque les points correspondants sont trop éloignés.
  \end{itemize}
  
\item La méthode \cmd{g:Dparametric(p, options)} fait le calcul des points et le dessin de la courbe paramétrée par \argu{p}. L'argument \argu{options} est une table dont les champs sont les options possibles. Celles-ci sont, avec leur valeur par défaut:

  \begin{itemize}
      \item \opt{t=\{g:Xinf(),g:Xsup()\}}, c'est l'intervalle pour le paramètre $t$ (par défaut, c'est tout l'intervalle des abscisses de la fenêtre), 
      \item \opt{nbdots=40}, 
      \item \opt{discont=false},
      \item \opt{nbdiv=5},
      \item \opt{draw\_options=""},  chaîne qui sera transmise telle quelle à l'instruction \drawcmd,
      \item \opt{clip=nil}, cette option est soit \nil (valeur par défaut), soit une table \code{\{x1,x2,y1,y2\}}, dans le premier cas la ligne est clippée par la fenêtre 2D courante \textbf{après} sa transformation par la matrice 2D du graphe, dans le second cas la ligne est clippée par la fenêtre $[x_1;x_2]\times[y_1;y_2]$ \textbf{avant} d'être transformée par la matrice du graphe.
  \end{itemize}
\end{itemize}  

\subsubsection{Polaires : Dpolar}

\begin{itemize}
\item La fonction \cmd{ld.polar(rho, t1, t2 \fac{, nbdots, discont, nbdiv})} fait le calcul des points et renvoie une ligne polygonale (pas de dessin). L'argument \argu{rho} est le paramétrage polaire de la courbe, ce doit être une fonction d'une variable réelle $t$ et à valeurs réelles, par exemple :

    \mintinline{Lua}{local rho = function(t) return 4*math.cos(2*t) end}

    Les autres arguments sont identiques aux courbes paramétrées.
\item La méthode \cmd{g:Dpolar(rho, options)} fait le calcul des points et le  dessin de la courbe polaire paramétrée par \emph{rho}. L'argument \argu{options} est une table dont les champs sont les options possibles. Celles-ci sont, avec leur valeur par défaut:

  \begin{itemize}
      \item \opt{t=\{-pi,pi\}}, c'est l'intervalle pour le paramètre $t$,
      \item \opt{nbdots=40}, 
      \item \opt{discont=false},
      \item \opt{nbdiv=5},
      \item \opt{draw\_options=""},  chaîne qui sera transmise telle quelle à l'instruction \drawcmd,
      \item \opt{clip=nil}, cette option est soit \nil (valeur par défaut), soit une table \code{\{x1,x2,y1,y2\}}, dans le premier cas la ligne est clippée par la fenêtre 2D courante \textbf{après} sa transformation par la matrice 2D du graphe, dans le second cas la ligne est clippée par la fenêtre $[x_1;x_2]\times[y_1;y_2]$ \textbf{avant} d'être transformée par la matrice du graphe.
  \end{itemize}
\end{itemize}

\subsubsection{Cartésiennes : Dcartesian}

\begin{itemize}
\item La fonction \cmd{ld.cartesian(f, x1, x2 \fac{, nbdots, discont, nbdiv})} fait le calcul des points de la courbe et renvoie une ligne polygonale (pas de dessin). L'argument \argu{f} est la fonction dont on veut la courbe, ce doit être une fonction d'une variable réelle $x$ et à valeurs réelles, par exemple :

    \mintinline{Lua}{local f = function(x) return 1+3*math.sin(x)*x end}

    Les arguments \argu{x1} et \argu{x2} sont obligatoires et forment les bornes de l'intervalle pour la variable. Les autres arguments sont identiques aux courbes paramétrées.

\item La méthode \cmd{g:Dcartesian(f, options)} fait le calcul des points et le dessin de la courbe de \argu{f}. L'argument \argu{options} est une table dont les champs sont les options possibles. Celles-ci sont, avec leur valeur par défaut:

  \begin{itemize}
      \item \opt{x=\{g:Xinf(),g:Xsup()\}}, c'est l'intervalle pour le paramètre $x$ (par défaut, c'est tout l'intervalle des abscisses de la fenêtre), 
      \item \opt{nbdots=40}, 
      \item \opt{discont=false},
      \item \opt{nbdiv=5},
      \item \opt{draw\_options=""},  chaîne qui sera transmise telle quelle à l'instruction \drawcmd,
      \item \opt{clip=nil}, cette option est soit \nil (valeur par défaut), soit une table \code{\{x1,x2,y1,y2\}}, dans le premier cas la ligne est clippée par la fenêtre 2D courante \textbf{après} sa transformation par la matrice 2D du graphe, dans le second cas la ligne est clippée par la fenêtre $[x_1;x_2]\times[y_1;y_2]$ \textbf{avant} d'être transformée par la matrice du graphe.
  \end{itemize}
\end{itemize}

\subsubsection{Fonctions périodiques : Dperiodic}

\begin{itemize}
\item La fonction \cmd{ld.periodic(f, period, x1, x2 \fac{, nbdots, discont, nbdiv})} fait le calcul des points de la courbe et renvoie une ligne polygonale (pas de dessin).

  \begin{itemize}
    \item L'argument \argu{f} est la fonction dont on veut la courbe, ce doit être une fonction d'une variable réelle $x$ et à valeurs réelles.
    \item L'argument \argu{period} est une table du type $\{a,b\}$ avec \(a<b\) représentant une période de la fonction \argu{f}.
    \item Les arguments \argu{x1} et \argu{x2} sont obligatoires et forment les bornes de l'intervalle pour la variable.
    \item Les autres arguments sont identiques aux courbes paramétrées.
  \end{itemize}
  
\item La méthode \cmd{g:Dperiodic(f, period, options)} fait le calcul des points de la courbe et le dessin de la courbe de \emph{f}. L'argument \argu{options} est une table dont les champs sont les options possibles. Celles-ci sont, avec leur valeur par défaut:

  \begin{itemize}
      \item \opt{x=\{g:Xinf(),g:Xsup()\}}, c'est l'intervalle pour le paramètre $x$ (par défaut, c'est tout l'intervalle des abscisses de la fenêtre), 
      \item \opt{nbdots=40}, 
      \item \opt{discont=false},
      \item \opt{nbdiv=5},
      \item \opt{draw\_options=""},  chaîne qui sera transmise telle quelle à l'instruction \drawcmd,
      \item \opt{clip=nil}, cette option est soit \nil (valeur par défaut), soit une table \code{\{x1,x2,y1,y2\}}, dans le premier cas la ligne est clippée par la fenêtre 2D courante \textbf{après} sa transformation par la matrice 2D du graphe, dans le second cas la ligne est clippée par la fenêtre $[x_1;x_2]\times[y_1;y_2]$ \textbf{avant} d'être transformée par la matrice du graphe.
  \end{itemize}

\end{itemize}

\subsubsection{Fonctions en escaliers : Dstepfunction}

\begin{itemize}
    \item La fonction \cmd{ld.stepfunction(def, discont)} fait le calcul des points de la courbe et renvoie une ligne polygonale (pas de dessin).

      \begin{itemize}
      \item  L'argument \argu{def} permet de définir la fonction en escaliers, c'est une table à deux éléments :
    
    \begin{TeXcode}
      { {x1,x2,x3,...,xn}, {c1,c2,...} }
    \end{TeXcode}
    
      Le premier élément \code{\{x1,x2,x3,\ldots,xn\}} doit être une subdivision du segment \([x_1;x_n]\).
      
      Le deuxième élément \code{\{c1,c2,\ldots\}} est la liste des constantes avec \code{c1} pour le morceau $[x_1;x_2]$, \code{c2} pour le morceau $[x_2;x_3]$, etc.
      
      \item   L'argument \argu{discont} est un booléen qui vaut \true par défaut.
      \end{itemize}
  
    \item La méthode \cmd{g:Dstepfunction(def, options)} fait le calcul des points et le dessin de la courbe de la fonction en escalier.

      \begin{itemize}
      \item L'argument \argu{def} est le même que celui décrit au-dessus (définition de la fonction en escalier).
      \item  L'argument \argu{options} est une table dont les champs sont les options possibles. Celles-ci sont, avec leur valeur par défaut:
    
      \begin{itemize}
          \item \opt{discont=true},
          \item \opt{draw\_options=""},  chaîne qui sera transmise telle quelle à l'instruction \drawcmd,
          \item \opt{clip=nil}, cette option est soit \nil (valeur par défaut), soit une table \code{\{x1,x2,y1,y2\}}, dans le premier cas la ligne est clippée par la fenêtre 2D courante \textbf{après} sa transformation par la matrice 2D du graphe, dans le second cas la ligne est clippée par la fenêtre $[x_1;x_2]\times[y_1;y_2]$ \textbf{avant} d'être transformée par la matrice du graphe.
      \end{itemize}
    \end{itemize}
\end{itemize}

\subsubsection{Fonctions affines par morceaux : Daffinebypiece}

\begin{itemize}
\item La fonction \cmd{ld.affinebypiece(def, discont)} fait le calcul des points de la courbe et renvoie une ligne polygonale (pas de dessin).

  \begin{itemize}
  \item  L'argument \argu{def} permet de définir la fonction en escaliers, c'est une table à deux champs :

\begin{TeXcode}
 { {x1,x2,x3,...,xn}, { {a1,b1}, {a2,b2},...} }
\end{TeXcode}

  Le premier élément \code{\{x1,x2,x3,\ldots,xn\}} doit être une subdivision du segment \([x_1;x_n]\).
  
  Le deuxième élément \code{\{ \{a1,b1\},\{a2,b2\},\ldots\}} signifie que sur $[x_1;x_2]$ la fonction est \(x\mapsto a_1x+b_1\), sur $[x_2;x_3]$ la fonction est \(x\mapsto a_2x+b_2\), etc.
  
  \item L'argument \argu{discont} est un booléen qui vaut \true par défaut.
  \end{itemize}
  
\item La méthode \cmd{g:Daffinebypiece(def, options)} fait le calcul des points et le dessin de la courbe de la fonction affine par morceaux.

  \begin{itemize}
  \item L'argument \argu{def} est le même que celui décrit au-dessus (définition de la fonction affine par morceaux).
  \item L'argument \argu{options} est une table dont les champs sont les options possibles. Celles-ci sont, avec leur valeur par défaut:
  \begin{itemize}
      \item \opt{discont=true},
      \item \opt{draw\_options=""},  chaîne qui sera transmise telle quelle à l'instruction \drawcmd,
      \item \opt{clip=nil}, cette option est soit \nil (valeur par défaut), soit une table \code{\{x1,x2,y1,y2\}}, dans le premier cas la ligne est clippée par la fenêtre 2D courante \textbf{après} sa transformation par la matrice 2D du graphe, dans le second cas la ligne est clippée par la fenêtre $[x_1;x_2]\times[y_1;y_2]$ \textbf{avant} d'être transformée par la matrice du graphe.
  \end{itemize}
\end{itemize}
\end{itemize}

\subsubsection{Équations différentielles : Dodesolve}

\begin{itemize}
\item La fonction \cmd{ld.odesolve(f, t0, Y0, tmin, tmax, nbdots \fac{, method})} permet une résolution approchée de l'équation différentielle \(Y'(t)=f(t,Y(t))\) dans l'intervalle $[t_{\mathrm{min}};t_{\mathrm{max}}]$ qui doit contenir \argu{t0}, avec la condition initiale $Y(t_0)=Y_0$.

\begin{itemize}
  \item L'argument \argu{f} est une fonction \(f\colon (t,Y) \mapsto f(t,Y)\) à valeurs dans \(\mathbf R^n\) et où $Y$ est également dans \(mathbf R^n\) : $Y=\{y_1, y_2,\ldots, y_n\}$, mais lorsque $n=1$, $Y$ est un réel.
  \item Les arguments \argu{t0} et \argu{Y0} donnent les conditions initiales avec $Y_0=\{y_1(t_0), \ldots, y_n(t_0)\}$ (les $y_i$ sont réels), ou $Y_0=y_1(t_0)$ lorsque $n=1$.
  \item Les arguments \argu{tmin} et \argu{tmax} définissent l'intervalle de résolution, celui-ci doit contenir \argu{t0}.
  \item L'argument \argu{nbdots} indique le nombre de points calculés de part et d'autre de \argu{t0}.
  \item L'argument optionnel \argu{method} est une chaîne qui peut valoir \val{"rkf45"} (valeur par défaut), ou \val{"rk4"}. Dans le premier cas, on utilise la méthode de Runge Kutta-Fehlberg (à pas variable), dans le second cas c'est la méthode classique de Runge-Kutta d'ordre $4$.
  \item En sortie, la fonction renvoie la matrice suivante (liste de listes de réels) :
  
\begin{TeXcode}
{ {tmin,...,tmax}, {y1(tmin),...,y1(tmax)}, {y2(tmin),...,y2(tmax)},...}
\end{TeXcode}

  La première composante est la liste des valeurs de $t$ (dans l'ordre croissant), la deuxième est la liste des valeurs (approchées) de la composante $y_1$ correspondant à ces valeurs de $t$, etc.
    \end{itemize}
    
\item La méthode \cmd{g:DplotXY(X, Y \fac{, draw\_options, clip})}, où les arguments \argu{X} et \argu{Y} sont deux listes de réels de même longueur, dessine la ligne polygonale constituée des points $(X[k],Y[k])$. L'argument \argu{draw\_options} est une chaîne (vide par défaut) qui sera passée telle quelle à l'instruction \drawcmd. L'argument \argu{clip} est soit \nil (valeur par défaut), soit une table \code{\{x1,x2,y1,y2\}}, dans le premier cas la ligne est clippée par la fenêtre 2D courante \textbf{après} sa transformation par la matrice 2D du graphe, dans le second cas la ligne est clippée par la fenêtre $[x_1;x_2]\times[y_1;y_2]$ \textbf{avant} d'être transformée par la matrice du graphe.


\begin{demo}{Un système différentiel de Lokta-Volterra}
\begin{luadraw}{name=lokta_volterra}
local ld = luadraw
local cpx = ld.cpx
local g = ld.graph:new{window={-5,50,-0.5,5},size={10,10,0}, border=true}
local i = cpx.I
local f = function(t,y) return {y[1]-y[1]*y[2],-y[2]+y[1]*y[2]} end
g:Labelsize("footnotesize")
g:Daxes({0,10,1},{limits={{0,50},{0,4}}, nbsubdiv={4,0}, legendsep={0.1,0},
originpos={"center","center"}, legend={"$t$",""}})
local y0 = {2,2}
local M = ld.odesolve(f,0,y0,0,50,250) -- résolution approchée
-- M est une table à 3 éléments: t, x et y
g:Lineoptions("solid","blue",8)
g:Dseg({5+3.5*i,10+3.5*i}); g:Dlabel("$x$",10+3.5*i,{pos="E"})
g:DplotXY(M[1],M[2]) -- points (t,x(t))
g:Linecolor("red"); g:Dseg({5+3*i,10+3*i}); g:Dlabel("$y$",10+3*i,{pos="E"})
g:DplotXY(M[1],M[3]) -- points (t,y(t))
g:Lineoptions(nil,"black",4)
g:Saveattr(); g:Viewport(20,50,3,5) -- changement de vue
g:Coordsystem(-0.5,3.25,-0.5,3.25) -- nouveau repère associé
g:Daxes({0,1,1},{legend={"$x$","$y$"},arrows="->"})
g:Lineoptions(nil,"ForestGreen",8); g:DplotXY(M[2],M[3]) -- points (x(t),y(t))
g:Restoreattr() -- retour à l'ancienne vue
g:Dlabel("$\\begin{cases}x'=x-xy\\\\y'=-y+xy\\end{cases}$", 5+4.75*i,{})
g:Show()
\end{luadraw}
\end{demo}
    
\item La méthode \cmd{g:Dodesolve(f, t0, Y0, options)} permet le dessin d'une solution à l'équation \(Y'(t)=f(t,Y(t))\).
  \begin{itemize}
  \item L'argument obligatoire \argu{f} est une fonction \(f\colon (t,Y) \mapsto f(t,Y)\) à valeurs dans \(\mathbf R^n\) et où $Y$ est également dans \(\mathbf R^n\) : $Y=\{y_1,y_2,\ldots,y_n\}$,  mais lorsque $n=1$, $Y$ est un réel.
  \item Les arguments \argu{t0} et \argu{Y0} donnent les conditions initiales avec $Y_0=\{y_1(t_0),\ldots,y_n(t_0)\}$ (les $y_i$ sont réels), ou $Y_0=y_1(t_0)$ lorsque $n=1$.
  \item L'argument \argu{options} est une table dont les champs sont les options possibles. Celles-ci sont, avec leur valeur par défaut:
      \begin{itemize}
        \item \opt{t=\{g:Xinf(), g:Xsup()\}}, détermine l'intervalle pour la variable $t$,
        \item \opt{out=\{1,2\}},  table de deux entiers $\{i_1,i_2\}$, si $M$ désigne la matrice renvoyée par la fonction \code{odesolve}, les points dessinés auront pour abscisses les $M[i_1]$ et pour ordonnées les $M[i_2]$. Par défaut on a $i_1=1$ et $i_2=2$, ce qui correspond à la fonction $y_1$ en fonction de $t$,
        \item \opt{nbdots=50}, détermine le nombre de points à calculer pour la fonction,
        \item \opt{method="rkf45"}, détermine la méthode à utiliser, les valeurs possibles sont \val{"rkf45"} ou \val{"rk4"},
        \item \opt{draw\_options=""}, est une chaîne qui sera passée telle quelle à l'instruction \drawcmd,
      \item \opt{clip=nil}, cette option est soit \nil (valeur par défaut), soit une table \code{\{x1,x2,y1,y2\}}, dans le premier cas la ligne est clippée par la fenêtre 2D courante \textbf{après} sa transformation par la matrice 2D du graphe, dans le second cas la ligne est clippée par la fenêtre $[x_1;x_2]\times[y_1;y_2]$ \textbf{avant} d'être transformée par la matrice du graphe.
      \end{itemize}
  \end{itemize}
\end{itemize}

\subsubsection{Courbes implicites : Dimplicit}

\begin{itemize}
\item La fonction \cmd{ld.implicit(f, x1, x2, y1, y2, grid)} calcule et renvoie une ligne polygonale constituant la courbe implicite d'équation $f(x,y)=0$ dans le pavé $[x_1,x_2]\times[y_1,y_2]$. Ce pavé est découpé en fonction de l'argument \argu{grid}.

      \begin{itemize}
      \item L'argument \argu{f} est une fonction \(f\colon (x,y) \mapsto f(x,y)\) à valeurs dans \(\mathbf R\).
      \item Les arguments \argu{x1}, \argu{x2} ,\argu{y1}, \argu{y2} définissent la fenêtre du tracé, qui sera le pavé $[x_1,x_2]\times[y_1,y_2]$, on doit avoir \(x_1<x_2\) et \(y_1<y_2\).
      \item L'argument \argu{grid} est une table contenant deux entiers positifs : $\{n_1,n_2\}$, le premier entier indique le nombre de subdivisions suivant l'axe des $x$, et le second le nombre de subdivisions suivant l'axe des $y$.
      \end{itemize}
  
\item La méthode \cmd{g:Dimplicit(f, options)} fait le dessin de la courbe implicite d'équations $f(x,y)=0$.
      \begin{itemize}
      \item L'argument \argu{f} est une fonction \(f\colon (x,y) \mapsto f(x,y)\) à valeurs dans \(\mathbf R\).
      \item L'argument \argu{options} est une table dont les champs sont les options possibles. Celles-ci sont, avec leur valeur par défaut:

            \begin{itemize}
                \item \opt{view=\{g:Xinf(), g:Xsup(), g:Yinf(), g:Ysup()\}}, table de la forme \code{\{x1,x2,y1,y2\}} qui détermine la zone de dessin $[x_1,x_2]\times[y_1,y_2]$.
                \item \opt{grid=\{50,50\}}, détermine les subdivisions,
                \item \opt{draw\_options=""}, est une chaîne qui sera passée telle quelle à l'instruction \drawcmd.
            \end{itemize}
      \end{itemize}
\end{itemize}

\subsubsection{Courbes de niveau : Dcontour}

La méthode \cmd{g:Dcontour(f, z, options)} fait le dessin de \textbf{lignes de niveau} de la fonction \(f\colon (x,y) \mapsto f(x,y)\) à valeurs réelles.

\begin{itemize}
  \item L'argument \argu{f} est la fonction.
  \item L'argument \argu{z} est la liste des différents niveaux à tracer.
  \item L'argument \argu{options} est une table dont les champs sont les options possibles. Celles-ci sont, avec leur valeur par défaut:
            \begin{itemize}
                \item \opt{view=\{g:Xinf(), g:Xsup(), g:Yinf(), g:Ysup()\}}, table de la forme \code{\{x1,x2,y1,y2\}} qui détermine la zone de dessin $[x_1,x_2]\times[y_1,y_2]$.
                \item \opt{grid=\{50,50\}}, détermine les subdivisions,
                \item \opt{colors \{\}}, la liste des couleurs par niveau, par défaut cette liste est vide et c'est la couleur courante de tracé qui est utilisée.
                \item \opt{draw\_options=""}, est une chaîne qui sera passée telle quelle à l'instruction \drawcmd.
            \end{itemize}  
\end{itemize}

\begin{demo}{Exemple avec Dcontour}
\begin{luadraw}{name=Dcontour}
local ld = luadraw
local cpx = ld.cpx
local g = ld.graph:new{window={-1,6.5,-1.5,11},size={10,10,0}}
local i, sin, cos = cpx.I, math.sin, math.cos
local f = function(x,y) return (x+y)/(2+cos(x)*sin(y)) end
local Lz = ld.range(1,10) -- niveaux à tracer
local Colors = ld.getpalette(ld.palRainbow,10) -- 10 couleurs équiréparties dans la palette ld.palRainbow
g:Dgradbox({0,5+10*i,1,1},{legend={"$x$","$y$"}, grid=true, title="$z=\\frac{x+y}{2+\\cos(x)\\sin(y)}$"})
g:Linewidth(12); g:Dcontour(f,Lz,{view={0,5,0,10}, colors=Colors})
for k = 1, 10 do
    local y = (2*k+4)/3*i
    g:Dseg({5.25+y,5.5+y},1,"color="..Colors[k])
    g:Labelcolor(Colors[k])
    g:Dlabel("$z="..k.."$",5.5+y,{pos="E"})
end
g:Show()
\end{luadraw}
\end{demo}

\subsubsection{Paramétrisation d'une ligne polygonale : \emph{curvilinear\_param}}
Soit $L$ une liste de complexe représentant une ligne \og continue \fg, il est possible d'obtenir une paramétrisation de cette ligne en fonction d'un paramètre $t$ entre $0$ et $1$ ($t$ est l'abscisse curviligne divisée par la longueur totale de $L$).

La fonction \cmd{ld.curvilinear\_param(L \fac{, close})} renvoie une fonction d'une variable $t\in[0;1]$ et à valeurs sur la ligne \argu{L} (liste de nombres complexes), la valeur en $t=0$ est le premier point de \argu{L}, et la valeur en $t=1$ est le dernier point; cette fonction est suivie d'un nombre qui représente la longueur total de \argu{L}. L'argument optionnel \argu{close} indique si la ligne \argu{L} doit être refermée (\false par défaut).

\begin{demo}{Points répartis sur une ligne polygonale}
\begin{luadraw}{name=curvilinear_param}
local ld = luadraw
local cpx = ld.cpx
local g = ld.graph:new{bbox=false,size={10,10}}
local i = cpx.I; g:Linewidth(8)
local L = ld.cartesian(math.sin,-5,5)[1]
ld.insert(L, {5-2*i, -5-2*i})
local f = ld.curvilinear_param(L, true)
local I = ld.map(f,ld.linspace(0,1,20)) -- 20 points répartis sur L
g:Shift(4*i)
g:Lineoptions(nil,"ForestGreen",6); g:Dpolyline(L,true)
g:Filloptions("full","white"); g:Ddots(I) -- le premier et le dernier point sont confondus car L est fermée
-- autre exemple d'utilisation:
local nb = 16 -- nombre de flèches
local t = ld.linspace(0,1,3*nb+1)
g:Shift(-4*i)
for k = 0,nb-1 do
g:Dparametric(f,{t={t[3*k+1],t[3*k+3]},nbdots=10,nbdiv=2,draw_options="-stealth"})
end
g:Show()
\end{luadraw}
\end{demo}


\subsection{Domaines liés à des courbes cartésiennes}

\subsubsection{Ddomain1}

\begin{itemize}
    \item La fonction \cmd{ld.domain1(f, a, b \fac{, nbdots, discont, nbdiv})} renvoie une liste de complexes qui représente le contour de la partie du plan délimitée par la courbe de la fonction \argu{f} sur un intervalle déterminé par \argu{a} et \argu{b}, l'axe $Ox$, et les droites \(x=a\), \(x=b\).

    \item La méthode \cmd{g:Ddomain1(f, options)} dessine ce contour. L'argument \argu{options} est une table dont les champs sont les options possibles. Celles-ci sont, avec leur valeur par défaut:
  \begin{itemize}
      \item \opt{x=\{g:Xinf(),g:Xsup()\}}, c'est l'intervalle pour le paramètre $x$ (par défaut, c'est tout l'intervalle des abscisses de la fenêtre), 
      \item \opt{nbdots=40}, 
      \item \opt{discont=false},
      \item \opt{nbdiv=5},
      \item \opt{draw\_options=""}, chaîne qui sera transmise telle quelle à l'instruction \drawcmd.
  \end{itemize}
\end{itemize}
    
\subsubsection{Ddomain2}

\begin{itemize}
    \item La fonction \cmd{ld.domain2(f, g, a, b \fac{, nbdots, discont, nbdiv})} renvoie une liste de complexes qui représente le contour de la partie du plan délimitée par la courbe de la fonction \argu{f}, la courbe de la fonction \argu{g}, et les droites \(x=a\), \(x=b\).
    
    \item La méthode \cmd{g:Ddomain2(f, g, options)} dessine ce contour. L'argument \argu{options} est une table dont les champs sont les options possibles. Celles-ci sont, avec leur valeur par défaut:
  \begin{itemize}
      \item \opt{x=\{g:Xinf(),g:Xsup()\}}, c'est l'intervalle pour le paramètre $x$ (par défaut, c'est tout l'intervalle des abscisses de la fenêtre), 
      \item \opt{nbdots=40}, 
      \item \opt{discont=false},
      \item \opt{nbdiv=5},
      \item \opt{draw\_options=""}, chaîne qui sera transmise telle quelle à l'instruction \drawcmd.
  \end{itemize}
\end{itemize}
  
\subsubsection{Ddomain3}

\begin{itemize}
    \item La fonction \cmd{ld.domain3(f, g, a, b \fac{, nbdots, discont, nbdiv})} renvoie une liste de complexes qui représente le contour de la partie du plan délimitée par la courbe de la fonction \argu{f} et celle de la fonction \argu{g} avec recherche des points d'intersection dans l'intervalle défini par \argu{a} et \argu{b}.
 
    \item La méthode \cmd{g:Ddomain3(f, g, options)} dessine ce contour. L'argument \argu{options} est une table dont les champs sont les options possibles. Celles-ci sont, avec leur valeur par défaut:
  \begin{itemize}
      \item \opt{x=\{g:Xinf(),g:Xsup()\}}, c'est l'intervalle pour le paramètre $x$ (par défaut, c'est tout l'intervalle des abscisses de la fenêtre), 
      \item \opt{nbdots=40}, 
      \item \opt{discont=false},
      \item \opt{nbdiv=5},
      \item \opt{draw\_options=""}, chaîne qui sera transmise telle quelle à l'instruction \drawcmd.
  \end{itemize}
\end{itemize}

\begin{demo}{Partie entière, fonctions Ddomain1 et Ddomain3}
\begin{luadraw}{name=courbe}
local ld = luadraw
local Z = ld.cpx.Z
local g = ld.graph:new{ window={-5,5,-5,5}, bg="", size={10,10} }
local f = function(x) return (x-2)^2-2 end
local h = function(x) return 2*math.cos(x-2.5)-2.25 end
g:Daxes( {0,1,1},{grid=true,gridstyle="dashed", arrows="->"})
g:Filloptions("full","brown",0.3)
g:Ddomain1( math.floor, { x={-2.5,3.5} })
g:Filloptions("none","white",1); g:Lineoptions("solid","red",12)
g:Dstepfunction( {ld.range(-5,5), ld.range(-5,4)},{draw_options="arrows={Bracket-Bracket[reversed]},shorten >=-2pt"})
g:Labelcolor("red")
g:Dlabel("Partie entière",Z(-3,3),{node_options="fill=white"})
g:Ddomain3(f,h,{draw_options="fill=blue,fill opacity=0.6"})
g:Dcartesian(f, {x={0,5}, draw_options="blue"})
g:Dcartesian(h, {x={0,5}, draw_options="green"})
g:Show()
\end{luadraw}
\end{demo}

\subsubsection{Dinequalities}

\begin{itemize}
    \item La fonction \cmd{ld.inequalities(constraints, x1, x2, y1, y2)} renvoie une liste de nombres complexes représentant le contour de la partie du plan située dans le pavé $[x_1;x_2]\times[y_1;y_2]$ et vérifiant les contraintes exprimées dans l'argument \argu{contraints}, celui-ci est une liste de la forme \code{\{f1,sg1,f2,sg2,...fn,sgn\}} où les \argu{fi} sont des fonctions ($f_i\colon x \mapsto f_i(x)\in \mathbf R$), et les \argu{sgi} sont soit \val{">"}, soit \val{"<"}, la première contrainte est donnée par \argu{f1} et \argu{sg1}, cette contrainte est soit $y>f_1(x)$, soit $y<f_1(x)$ en fonction de la valeur de \argu{sg1}. Il en va de même pour les contraintes suivantes.
 
    \item La méthode \cmd{g:Dinequalities(constraints, options)} permet de dessiner ce contour. L'argument \argu{options} est une table dont les champs sont les options possibles. Celles-ci sont, avec leur valeur par défaut:
  \begin{itemize}
      \item \opt{view=\{g:Xinf(),g:Xsup(),g:Yinf(),g:Ysup()\}}, c'est la fenêtre dans laquelle va se faire la résolution. Par défaut c'est la fenêtre 2D courante.
      \item \opt{draw\_options=""}, chaîne qui sera transmise telle quelle à l'instruction \drawcmd.
      \item \opt{useclip=\false}, avec la valeur \false la méthode utilise la fonction précédente (qui calcule le contour de la solution), avec la valeur \true la méthode ne calcule pas le contour mais utilise des "clipping" (un par contrainte).
  \end{itemize}
\end{itemize}

\begin{demo}{Dinequalities}
\begin{luadraw}{name=Dinequalities}
local ld = luadraw
local Z = ld.cpx.Z
local g = ld.graph:new{window={-2,5,-4,9}, size={10,10}}
g:Daxes({0,1,1}, {grid=true, arrows="-Stealth", legend={"$x$","$y$"}})
local f = function(x) return x*x - 2*x - 2 end
local h = function(x) return x*x + 2 end
local k = function(x) return -x*x+5*x + 2 end
g:Dinequalities({f,'>', h,'<', k,'<'}, {draw_options="draw=none,fill=pink,fill opacity=0.6"})
g:Dcartesian(h, {draw_options="blue, line width=0.8pt"}); g:Dlabel("$C_h$",Z(3,8.75),{node_options="blue"})
g:Dcartesian(f, {draw_options="red, line width=0.8pt"}); g:Dlabel("$C_f$",Z(2,-2),{pos="E",node_options="red"})
g:Dcartesian(k, {draw_options="green, line width=0.8pt"}); g:Dlabel("$C_k$",Z(1,7),{node_options="green"})
g:Show()
\end{luadraw}
\end{demo}

\subsubsection{Dimplicit\_inequalities}

\begin{itemize}
    \item La fonction \cmd{ld.implicit\_inequality(f, sg, x1, x2, y1, y2, grid)} renvoie un \textbf{chemin} représentant le contour de la partie du plan située dans le pavé $[x_1;x_2]\times[y_1;y_2]$ et vérifiant la condition $f(x,y)>=0$ si \argu{sg}=">", ou $f(x,y)<=0$ si \argu{sg}="<". L'argument \argu{f} est une fonction à deux variables \argu{f}$\colon (x,y)\mapsto f(x,y)\in \mathbf R$. L'argument \argu{grid} est une table de deux entiers : \argu{grid}=\{\argu{n1},\argu{n2}\}, le premier indique le nombre de subdivisions de l'intervalle $[x_1;x_2]$, et le second, le nombre de subdivisions de l'intervalle $[y_1;y_2]$. 
 
    \item La méthode \cmd{g:Dimplicit\_inequalities(constraints, options)} permet de remplir une partie du plan vérifiant les contraintes exprimées dans l'argument \argu{contraints}, celui-ci est une liste de la forme \code{\{f1,sg1,f2,sg2,...fn,sgn\}} où les \argu{fi} sont des fonctions ($f_i\colon (x,y) \mapsto f_i(x,y)\in \mathbf R$), et les \argu{sgi} sont soit \val{">"}, soit \val{"<"}, la première contrainte est donnée par \argu{f1} et \argu{sg1}, cette contrainte est soit $f_1(x,y)>0$, soit $f_1(x,y)<0$ en fonction de la valeur de \argu{sg1}. L'argument \argu{options} est une table dont les champs sont les options possibles. Celles-ci sont, avec leur valeur par défaut:
    
  \begin{itemize}
      \item \opt{view=\{g:Xinf(),g:Xsup(),g:Yinf(),g:Ysup()\}}, c'est la fenêtre dans laquelle va se faire la résolution. Par défaut c'est la fenêtre 2D courante.
      
      \item \opt{draw\_options=""}, chaîne qui sera transmise telle quelle à l'instruction \drawcmd.
      
      \item \opt{grid=\{50,50\}} : table de deux entiers, le premier indique le nombre de subdivisions de l'intervalle des $x$, et le second, le nombre de subdivisions de l'intervalle ds $y$.
  \end{itemize}
\end{itemize}

\textbf{NB}: la méthode ne calcule pas le contour mais utilise des "clipping" (un par contrainte).

\begin{demo}{La méthode \cmd{g:Dimplicit\_inequalities}}
\begin{luadraw}{name=implicit_inequalities}
local ld = luadraw
local cpx = ld.cpx
local Z = cpx.Z
local g = ld.graph:new{ window = {-3, 4, -3, 5}, size = {10, 10} }
local f1 = function(x,y) return -x*y+y^2-x^2-1  end
local f2 = function(x,y) return x^3-x-y^2+4 end
local filloptions = "draw=none,pattern= north west lines, pattern color=cyan"
g:Daxes()
g:Dimplicit_inequalities( {f1,'<',f2,">"}, {view={-2,5,-5,5}, draw_options=filloptions}) 
g:Dimplicit(f1, {draw_options="red,thick"})
g:Dimplicit(f2, {view={-2,5,-5,5}, draw_options="violet,thick"})
local eq = \luastringO{$\begin{cases}y^2-x^2-xy < 1\\ x^3-x+4 > y^2\end{cases}$}
g:Dlabel( eq, Z(2.25,1), {node_options="fill=white"})
g:Show()
\end{luadraw}
\end{demo}


\subsection{Points (Ddots) et labels (Dlabel)}

\begin{itemize}
\item La méthode pour dessiner un ou plusieurs points est : \cmd{g:Ddots(dots \fac{, mark\_options})}.

    \begin{itemize}
    \item L'argument \argu{dots} peut être soit un seul point (donc un complexe), soit une liste (une table) de complexes, soit une liste de listes de complexes. Les points sont dessinés dans la couleur courante du tracé de lignes.
    \item L'argument \argu{mark\_options} est une chaîne de caractères (vide par défaut) qui sera passée telle quelle à l'instruction \drawcmd (modifications locales), exemple :
\begin{TeXcode}
    mark_options = "color=green, line width=1.2, scale=0.25"
\end{TeXcode}

    \item Deux méthodes pour modifier globalement l'apparence des points :
        \begin{itemize}
        \item La méthode \cmd{g:Dotstyle(style)} qui définit le style de point, l'argument \argu{style} est une chaîne de caractères qui vaut par défaut \val{"*"}. Les styles possibles sont ceux de la librairie \emph{plotmarks}.
        \item La méthode \cmd{g:Dotscale(scale)} permet de jouer sur la taille du point, l'argument \argu{scale} est un entier positif qui vaut $1$ par défaut, il sert à multiplier la taille par défaut du point. La largeur courante de tracé de ligne intervient également dans la taille du point. Pour les style de points "pleins" (par exemple le style \val{"triangle*"}), le style et la couleur de remplissage courants sont utilisés par la librairie.
        \end{itemize}
    \end{itemize}
    
\item La méthode pour placer un label est : 
\cmdln{g:Dlabel(text1, anchor1, options1, text2, anchor2, options2, \ldots)}

    \begin{itemize}
    \item Les arguments \argu{text1}, \argu{text2}, \ldots, sont des chaînes de caractères, ce sont les labels.
    \item Les arguments \argu{anchor1}, \argu{anchor2}, \ldots, sont des complexes représentant les points d'ancrage des labels.
    \item  Les arguments \argu{options1}, \argu{options2}, \ldots,  permettent de définir localement les paramètres des labels, ces paramètres sont des tables dont les champs définissent les options, qui sont:
        \begin{itemize}
            \item \opt{pos="center"}, indique la position du label par rapport au point d'ancrage, il peut valoir \val{"center"} (centré), \val{"N"} (nord), \val{"NE"} (nord est), \val{"E"} (est), \val{"SE"} (sud est), \val{"S"} (sud), \val{"SW"} (sud ouest), \val{"W"} (ouest), \val{"NW"} (nord ouest). Par défaut, il vaut \val{center}, et dans ce cas le label est centré sur le point d'ancrage,
            \item \opt{dist=0}, distance en cm entre le label et son point d'ancrage lorsque \opt{pos} n'est pas égal à \val{"center"},
            \item \opt{dir= \{\}}, cette option est une table de la forme \code{\{dirX \fac{,dirY,dep}\}} qui indique la direction de l'écriture. Les 3 valeurs \code{dirX}, \code{dirY} et \code{dep} sont trois complexes représentant 3 vecteurs, les deux premiers indiquent le sens de l'écriture, le troisième un déplacement (translation) du label par rapport au point d'ancrage. Le vecteur \code{dep} est nul par défaut, et le vecteur \code{dirY}, s'il est absent, est égal au vecteur \code{dirX} tourné de 90 degrés dans le sens direct.
            Par défaut l'option \opt{dir} est égale à la valeur courante de la direction de l'écriture,
            \item \opt{node\_options=""} est une chaîne (vide par défaut) destinée à recevoir des options qui seront directement passées à TikZ dans l'instruction \emph{node[]}.
            \item Les labels sont dessinés dans la couleur courante du texte du document, mais on peut changer de couleur avec l'option \opt{node\_options} en mettant par exemple : \opt{node\_options="color=blue"}.
            
            \textbf{Attention} : les options choisies pour un label s'appliquent aussi aux labels suivants si elles sont inchangées.
        \end{itemize}
  \end{itemize}

Options globales pour les labels :

    \begin{itemize}
        \item la méthode \cmd{g:Labelstyle(position)} permet de préciser la position des labels par rapport aux points d'ancrage. L'argument \argu{position} est une chaîne qui peut valoir : \val{"N"} pour  nord, \val{"NE"} pour nord-est, \val{"NW"} pour nord-ouest, ou  encore \val{"S"}, \val{"SE"}, \val{"SW"}. Par défaut, il vaut \val{"center"}, et dans ce cas le label est centré sur le point  d'ancrage.
  
    \item La méthode \cmd{g:Labelcolor(color)} permet de définir la couleur des labels. L'argument \argu{color} est une chaîne représentant une couleur pour TikZ. Par défaut l'argument est une chaîne vide ce qui représente la couleur courante du document.
    
    \item La méthode \cmd{g:Labelangle(angle)} permet de préciser un \argu{angle} (en degrés) de rotation des labels autour du point d'ancrage, cet angle est nul par défaut.
    
    \item La méthode \cmd{g:Labelsize(size)} permet de gérer la taille des labels. L'argument \argu{size} est une chaîne qui peut valoir: \val{"tiny"}, ou \val{"scriptsize"} ou \val{"footnotesize"}, etc. Par défaut l'argument est une chaîne vide, ce qui représente la taille \val{"normalsize"}.
    
    \item La méthode \val{g:Labeldir(dir)} permet de gérer la direction l'écriture. l'argument \argu{dir} est une table de la forme \code{\{dirX \fac{,dirY,dep}\}}, les 3 valeurs \code{dirX}, \emph{dirY} et \code{dep} sont trois complexes représentant 3 vecteurs, les deux premiers indiquent le sens de l'écriture, le troisième un déplacement (translation) du label par rapport au point d'ancrage. Le vecteur \code{dep} est nul par défaut, et le vecteur \code{dirY}, s'il est absent, est égal au vecteur \code{dirX} tourné de 90 degrés dans le sens direct. Lorsque \argu{dir} est une liste vide, cela représente le sens usuel de l'écriture.
  \end{itemize}
  
  
\item La méthode \cmd{g:Dlabeldot(text, anchor, options)} permet de placer un label et de dessiner le point d'ancrage en même temps.
    \begin{itemize}
    \item L'argument \argu{texte} est une chaîne de caractères, c'est le label,
    \item L'argument \argu{anchor} est un complexe représentant le point d'ancrage du label,
    \item L'argument \argu{options} est une table dont les champs sont les options possibles. Celles-ci sont à celles de la méthode \cmd{Dlabel}, plus l'option \opt{mark\_options=""} qui est une chaîne de caractères qui sera passée telle quelle à l'instruction \drawcmd lors du dessin du point d'ancrage.
    \end{itemize}
\end{itemize}

\subsection{Chemins : Dpath, Dspline et Dtcurve}\label{Dpath}

\subsubsection{Qu'est ce qu'un chemin}

Un chemin est une table de nombres complexes et d'instructions (sous forme de chaînes), cette table représente une succession de différents "morceaux", chaque morceau est une succession de données (points 2D et parfois valeurs numériques) et se termine par une chaîne de caractères qui représente une instruction. Le chemin est régi par la règle suivante: \par
\hfil\textbf{le dernier point d'un morceau est le premier point du morceau suivant (il n'est donc pas répété) }\hfil\par

\paragraph{Exemple:}
\begin{Luacode}
local Z = cpx.Z
local L = { Z(-3,2),"m",-3,-2,"l",0,2,2,-1,"ca",3,Z(3,3),0.5,"la",1,Z(-1,5),Z(-3,2),"b" } 
\end{Luacode}
Le chemin \emph{L} est composé de cinq morceaux, qui sont:
    \begin{enumerate}
        \item \code{\{Z(-3,2),"m"\}}: il y a une donnée et l'instruction \code{"m"} qui signifie \emph{moveto}, cette instruction ne fait pas de dessin proprement dit, mais elle permet de commencer une nouvelle composante dont le premier point est $Z(-3,2)$ (le dernier point du morceau précédent, s'il y en a un, n'est pas pris en compte par cette instruction, c'est la seule exception).
        
        \item \code{\{Z(-3,2),-3,-2,"l"\}}:  le premier point du deuxième morceau est bien $Z(-3,2)$ et non pas $-3$, car $Z(-3,2)$ est le dernier point du morceau précédent. Il y a donc trois données et l'instruction "l" qui signifie \emph{lineto}, c'est comme si on exécutait l'instruction \code{g:Dpolyline({Z(-3,2),-3,-2})}, ces trois points sont donc reliés par un segment. Le dernier point de ce morceau est $-2$.
        
        \item \code{\{-2,0,2,2,-1,"ca"\}}: le premier point du troisième morceau est bien $-2$ et non pas $0$, car $-2$ est le dernier point du morceau précédent. Il y a cinq données et l'instruction \code{"ca"} qui signifie \emph{circle arc}, c'est comme si on exécutait l'instruction \code{g:Darc(-2,0,2,2,-1)}, donc le centre est $0$, l'arc va de $-2$ à $2$ avec un rayon égal à $2$ et dans le sens des aiguilles d'un montre (dernière valeur $-1$). Le dernier point de ce morceau est $2$.
        \item \code{\{2,3,Z(3,3),0.5,"la"\}}: le premier point du quatrième morceau est $2$ (et non pas $3$), il y a quatre données et l'instruction \code{"la"} qui signifie \emph{line arc}, c'est une ligne polygonale aux angles arrondis avec un arc de cercle de rayon $0.5$ (valeur qui précède l'instruction). Les points de cette ligne sont \code{\{2,3,Z(3,3)\}}, il y aura donc un arrondi en $3$. Le dernier point de ce morceau est $Z(3,3)$.
        
        \item \code{\{Z(3,3),1,Z(-1,5),Z(-3,2),"b"\}}: le premier point du cinquième morceau est $Z(3,3)$ (et non pas $1$), l'instruction \code{"b"} signifie \emph{bezier}, on dessine donc une courbe de Bézier allant de $Z(3,3)$ jusqu'à $Z(-3,2)$, les deux autres points, $1$ et $Z(-1,5)$ sont le premier et le deuxième point de contrôle de la courbe.
    \end{enumerate}

Voici ce que donne ce chemin:
\begin{demo}{Path example}
\begin{luadraw}{name=path_example}
local ld = luadraw
local Z = ld.cpx.Z
local g = ld.graph:new{window={-4,4,-0.5,3}, size={10,10}}
local L = { Z(-3,2),"m",-3,-2,"l",0,2,2,-1,"ca",3,Z(3,3),0.5,"la",1,Z(-1,5),Z(-3,2),"b" }
g:Dpath(L, "line width=0.8pt, blue")
g:Ddots({Z(-3,2),-3,-2,0,2,3,Z(3,3)})
g:Dlabel("$Z(-3,2)$",Z(-3,2),{pos="W"}, "$-3$",-3,{pos="S"}, "$-2$",-2,{},
"$0$",0,{}, "$2$",2,{}, "$3$",3,{}, "$Z(3,3)$",Z(3,3),{pos="E"})
g:Show()
\end{luadraw}
\end{demo}

\noindent\textbf{Remarque}: dans l'exemple ci-dessus, vous pouvez remplacer la partie: \code{Z(-3,2),"m",-3,-2,"l"}, par: \code{Z(-3,2),-3,-2,"l"} car il n'y a pas d'autre morceau avant le \emph{moveto}.

\noindent\textbf{Instructions disponibles et leur syntaxe}, le mot \emph{last} représente le dernier point du morceau précédent:
      \begin{itemize}
      \item \code{z1,"m"} (moveto), permet de commencer une nouvelle composante du chemin au point d'affixe $z_1$.
      \item \code{z1,...,zn,"l"} (lineto), dessine la ligne polygonale \code{\{last,z1,...,zn\}}.
      \item \code{c1,c2,z2,"b"} (bézier) dessine la courbe de Bézier \code{\{last,c1,c2,z2\}}, où $c_1$ et $c_2$ sont les deux points de contrôle.
      \item \code{z1,...,zn,"s"} (spline), dessine la spline cubique naturelle passant par les points \code{\{last,z1,...,zn\}}.
      \item \code{z1,"c"} (cercle), dessine le cercle de centre $z_1$ et passant par le point \emph{last}. Il y a une autre syntaxe possible pour le cercle: \code{z1,z2,"c"}, dans ce cas on dessine le cercle passant par les points \emph{last}, $z_1$ et $z_2$.
      \item \code{z1,z2,r,sens,"ca"} (arc de cercle), dessine un arc de cercle de centre $z_1$, de rayon $r$, allant de \emph{last} vers $z_2$, dans le sens trigonométrique lorsque \code{sens=1} (et donc dans le sens inverse si \code{sens=-1}).
      \item \code{z1,z2,rx,ry,sens,inclinaison,"ea"} (arc d'ellipse), dessine un arc d'ellipse de centre $z_1$ allant de \emph{last} vers $z_2$, $rx$ et $ry$ sont les deux rayons suivant les deux axes de l'ellipse, \code{inclinaison} est l'angle en degrés que fait le premier axe de l'ellipse (celui qui porte $rx$) avec l'horizontal. Le paramètre \code{inclinaison} est facultatif, il vaut $0$ par défaut.
      \item \code{z1,rx,ry,inclinaison,"e"} (ellipse), dessine l'ellipse de centre $z_1$, passant par \emph{last}, $rx$ et $ry$ sont les deux rayons suivant les deux axes de l'ellipse, \code{inclinaison} est l'angle en degrés que fait le premier axe de l'ellipse (celui qui porte $rx$) avec l'horizontal. Le paramètre \code{inclinaison} est facultatif, il vaut $0$ par défaut. Lorsque le point \emph{last} n'est pas sur cette ellipse, alors un segment est tracé entre ce point et un point de l'ellipse.
      \item \code{z1,...,zn,r,"la"} (line arc), dessine la ligne polygonale \code{\{last,z1,...,zn\}} en remplaçant chaque "angle" par un arc de cercle de rayon $r$.
      \item \code{z1,...,zn,r,"cla"} (closed line arc), même chose que l'instruction précédente, sauf que la ligne est refermée. 
      \item \code{"cl"} (closepath), cette instruction s'utilise seule, elle permet de refermer la composante courante en traçant un segment reliant le dernier point au premier point (de la composante courante).
      \end{itemize}
  
\subsubsection{Dessiner un chemin}

\begin{itemize}
\item La fonction \cmd{ld.path(chemin \fac{, nbdots})} renvoie une ligne polygonale contenant les points constituant le \argu{chemin}. L'argument facultatif, \argu{nbdots} est le nombre minimal de points calculés pour chaque courbe de Bézier contenue dans le chemin, sa valeur par défaut est la variable globale \varglob{ld.bezier\_nbdots} qui est initialisée à $12$.

\item La méthode \cmd{g:Dpath(chemin \fac{, draw\_options})} fait le dessin du \argu{chemin} (qui été décrit ci-dessus) en utilisant au maximum les courbes de Bézier, y compris pour les arcs, les ellipses, etc. L'argument \argu{draw\_options} est une chaîne de caractères qui sera passée directement à l'instruction \drawcmd.
 
\item La fonction \cmd{ld.spline(points \fac{, v1, v2})} renvoie sous forme de chemin (à dessiner avec \cmd{Dpath}) la spline cubique passant par les points de l'argument \argu{points} (qui doit être une liste de complexes). Les arguments \argu{v1} et \argu{v2} sont vecteurs tangents imposés aux extrémités (contraintes), lorsque ceux-ci sont égaux à \nil (valeur par défaut), c'est une spline cubique naturelle (c'est-à-dire sans contrainte) qui est calculée.

\item La méthode \cmd{g:Dspline(points \fac{, v1, v2, draw\_options})} fait le dessin de la spline décrite ci-dessus. L'argument \argu{draw\_options} est une chaîne de caractères qui sera passée directement à l'instruction \drawcmd.

\begin{demo}{Path et Spline}
\begin{luadraw}{name=path_spline}
local ld = luadraw
local g = ld.graph:new{window={-5,5,-5,5},size={10,10},bg="Beige"}
local i = ld.cpx.I
local p = {-3+2*i,"m",-3,-2,"l",0,2,2,1,"ca",3,3+3*i,0.5,"la",1,-1+5*i,-3+2*i,"b",-1,"m",0,"c"}
g:Daxes( {0,1,1} )
g:Filloptions("full","blue!30",1,true); g:Dpath(p,"line width=0.8pt")
g:Filloptions("none")
local A,B,C,D,E = -4-i,-3*i,4,3+4*i,-4+2*i
g:Lineoptions(nil,"ForestGreen",12); g:Dspline({A,B,C,D,E},nil,-5*i) -- contrainte en E
g:Ddots({A,B,C,D,E},"fill=white,scale=1.25")
g:Show()
\end{luadraw}
\end{demo}

\item La fonction \cmd{ld.tcurve(L)} renvoie sous forme de chemin une courbe passant par des points donnés avec des vecteurs tangents (à gauche et à droite) imposés à chaque point. L'argument \argu{L} est une table de la forme :
\begin{Luacode}
L = {point1,{t1,a1,t2,a2}, point2,{t1,a1,t2,a2}, ..., pointN,{t1,a1,t2,a2}}
\end{Luacode}
\argu{point1}, \ldots, \argu{pointN} sont les points d'interpolation de la courbe (affixes), et chacun d'eux est suivi d'une table de la forme \code{\{t1,a1,t2,a2\}} qui précise les vecteurs tangents à la courbe à gauche du point (avec \argu{t1} et \argu{a1}) et à droite du point (avec \argu{t2} et \argu{a2}). Le vecteur tangent à gauche est donné par la formule $V_g = t_1\times e^{ia_1\pi/180}$, donc \argu{t1} représente le module et \argu{a1} est un argument \textbf{en degrés} de ce vecteur. C'est la même chose avec \argu{t2} et \argu{a2} pour le vecteur tangent à droite, \textbf{mais ceux-ci sont facultatifs}, et s'ils ne sont pas précisés alors ils prennent les mêmes valeurs que \argu{t1} et \argu{a1}.

Deux points consécutifs seront reliés par une courbe de Bézier, la fonction calcule les points de contrôle pour avoir les vecteurs tangents souhaités.

\item La méthode \cmd{g:Dtcurve(L, options)} fait le dessin du chemin obtenu par \argu{tcurve} décrit ci-dessus. L'argument \argu{options} est une table dont les champs sont les options possibles, qui sont, avec leur valeur par défaut:
    \begin{itemize}
        \item \opt{showdots=false}, cette option permet de dessiner ou non les points d'interpolation donnés ainsi que les points de contrôles calculés, ce qui permet une visualisation des contraintes.
        \item \opt{draw\_options=""}, c'est une chaîne de caractères qui sera passée directement à l'instruction \drawcmd.
    \end{itemize}
\end{itemize}

\begin{demo}{Courbe d'interpolation avec vecteurs tangents imposés}
\begin{luadraw}{name=tcurve}
local ld = luadraw
local g = ld.graph:new{window={-0.5,10.5,-0.5,6.5},size={10,10,0}}
local i = ld.cpx.I
local L = {
1+4*i,{2,-20},
2+3*i,{2,-70},
4+i/2,{3,0},
6+3*i,{4,15},
8+6*i,{4,0,4,-90}, -- point anguleux
10+i,{3,-15}}
g:Dgrid({0,10+6*i},{gridstyle="dashed"})
g:Daxes(nil,{limits={{0,10},{0,6}},originpos={"center","center"}, arrows="->"})
g:Dtcurve(L,{showdots=true,draw_options="line width=0.8pt,red"})
g:Show()
\end{luadraw}
\end{demo}

\subsection{Chemins et clipping : Beginclip() et Endclip()}

Un chemin peut être utilisé pour faire du clipping grâce à deux fonctions : \cmd{g:Beginclip(chemin \fac{, inverse})} et \cmd{g:Endclip()}. La première ouvre un groupe \emph{scope} et passe le \argu{chemin} comme argument à la fonction \emph{\textbackslash clip} de TikZ. La seconde referme le groupe \emph{scope}, elle est indispensable (sinon il y aura une erreur de compilation).
L'argument \argu{inverse} est un booléen qui vaut \false par défaut, lorsqu'il a la valeur \true le clipping est inversé, c'est à dire que seul ce qui est à l'extérieur du \argu{chemin} sera dessiné, mais il faut pour cela que celui-ci soit dans le sens trigonométrique.

\begin{demo}{Exemple de clipping}
\begin{luadraw}{name=polygon_with_different_line_color_and_rounded_corners}
local ld = luadraw
local g = ld.graph:new{window={-5,5,-5,5},size={10,10}}
local i = ld.cpx.I
local Dcolored_polyreg = function(c,a,nb,r,wd,colors)
-- c=centre, a=sommet, nb=nombre de côtés, r=rayon, wd=épaisseur en point, colors=liste de couleurs
    local L = ld.polyreg(c,a,nb)
    ld.insert(L,{r,"cla"}) --polygone aux angles arrondis (radius=r)
    local angle = 360/nb
    local b = a
    for k = 1, nb do
        a = b; b = ld.rotate(a,angle,c)
        g:Beginclip({2*a-c,c,2*b-c,"l"}) -- définition d'un secteur angulaire pour clipper
        g:Dpath(L,"line width="..wd.."pt,color="..colors[k])
        g:Endclip()
    end
end
Dcolored_polyreg(3+2*i,5+2*i,5,0.8,12,{"red","blue","orange","green","yellow"}) -- pentagon
Dcolored_polyreg(-2.5-2*i,-5-2*i,7,1,24,ld.getpalette(ld.palGasFlame,7)) -- heptagon
g:Show()
\end{luadraw}
\end{demo}


\subsection{Axes et grilles}

Variables globales utilisées pour les axes et les grilles :
\begin{itemize}
    \item \varglob{ld.maxGrad = 100} : nombre max de graduations sur un axe.
    \item \varglob{ld.defaultlabelshift = 0.1875} : lorsqu'une grille est dessinée avec les axes (option \opt{grid=true}) les labels sont automatiquement décalés le long de l'axe avec cette variable.
    \item \varglob{ld.defaultxylabelsep = 0} : définit la distance par défaut entre les labels et les graduations.
    \item \varglob{ld.defaultlegendsep = 0.2} : définit la distance par défaut entre la légende et l'axe.
    \item \varglob{ld.digits = 4} : nombre de décimales par défaut dans les conversions en chaînes de caractères, les $0$ terminaux sont supprimés.
    \item \varglob{ld.dollar = true} :  pour ajouter des dollars autour des labels des graduations.
    \item \varglob{ld.siunitx = false} :  avec la valeur \true les labels sont formatés avec la macro \verb|\num{...}| du package \emph{siunitx}, ce qui permet d'utiliser certaines options de ce package, comme remplacer le point décimal par une virgule en faisant :\par
    \begin{TeXcode}
    \usepackage[local=FR]{siunitx}
    \end{TeXcode}
    ou bien en faisant :
    \begin{TeXcode}
    \usepackage{siunitx}
    \sisetup{output-decimal-marker={,}}
    \end{TeXcode}
\end{itemize}
Pour les axes, en 2D comme en 3D, tous les labels sont formatés en chaînes de caractères avec la fonction \cmd{ld.num(x)}, celle-ci transforme le nombre $x$ en une chaîne \emph{str} avec le nombre de décimales fixées par la variable globale \varglob{ld.digits}, lorsque la variable \varglob{ld.siunitx} a la valeur \true, la fonction renvoie \val{"\textbackslash num\{str\}"}, sinon elle renvoie simplement \emph{str}. Ceci vaut pour également pour les axes en 3D. Voici le code de cette fonction:
\begin{Luacode}
function ld.num(x) -- x is a real, returns a string
    local rep = ld.strReal(x) -- conversion to string with digits decimals max
    if ld.siunitx then rep = "\\num{"..rep.."}" end --needs \usepackage{siunitx}
    return rep
end
\end{Luacode}

\subsubsection{Daxes}
%\def\opt#1{\textcolor{blue}{\texttt{#1}}}%
Le tracé des axes s'obtient avec la méthode \cmd{g:Daxes( \fac{\{A, xpas, ypas\}, options})}.
\begin{itemize}
    \item Le premier argument précise le point d'intersection des deux axes (c'est le complexe \argu{A}), le pas des graduations sur l'axe $Ox$ (c'est \argu{xpas}) et le pas des graduations sur $Oy$ (c'est \argu{ypas}). Par défaut le point \argu{A} est l'origine $Z(0,0)$, et les deux pas sont égaux à $1$.
    \item L'argument \argu{options} est une table précisant les options possibles. Voici ces options avec leur valeur par défaut :
        \begin{itemize}
            \item \opt{showaxe=\{1,1\}}. Cette option précise si les axes doivent être tracés ou pas ($1$ ou $0$). La première valeur est pour l'axe $Ox$ et la seconde pour l'axe $Oy$.
            \item \opt{arrows="-"}. Cette option permet d'ajouter ou non une flèche aux axes (pas de flèche par défaut, mettre \val{"->"} par exemple pour ajouter une flèche).
            \item \opt{limits=\{"auto","auto"\}}. Cette option permet de préciser l'étendue des deux axes (première valeur pour $Ox$, seconde valeur pour $Oy$). La valeur "auto" signifie que c'est la droite en entier, mais on peut préciser les abscisses extrêmes, par exemple : \opt{limits=\{\{-4,4\},"auto"\}}.
            \item \opt{gradlimits=\{"auto","auto"\}}. Cette option permet de préciser l'étendue des graduations sur les deux axes (première valeur pour $Ox$, seconde valeur pour $Oy$). La valeur \val{"auto"} signifie que c'est la droite en entier, mais on peut préciser les graduations extrêmes, par exemple : \opt{gradlimits=\{\{-4,4\},\{-2,3\}\}}.
            \item \opt{unit=\{"",""\}}. Cette option permet de préciser de combien en combien vont les graduations sur les axes. La valeur par défaut signifie qu'il faut prendre la valeur du pas (\argu{xpas} sur $Ox$, ou \argu{ypas} sur $Oy$), SAUF lorsque l'option \opt{labeltext} n'est pas la chaîne vide, dans ce cas \opt{unit} prend la valeur $1$.
            \item \opt{nbsubdiv=\{0,0\}}. Cette option permet de préciser le nombre de subdivisions entre deux graduations principales sur l'axe.
            \item \opt{tickpos=\{0.5,0.5\}}. Cette option précise la position des graduations par rapport à chaque axe, ce sont deux nombres entre $0$ et $1$, la valeur par défaut de $0.5$ signifie qu'ils sont centrés sur l'axe ($0$ et $1$ représentent les extrémités).
            \item \opt{tickdir=\{"auto","auto"\}}. Cette option indique la direction des graduations sur l'axe. Cette direction est un vecteur (complexe) non nul. La valeur par défaut \val{"auto"} signifie que les graduations sont orthogonales à l'axe.
            \item \opt{xyticks=\{0.2,0.2\}}. Cette option précise la longueur des graduations sur l'axe.
            \item \opt{xylabelsep=\{0,0\}}. Cette option précise la distance entre les labels et les graduations sur l'axe.
            \item \opt{originpos=\{"right","top"\}}. Cette option précise la position du label à l'origine sur l'axe, les valeurs possibles sont : \val{"none"}, \val{"center"}, \val{"left"}, \val{"right"} pour $Ox$, et \val{"none"}, \val{"center"}, \val{"bottom"}, \val{"top"} pour $Oy$.
            
            \item \opt{originnum=\{A.re,A.im\}}. Cette option précise la valeur au point origine des graduations (graduation numéro $0$). 
            
            La formule qui définit le label à la graduation numéro $n$ est : 
            \codeln{(originnum + unit*n)"labeltext"/labelden.}
            
            \item \opt{originloc=A}. Cette option précise le point origine des graduations.
            
            \item \opt{legend=\{"",""\}}. Cette option permet de préciser une légende pour l'axe.
            \item \opt{legendpos=\{0.975,0.975\}}. Cette option précise la position (entre $0$ et $1$) de la légende par rapport à chaque axe.
            
            \item \opt{legendsep=\{ld.defaultlegendsep,ld.defaultlegendsep\}}. Cette option précise la distance entre la légende et l'axe. La légende est de l'autre côté de l'axe par rapport aux graduations.
            
            \item \opt{legendangle=\{"auto","auto"\}}. Cette option précise l'angle (en degrés) que doit faire la légende pour l'axe. La valeur \val{"auto"} par défaut signifie que la légende doit être parallèle à l'axe si l'option \opt{labelstyle} est aussi à \val{"auto"}, sinon la légende est horizontale.
            
            \item \opt{legendstyle=\{"auto","auto"\}}. Précise le style de label pour les légendes, avec la valeur \val{"auto"} celui-ci est déterminé automatiquement, sinon on peut utiliser les valeurs: \val{"N"}, \val{"NW"}, \val{"W"}, \val{"SW"}, \val{"S"}, \val{"SE"}, \val{"E"}, \val{"NE"}.
            
            \item \opt{labelpos=\{"bottom","left"\}}. Cette option précise la position des labels par rapport à l'axe. Pour l'axe $Ox$, les valeurs possibles sont : \val{"none"}, \val{"bottom"} ou \val{"top"}, pour l'axe $Oy$ c'est :  \val{"none"}, \val{"right"} ou \val{"left"}.
            
            \item \opt{labelden=\{1,1\}}. Cette option précise le dénominateur des labels (entier) pour l'axe. Rappelons que la formule qui définit le label à la graduation numéro $n$ est : 
            \codeln{(originnum + unit*n)"labeltext"/labelden.}
            
            \item \opt{labeltext=\{"",""\}}. Cette option définit le texte qui sera ajouté au numérateur des labels pour l'axe.
            
            \item \opt{labelstyle=\{"S","W"\}}. Cette option définit le style des labels pour chaque axe. Les valeurs possibles sont : \val{"auto"}, \val{"N"}, \val{"NW"}, \val{"W"}, \val{"SW"}, \val{"S"}, \val{"SE"}, \val{"E"}, \val{"NE"}.
            
            \item \opt{labelangle=\{0,0\}}. Cette option définit pour chaque axe l'angle des labels en degrés par rapport à l'horizontale.
            
            \item \opt{labelcolor=\{"",""\}}. Cette option permet de choisir une couleur pour les labels sur chaque axe. La chaîne vide représente la couleur par défaut.
            
            \item \opt{labelshift=\{0,0\}}. Cette option permet de définir un décalage systématique pour les labels sur l'axe (décalage de long de l'axe), avec l'option \opt{grid=false} les valeurs par défaut sont $0$ et $0$, mais avec l'option \opt{grid=true} on a par défaut \opt{labelshift=\{ld.defaultlabelshift,ld.defaultlabelshift\}} où \varglob{ld.defaultlabelshift} est une variable globale initialisée à $0.1875$.
            
            \item \opt{xynode\_options=""}. Chaîne de caractères qui sera transmise telle quelle à l'instruction \verb|\node{}| pour tous les labels sur les deux axes (mais pas pour les légendes).
            
            \item \opt{xnode\_options=xynode\_options}. Chaîne de caractères qui sera transmise telle quelle à l'instruction \verb|\node{}| pour tous les labels sur l'axe des $x$ (mais pas pour la légende).
            
             \item \opt{ynode\_options=xynode\_options}. Chaîne de caractères qui sera transmise telle quelle à l'instruction \verb|\node{}| pour tous les labels sur l'axe des $y$ (mais pas pour la légende).
            
            \item \opt{nbdeci=\{2,2\}}. Cette option précise le nombre de décimales pour les valeurs numériques sur l'axe.
            
            \item \opt{use\_siunitx=\{ld.siunitx,ld.siunitx\}}. Cette option précise si les valeurs numériques doivent être formatées en utilisant le paquet \emph{siunitx}, la valeur par défaut est celle de la variable globale \varglob{ld.siunitx} qui vaut \false par défaut.
            
            \item \opt{myxlabels=""}. Cette option permet d'imposer des labels personnels sur l'axe $Ox$. Lorsqu'il y en a, la valeur passée à l'option doit être une liste du type : \code{\{pos1,"text1",pos2,"text2",\ldots\}}. Le nombre \argu{pos1} représente une abscisse dans le repère \code{(A,xpas)}, ce qui correspond au point d'affixe \code{A+pos1*xpas}.
            
            \item \opt{myylabels=""}. Cette option permet d'imposer des labels personnels sur l'axe $Oy$. Lorsqu'il y en a, la valeur passée à l'option doit être une liste du type :  \code{\{pos1,"text1",pos2,"text2",\ldots\}}. Le nombre \argu{pos1} représente une abscisse dans le repère \code{(A,i*ypas)}, ce qui correspond au point d'affixe \code{A+pos1*ypas*i}.
            
            \item \opt{grid=false}. Cette option permet d'ajouter ou non une grille.
            
            \item \opt{showlines=\{\true,\true\}}. Lorsque \opt{grid=\true}, cette option permet d'afficher ou non les traits verticaux de la grille (correspondant à l'axe des $x$, ainsi que les traits horizontaux de la grille (correspondant à l'axe des $y$).
            
            \item \opt{drawbox=false}. Cette option de dessiner les axes sous la forme d'une boite, dans ce cas, les graduations sont sur le côté gauche et le côté bas.
            \item \opt{gridstyle="solid"}. Cette option définit le style ligne pour la grille principale.
            \item \opt{subgridstyle="solid"}. Cette option définit le style ligne pour la grille secondaire. Une grille secondaire apparaît lorsqu'il y a des subdivisions sur un des axes.
            \item  \opt{gridcolor="gray"}. Ceci définit la couleur de la grille principale.
            \item \opt{subgridcolor="lightgray"}. Ceci définit la couleur de la grille secondaire.
            \item \opt{gridwidth=4}. Épaisseur de trait de la grille principale (ce qui fait 0.4pt).
            \item \opt{subgridwidth=2}. Épaisseur de trait de la grille secondaire (ce qui fait 0.2pt).
        \end{itemize}
\end{itemize}

\begin{demo}{Exemple avec axes avec grille}
\begin{luadraw}{name=axes_grid}
local ld = luadraw
local g = ld.graph:new{window={-6.5,6.5,-3.5,3.5}, size={10,10,0}}
local i, pi, a = ld.cpx.I, math.pi, math.sqrt(2)
local f = function(x) return 2*a*math.sin(x) end
g:Labelsize("footnotesize"); g:Linewidth(8)
g:Daxes({0,pi/2,a},{labeltext={"\\pi","\\sqrt{2}"}, labelden={2,1},nbsubdiv={1,1},grid=true,arrows="->"})
g:Lineoptions("solid","Crimson",12); g:Dcartesian(f, {x={-2*pi,2*pi}})
g:Show()
\end{luadraw}
\end{demo}

\subsubsection{DaxeX et DaxeY}

Les méthodes \cmd{g:DaxeX(\fac{\{A, xpas\}, options})} et \cmd{g:DaxeY(\fac{\{A, ypas\}, options})} permettent de tracer les axes séparément.
\begin{itemize}
    \item Le premier argument précise le point servant d'origine (c'est le complexe \argu{A}) et le pas des graduations sur l'axe. Par défaut le point \argu{A} est l'origine $Z(0,0)$, et le pas est égal à $1$.
    \item L'argument \argu{options} est une table précisant les options possibles. Voici ces options avec leur valeur par défaut :
        \begin{itemize}
            \item \opt{showaxe=1}. Cette option précise si l'axe doit être tracé ou non ($1$ ou $0$).
            
            \item \opt{arrows="-"}. Cette option permet d'ajouter ou non une flèche à l'axe (pas de flèche par défaut, mettre \val{"->"} pour ajouter une flèche).
            
            \item \opt{limits="auto"}. Cette option permet de préciser l'étendue des deux axes. La valeur "auto" signifie que c'est la droite en entier, mais on peut préciser les abscisses extrêmes, par exemple : \opt{limits=\{-4,4\}}.
            
            \item \opt{gradlimits="auto"}. Cette option permet de préciser l'étendue des graduations sur les deux axes. La valeur \val{"auto"} signifie que c'est la droite en entier, mais on peut préciser les graduations extrêmes, par exemple : \opt{gradlimits=\{-2,3\}}.
            \item \opt{unit=""}. Cette option permet de préciser de combien en combien vont les graduations sur l'axe. La valeur par défaut signifie qu'il faut prendre la valeur du pas, SAUF lorsque l'option \opt{labeltext} n'est pas la chaîne vide, dans ce cas \opt{unit} prend la valeur $1$.
            
            \item \opt{nbsubdiv=0}. Cette option permet de préciser le nombre de subdivisions entre deux graduations principales.
            
            \item \opt{tickpos=0.5}. Cette option précise la position des graduations par rapport à l'axe, ce sont deux nombres entre $0$ et $1$, la valeur par défaut de $0.5$ signifie qu'ils sont centrés sur l'axe. ($0$ et $1$ représentent les extrémités).
            
            \item \opt{tickdir="auto"}. Cette option indique la direction des graduations sur l'axe. Cette direction est un vecteur (complexe) non nul. La valeur par défaut \val{"auto"} signifie que les graduations sont orthogonales à l'axe.
            
            \item \opt{xyticks=0.2}. Cette option précise la longueur des graduations.
            
            \item \opt{xylabelsep=0}. Cette option précise la distance entre les labels et les graduations.
            
            \item \opt{originpos="center"}. Cette option précise la position du label à l'origine sur l'axe, les valeurs possibles sont : \val{"none"}, \val{"center"}, \val{"left"}, \val{"right"} pour $Ox$, et \val{"none"}, \val{"center"}, \val{"bottom"}, \val{"top"} pour $Oy$.
            
            \item \opt{originnum=A.re} pour $Ox$ et \opt{originnum=A.im} pour $Oy$. Cette option précise la valeur de la graduation à l'origine (graduation numéro $0$). 
            
            La formule qui définit le label à la graduation numéro $n$ est : 
            \codeln{(originnum + unit*n)"labeltext"/labelden.}

            \item \opt{legend=""}. Cette option permet de préciser une légende pour l'axe.
            
            \item \opt{legendpos=0.975}. Cette option précise la position (entre $0$ et $1$) de la légende par rapport à l'axe.
            
            \item \opt{legendsep=ld.defaultlegendsep}. Cette option précise la distance entre la légende et l'axe. La légende est de l'autre côté de l'axe par rapport aux graduations.
            
            \item \opt{legendangle="auto"}. Cette option précise l'angle (en degrés) que doit faire la légende pour l'axe. La valeur "auto" par défaut signifie que la légende doit être parallèle à l'axe si l'option \opt{labelstyle} est aussi à \val{"auto"}, sinon la légende est horizontale.
            
            \item \opt{legendstyle="auto"}. Précise le style de label pour la légende, avec la valeur \val{"auto"} celui-ci est déterminé automatiquement, sinon on peut utiliser les valeurs: \val{"N"}, \val{"NW"}, \val{"W"}, \val{"SW"}, \val{"S"}, \val{"SE"}, \val{"E"}, \val{"NE"}.
            
            \item \opt{labelpos="bottom"} pour $Ox$ et \opt{labelpos="left"} pour $Oy$. Cette option précise la position des labels par rapport à l'axe. Pour l'axe $Ox$, les valeurs possibles sont : \val{"none"}, \val{"bottom"} ou \val{"top"}, pour l'axe $Oy$ c'est :  \val{"none"}, \val{"right"} ou \val{"left"}.
            
            \item \opt{labelden=1}. Cette option précise le dénominateur des labels (entier) pour l'axe. La formule qui définit le label à la graduation numéro $n$ est : 
            \codeln{(originnum + unit*n)"labeltext"/labelden.}
            
            \item \opt{labeltext=""}. Cette option définit le texte qui sera ajouté au numérateur des labels.
            
            \item \opt{labelstyle="S"} pour $Ox$ et \opt{labelstyle="W"} pour $Oy$. Cette option définit le style des labels. Les valeurs possibles sont: \val{"auto"}, \val{"N"}, \val{"NW"}, \val{"W"}, \val{"SW"}, \val{"S"}, \val{"SE"}, \val{"E"}, \val{"NE"}.
            
            \item \opt{labelangle=0}. Cette option définit l'angle des labels en degrés par rapport à l'horizontale.
            
            \item \opt{labelcolor=""}. Cette option permet de choisir une couleur pour les labels. La chaîne vide représente la couleur courante du texte.
            
            \item \opt{labelshift=0}. Cette option permet de définir un décalage systématique pour les labels sur l'axe (décalage de long de l'axe).
            
            \item \opt{nbdeci=2}. Cette option précise le nombre de décimales pour les labels numériques.
            
            \item \opt{use\_siunitx=ld.siunitx}. Cette option précise si les valeurs numériques doivent être formatées en utilisant le paquet \emph{siunitx}, la valeur par défaut est celle de la variable globale \varglob{ld.siunitx} qui vaut \false par défaut.
            
            \item \opt{mylabels=""}. Cette option permet d'imposer des labels personnels. Lorsqu'il y en a, la valeur passée à l'option doit être une liste du type : \code{\{pos1,"text1",pos2,"text2",\ldots\}}. Le nombre \argu{pos1} représente une abscisse dans le repère \code{(A,xpas)} pour $Ox$, ou \code{(A,ypas*i)} pour $Oy$, ce qui correspond au point d'affixe \code{A+pos1*xpas} pour $Ox$, et \code{A+pos1*ypas*i} pour $Oy$.
        \end{itemize}
\end{itemize}

\subsubsection{Dgradline}

Les méthodes de tracé des axes s'appuient sur la méthode \cmd{g:Dgradline(\{A, u\}, options)}, où \argu{\{A,u\}} représente la droite passant par \argu{A} (un complexe) et dirigé par le vecteur \argu{u} (un complexe non nul), le couple $(A,u)$ sert de repère sur cette droite (et oriente cette droite), donc chaque point $M$ de cette droite a une abscisse $x$ telle $M=A+xu$. Cette méthode permet de dessiner cette droite graduée, l'argument \argu{options} est une table précisant les options possibles, qui sont (avec leur valeur par défaut) :
\begin{itemize}
    \item \opt{showaxe=1}. Cette option précise si l'axe doit être tracé ou non ($1$ ou $0$).
    
    \item \opt{arrows="-"}. Cette option permet d'ajouter ou non une flèche à l'axe (pas de flèche par défaut, mettre \val{"->"} pour ajouter une flèche).
    
    \item \opt{limits="auto"}. Cette option permet de préciser l'étendue des deux axes. La valeur \val{auto"} signifie que c'est la droite en entier, mais on peut préciser les abscisses extrêmes, par exemple : \opt{limits=\{-4,4\}}.
    
    \item \opt{gradlimits="auto"}. Cette option permet de préciser l'étendue des graduations sur les deux axes. La valeur \val{"auto"} signifie que c'est la droite en entier, mais on peut préciser les graduations extrêmes, par exemple : \opt{gradlimits=\{-2,3\}}.
    
    \item \opt{unit=1}. Cette option permet de préciser de combien en combien vont les graduations sur l'axe.
    
    \item \opt{nbsubdiv=0}. Cette option permet de préciser le nombre de subdivisions entre deux graduations principales.
    
    \item \opt{tickpos=0.5}. Cette option précise la position des graduations par rapport à l'axe, ce sont deux nombres entre $0$ et $1$, la valeur par défaut de $0.5$ signifie qu'ils sont centrés sur l'axe. ($0$ et $1$ représentent les extrémités).
    
    \item \opt{tickdir="auto"}. Cette option indique la direction des graduations sur l'axe. Cette direction est un vecteur (complexe) non nul. La valeur par défaut \val{"auto"} signifie que les graduations sont orthogonales à l'axe.
    
    \item \opt{xyticks=0.2}. Cette option précise la longueur des graduations.
    
    \item \opt{xylabelsep=defaultxylabelsep}. Cette option précise la distance entre les labels et les graduations, \varglob{defaultxylabelsep} est une variable globale valant $0$ par défaut.
    
    \item \opt{originpos="center"}. Cette option précise la position du label à l'origine sur l'axe, les valeurs possibles sont : \val{"none"}, \val{"center"}, \val{"left"}, \val{"right"}.
    
    \item \opt{originnum=0}. Cette option précise la valeur de la graduation à l'origine $A$ (graduation numéro $0$).
    
    La formule qui définit le label à la graduation numéro $n$ (au point $A+nu$) est : 
    \codeln{(originnum + unit*n)"labeltext"/labelden.}

    \item \opt{legend=""}. Cette option permet de préciser une légende pour l'axe.
    
    \item \opt{legendpos=0.975}. Cette option précise la position (entre $0$ et $1$) de la légende par rapport à l'axe.
    
    \item \opt{legendsep=ld.defaultlegendsep}. Cette option précise la distance entre la légende et l'axe. La légende est de l'autre côté de l'axe par rapport aux graduations, \varglob{ld.defaultlegendsep} est une variable globale qui vaut $0.2$ par défaut.
    
    \item \opt{legendangle="auto"}. Cette option précise l'angle (en degrés) que doit faire la légende pour l'axe. La valeur \val{"auto"} par défaut signifie que la légende doit être parallèle à l'axe si l'option \opt{labelstyle} est aussi à \val{"auto"}, sinon la légende est horizontale.
    
    \item \opt{legendstyle="auto"}. Précise le style de label pour la légende, avec la valeur \val{"auto"} celui-ci est déterminé automatiquement, sinon on peut utiliser les valeurs: \val{"N"}, \val{"NW"}, \val{"W"}, \val{"SW"}, \val{"S"}, \val{"SE"}, \val{"E"}, \val{"NE"}.
    
    \item \opt{labelpos="bottom"}. Cette option précise la position des labels par rapport à l'axe, les valeurs possibles sont : \val{"none"}, \val{"bottom"}, or \val{"top"}. Cette position détermine en même temps celle de la légende : de l'autre côté de l'axe.
    
    \item \opt{labelden=1}. Cette option précise le dénominateur des labels (entier) pour l'axe. La formule qui définit le label à la graduation numéro $n$ est : 
    \codeln{(originnum + unit*n)"labeltext"/labelden.}
    
    \item \opt{labeltext=""}. Cette option définit le texte qui sera ajouté au numérateur des labels.
    
    \item \opt{labelstyle="auto"}. Cette option définit le style des labels. Les valeurs possibles sont: \val{"auto"}, \val{"N"}, \val{"NW"}, \val{"W"}, \val{"SW"}, \val{"S"}, \val{"SE"}, \val{"E"}, \val{"NE"}.
    
    \item \opt{labelangle=0}. Cette option définit l'angle des labels en degrés par rapport à l'horizontale.
    
    \item \opt{labelcolor=""}. Cette option permet de choisir une couleur pour les labels. La chaîne vide représente la couleur courante du texte.
    
    \item \opt{labelshift=0}. Cette option permet de définir un décalage systématique pour les labels sur l'axe (décalage de long de l'axe).
    
    \item \opt{nbdeci=2}. Cette option précise le nombre de décimales pour les labels numériques.
    
    \item \opt{use\_siunitx=ld.siunitx}. Cette option précise si les valeurs numériques doivent être formatées en utilisant le paquet \emph{siunitx}, la valeur par défaut est celle de la variable globale \varglob{ld.siunitx} qui vaut \false par défaut.
    
    \item \opt{mylabels=""}. Cette option permet d'imposer des labels personnels. Lorsqu'il y en a, la valeur passée à l'option doit être une liste du type : \code{\{x1,"text1",x2,"text2",\ldots\}}. Les nombres \argu{x1}, \argu{x2}, \ldots, représentent des abscisses dans le repère $(A,u)$.
\end{itemize}

\begin{demo}{Exemples de droites graduées}
\begin{luadraw}{name=gradline}
local ld = luadraw
local g = ld.graph:new{window={-5,5,-5,5},size={10,10}}
g:Labelsize("footnotesize")
local i = ld.cpx.I
g:Dgradline({3.25*i,1+i/2}, {limits={-4,4}, legend="Axe", legendpos=0.5, arrows="-stealth"})
g:Dgradline({-3,1}, {legend="demo", labeltext="\\pi", labelden=3, unit=2, nbsubdiv=1, arrows="-latex"})
g:Dgradline({3-4*i,-1.25+i/5}, {legend="A", labelstyle="N", gradlimits={-1,5},
nbsubdiv=3, unit=1.411, nbdeci=3, arrows="-Latex"})
g:Show()
\end{luadraw}
\end{demo}

\subsubsection{Dgrid}

La méthode \cmd{g:Dgrid(\{A, B\}, options} permet le dessin d'une grille.
\begin{itemize}
    \item Le premier argument est obligatoire, il précise le coin inférieur gauche (c'est le complexe \argu{A}), le coin supérieur droit (c'est le complexe \argu{B}) de la grille.
    \item L'argument \argu{options} est une table précisant les options possibles. Voici ces options avec leur valeur par défaut :
        \begin{itemize}
            \item \opt{unit=\{1,1\}}. Cette option définit les unités sur les axes pour la grille principale.
            \item \opt{showlines=\{\true,\true\}}. Cette option permet d'afficher ou non les traits verticaux de la grille (correspondant à l'axe des $x$, ainsi que les traits horizontaux de la grille (correspondant à l'axe des $y$).
            \item \opt{gridwidth=4}. Cette option définit l'épaisseur du trait de la grille principale (0.4pt par défaut).
            \item \opt{gridcolor="gray"}. Couleur grille de la grille principale.
            \item \opt{gridstyle="solid"}. Style de trait pour la grille principale.
            \item \opt{nbsubdiv=\{0,0\}}. Nombre de subdivisions (pour chaque axe) entre deux traits de la grille principale. Ces subdivisions déterminent la grille secondaire.
            \item \opt{subgridcolor="lightgray"}. Couleur de la grille secondaire.
            \item \opt{subgridwidth=2}. Épaisseur du trait de la grille secondaire (0.2pt par défaut).
            \item \opt{subgridstyle="solid"}. Style de trait pour la grille secondaire.
            \item \opt{originloc=A}. Localisation de l'origine de la grille.
        \end{itemize}
\end{itemize}

\paragraph{Exemple :} il est possible de travailler dans un repère non orthogonal. Voici un exemple où l'axe $Ox$ est conservé, mais la première bissectrice devient le nouvel axe $Oy$, on modifie pour cela la matrice de transformation du graphe. À partir de cette modification les affixes représentent les coordonnées dans le nouveau repère.
\begin{demo}{Exemple de repère non orthogonal}
\begin{luadraw}{name=axes_non_ortho}
local ld = luadraw
local g = ld.graph:new{window={-5.25,5.25,-4,4},size={10,10}}
local i, pi, Z = ld.cpx.I, math.pi, ld.cpx.Z
local f = function(x) return 2*math.sin(x) end
g:Setmatrix({0,1,1+i}); g:Labelsize("small")
g:Dgrid({-5-4*i,5+4*i},{gridstyle="dashed"})
g:Daxes({0,1,1}, {arrows="-Stealth"})
g:Lineoptions("solid","ForestGreen",12); g:Dcartesian(f,{x={-5,5}})
g:Dcircle(0,3,"Crimson")
g:DlineEq(1,0,3,"Navy") -- droite d'équation x=-3
g:Lineoptions("solid","black",8); g:DtangentC(f,pi/2,1.5,"<->")
g:Dpolyline({pi/2,pi/2+2*i,2*i},"dotted")
g:Ddots(Z(pi/2,2))
g:Dlabeldot("$\\frac{\\pi}2$",pi/2,{pos="SW"})
g:Show()
\end{luadraw}
\end{demo}

\subsubsection{Dgradbox}

La méthode \cmd{g:Dgradbox(\{A, B, xpas, ypas\}, options)} permet le dessin d'une boîte graduée.
\begin{itemize}
    \item Le premier argument est obligatoire, il précise le coin inférieur gauche (c'est le complexe \argu{A}) et le coin supérieur droit (c'est le complexe \argu{B}) de la boîte, ainsi que le pas sur chaque axe.
    \item L'argument \argu{options} est une table précisant les options possibles. Ce sont les mêmes que pour les axes, mises à part certaines valeurs par défaut. À celles-ci s'ajoute l'option suivante : \opt{title=""} qui permet d'ajouter un titre en haut de la boite, attention cependant à laisser suffisamment de place pour cela.
\end{itemize}

\begin{demo}{Utilisation de Dgradbox}
\begin{luadraw}{name=gradbox}
local ld = luadraw
local g = ld.graph:new{window={-5,4,-5.5,5},size={10,10}}
local i, pi = ld.cpx.I, math.pi
local h = function(x) return x^2/2-2 end
local f = function(x) return math.sin(3*x)+h(x) end
g:Dgradbox({-pi-4*i,pi+4*i,pi/3,1},{grid=true,originloc=0, originnum={0,0}, labeltext={"\\pi",""},
labelden={3,1}, title="\\textbf{Title}", legend={"Legende $x$","Legende $y$"}})
g:Saveattr(); g:Viewport(-pi,pi,-4,4) -- on limite la vue (clip)
g:Filloptions("full","blue",0.6); g:Linestyle("noline"); g:Ddomain2(f,h,{x={-pi/2,2*pi/3}})
g:Filloptions("none",nil,1); g:Lineoptions("solid",nil,8); g:Dcartesian(h,{x={-pi,pi}, draw_options="DarkBlue"})
g:Dcartesian(f,{x={-pi,pi},draw_options="Crimson"})
g:Restoreattr()
g:Show()
\end{luadraw}
\end{demo}

\subsection{Dessins d'ensembles (diagrammes de Venn)}

\subsubsection{Dessiner un ensemble}

La fonction \cmd{ld.set(center \fac{, angle, scale})} renvoie un chemin représentant un ensemble (en forme d'œuf), dont le centre est \argu{center} (complexe), l'argument \argu{angle} représente l'inclinaison (en degrés) de l'axe vertical de l'ensemble ($0$ par défaut), et l'argument \argu{scale} est un facteur d'échelle permettant de modifier la taille de l'ensemble ($1$ par défaut). Un tel chemin peut être dessiné avec la méthode \cmd{g:Dpath()}.

\begin{demo}{Dessiner un ensemble}
\begin{luadraw}{name=set}
local ld = luadraw
local g = ld.graph:new{window={-5.25,5.25,-5,5},size={10,10}}
local i = ld.cpx.I
local A, B, C = ld.set(i,0), ld.set(-2-i,25), ld.set(2-i,-25)
g:Fillopacity(0.3)
g:Dpath(A,"fill=orange"); g:Dpath(B,"fill=blue")
g:Dpath(C,"fill=green")
g:Fillopacity(1)
g:Dlabel("$A$",5*i,{pos="N"},"$B$",-4+3*i,{pos="W"},"$C$",4+3*i,{pos="E"})
g:Show()
\end{luadraw}
\end{demo}

\subsubsection{Opérations sur les ensembles}

Notons $C_1$ et $C_2$ deux listes de complexes représentant le contour de deux ensembles (courbes fermées simples, d'un seul tenant). Les opérations possibles sont au nombre de trois :
\begin{itemize}
    \item La fonction \cmd{ld.cap(C1, C2)} renvoie une liste de complexes représentant le contour de l'intersection des ensembles correspondant à $C_1$ et $C_2$.
    \item La fonction \cmd{ld.cup(C1, C2)} renvoie une liste de complexes représentant le contour de la réunion des ensembles correspondant à $C_1$ et $C_2$.
        \item La fonction \cmd{ld.setminus(C1, C2)} renvoie une liste de complexes représentant le contour de la différence des ensembles correspondant à $C_1$ et $C_2$ ($C_1\setminus C_2$).    
\end{itemize}
Le résultat de ces opérations, étant une liste de complexes, peut être dessiné avec la méthode \cmd{g:Dpolyline()}.

\begin{demo}{Opérations sur les ensembles}
\begin{luadraw}{name=cap_and_cup}
local ld = luadraw
local g = ld.graph:new{window={-5.5,5.5,-5,5},size={10,10}}
local i = ld.cpx.I
local A, B, C = ld.set(i,0), ld.set(-2-i,25), ld.set(2-i,-25)
g:Fillopacity(0.3)
g:Dpath(A,"fill=orange"); g:Dpath(B,"fill=blue"); g:Dpath(C,"fill=green")
g:Fillopacity(1)
local C1, C2, C3 = ld.path(A), ld.path(B), ld.path(C) -- conversion: chemin -> liste de complexes
local I = ld.cap(ld.cup(C1,C2),C3)
g:Linecolor("red"); g:Filloptions("full","white")
g:Dpolyline(I,true,"line width=0.8pt,fill opacity=0.8")
g:Dlabel("$A$",5*i,{pos="N"},"$B$",-4+3*i,{pos="W"},"$C$",4+3*i,{pos="E"},
"$(A\\cup B) \\cap C$",-i,{pos="NE",node_options="red,draw"})
g:Show()
\end{luadraw}
\end{demo}

\paragraph{NB} : le résultat n'est pas toujours satisfaisant lorsque les contours deviennent trop complexes, ou lorsque les contours ont des tronçons en commun.


\subsection{Importer une image}

\subsubsection{Placer une image dans le graphique}

Pour importer une image dans le graphique en cours (et éventuellement dessiner par dessus), il y a la méthode : \cmdln{g:Dimage(filename, anchor, options)}
\begin{itemize}
    \item \argu{filename} est le nom complet du fichier image, celle-ci sera affichée par la commande \verb|\includegraphics[]{}| dans un node,
    \item \argu{anchor} est un nombre complexe représentant le point d'ancrage de l'image dans le graphique,
    \item \argu{options} est une table dont les champs définissent les paramètres d'affichage, qui sont (avec leur valeur par défaut):
    \begin{itemize}
        \item \opt{pos="center"} qui indique la position de l'image par rapport au point d'ancrage, sur le même principe que pour les labels. Les valeurs possibles sont: \val{"center"}, \val{"N"}, \val{"NE"}, \val{"E"}, \val{"SE"}, \val{"S"}, \val{"SW"}, \val{"W"}, \val{"NW"}.
        \item \opt{name=""} qui permet éventuellement de donner un nom au node qui va être créé (nom qui pourra être utilisé par TikZ ensuite),
        \item \opt{graphics\_options=""} qui est une chaîne contenant des options pour la commande \verb|\includegraphics|,
        \item \opt{matrix=nil} qui est une table représentant une matrice de transformation affine, celle-ci sera appliquée localement au node créé (pour le node, l'origine est le point d'ancrage). Il est à noter que la matrice 2D de transformation globale du graphique en cours est également prise en compte par cette méthode.
    \end{itemize}
\end{itemize}

\begin{demo}{Importer une image}
\begin{luadraw}{name=Dimage}
local ld = luadraw
local Z, i = ld.cpx.Z, ld.cpx.I
local g = ld.graph:new{margin={0,0,0,0}, size={10,10},bg="green!30"}
local filename = ld.cachedir.."flower.jpg"
local g_options = "width=4.5cm,height=4.5cm"
g:Dimage(filename,Z(-5,5),{pos="SE", graphics_options=g_options})
g:Dimage(filename,Z(5,5),{pos="SE", matrix={0,-1,i}, graphics_options=g_options})
g:Rotate(180)
g:Dimage(filename,Z(-5,5),{pos="SE", graphics_options=g_options})
g:IDmatrix()
g:Dimage(filename,Z(-5,-5),{pos="SE", matrix={0,1,-i}, graphics_options=g_options})
g:Dlabel( "Originale", Z(-5,5),{pos="SE"},
"Symétrie verticale", Z(5,5), {pos="SW"},
"Rotation (globale)", Z(5,-5), {pos="NW"},
"Symétrie horizontale", Z(-5,-5), {pos="NE"} )
g:Show()
\end{luadraw}
\end{demo}

\subsubsection{Mapper une image sur un parallélogramme}

Ceci est possible avec la méthode : \cmd{g:Dmapimage(filename, parallelo, options)}, où :
\begin{itemize}
    \item  \argu{filename} est le nom complet du fichier image, 
    \item \argu{parallelo} est une liste de la forme $\{a,u,v\}$ constituée de trois nombres complexes, cette liste représente le parallélogramme de sommets $\{a,a+u,a+u+v,a+v\}$.
    \item \argu{options} est une table dont les champs définissent les paramètres, qui sont (avec leur valeur par défaut):
    \begin{itemize}
        \item \opt{name=""} qui permet éventuellement de donner un nom au node qui va être créé (nom qui pourra être utilisé par TikZ ensuite).
        \item \opt{clip=false} qui est un booléen indiquant si l'image doit être clippée avec le parallélogramme, mais normalement cela n'est pas nécessaire.
        \item \opt{border\_options=nil}, lorsque cette option n'est pas égale à \nil, ce doit être une chaîne de caractères contenant les options de dessin le contour du parallélogramme, \opt{border\_options} sera alors transmise directement à l'instruction \drawcmd pour dessiner le contour.
    \end{itemize}    
\end{itemize}

\begin{demo}{Mapper une image}
\begin{luadraw}{name=Dmapimage}
local ld = luadraw
local g = ld.graph:new{window={-2,3.5,-0.5,7}, margin={0,0,0,0}, size={10,10},bg="green!30"}
local filename = ld.cachedir.."flower.jpg"
local i = ld.cpx.I
local a, u, v, w = 0, 3, -1.5+i, 4*i
local facets = {{a,u,w}, {a+v,-v,w}, {a+v+w,-v,u} }
local portrait = {a+w+(2*v+u)/3+0.2, (u-v)/4, 1.5*i+0.2}
for _,f in ipairs(facets) do
    g:Dmapimage(filename,f,{border_options="Navy"})
end
g:Dmapimage(ld.cachedir.."russell.jpg", portrait, {border_options="lightgray,line width=1.5mm"})
g:Show()
\end{luadraw}
\end{demo}


\subsection{Les couleurs}

Dans l'environnement \luadrawenv les couleurs sont des chaînes de caractères qui doivent correspondre à des couleurs connues de TikZ. Le package \emph{xcolor} est fortement conseillé pour ne pas être limité aux couleurs de bases.

\subsubsection{Calculs sur les couleurs}

Afin de pouvoir faire des manipulations sur les couleurs, celles-ci ont été définies (dans le module \emph{luadraw\_colors.lua}) sous la forme de tables de trois composantes : rouge, vert, bleu, chaque composante étant un nombre entre $0$ et $1$, et avec leur nom au format \emph{svgnames} du package \emph{xcolor}, par exemple on y trouve (entre autres) les déclarations :
\begin{Luacode}
local ld = luadraw
ld.AliceBlue = {0.9412, 0.9725, 1}
ld.AntiqueWhite = {0.9804, 0.9216, 0.8431}
ld.Aqua = {0.0, 1.0, 1.0}
ld.Aquamarine = {0.498, 1.0, 0.8314}
\end{Luacode}
On pourra se référer à la documentation de \emph{xcolor} pour avoir la liste de ces couleurs.

Pour utiliser celles-ci dans l'environnement \luadrawenv, on peut :
\begin{itemize}
    \item soit les utiliser avec leur nom si on a déclaré dans le préambule : \verb|\usepackage[svgnames]{xcolor}|, par exemple : \code{g:Linecolor("AliceBlue")},
    
    \item soit les utiliser avec la fonction \cmd{ld.rgb()} de \luadrawenv, par exemple : \code{g:Linecolor(ld.rgb(AliceBlue))}. Par contre, avec cette fonction \cmd{ld.rgb()}, pour changer localement de couleur il faut faire comme ceci (exemple) : \par
    \code{g:Dpolyline(L,"color="..ld.rgb(AliceBlue))}, ou \code{g:Dpolyline(L,"fill="..ld.rgb(AliceBlue))}. Car la fonction \cmd{rgb()} ne renvoie pas un nom de couleur, mais une définition de couleur.
\end{itemize}

\paragraph{Fonctions pour la gestion des couleurs :}
\begin{itemize}
    \item La fonction \cmd{ld.rgb(r, g, b)} ou \cmd{ld.rgb(\{r, g, b\})}, renvoie la couleur sous forme d'une chaîne de caractères compréhensible par TikZ dans les options \code{color=\ldots} et \code{fill=\ldots}. Les valeurs de $r$, $g$ et $b$ doivent être entre $0$ et $1$.
    
    \item La fonction \cmd{ld.hsb(h, s, b \fac{, table})} renvoie la couleur sous forme d'une chaîne de caractères compréhensible par TikZ. L'argument \argu{h} (hue) doit être un nombre enter $0$ et $360$, l'argument \argu{s} (saturation) doit être entre $0$ et $1$, et l'argument \argu{b} (brightness) doit être aussi entre $0$ et $1$. L'argument (facultatif) \argu{table} est un booléen (\false par défaut) qui indique si le résultat doit être renvoyé sous forme de table $\{r,g,b\}$ ou non (par défaut c'est sous forme d'une chaîne).
    
    \item La fonction \cmd{ld.mixcolor(color1, proportion1, color2, proportion1, ..., colorN, proportionN)} mélange les couleurs \argu{color1}, \ldots, \argu{colorN} dans les proportions demandées et renvoie la couleur qui en résulte sous forme d'une chaîne de caractères compréhensible par TikZ, suivie de cette même couleur sous forme de table $\{r,g,b\}$. Chacune des couleurs doit être une table de trois composantes $\{r,g,b\}$.
    
    \item La fonction \cmd{ld.mixpalette(pal, percent, color)} renvoie une nouvelle palette après avoir mixer chaque couleur de la palette \argu{pal} avec \argu{color}. L'argument \argu{pal} est une liste de couleurs, chacune d'elles étant une table de trois composantes $\{r,g,b\}$, l'argument \argu{percent} est un nombre entre $0$ et $100$, et l'argument \argu{color} est une couleur sous forme de table $\{r,g,b\}$.
    
    \item La fonction \cmd{ld.palette(colors, pos \fac{, table})} : l'argument \argu{colors} est une liste (table) de couleurs au format $\{r,g,b\}$, l'argument \argu{pos} est un nombre entre $0$ et $1$, la valeur $0$ correspond à la première couleur de la liste et la valeur $1$ à la dernière. La fonction calcule et renvoie la couleur correspondant à la position \argu{pos} dans la liste par interpolation linéaire. L'argument (facultatif) \argu{table} est un booléen (\false par défaut) qui indique si le résultat doit être renvoyé sous forme de table $\{r,g,b\}$ ou non (par défaut c'est sous forme d'une chaîne).
    
    \item La fonction \cmd{ld.getpalette(colors, nb \fac{, table})} : l'argument \argu{colors} est une liste (table) de couleurs au format $\{r,g,b\}$, l'argument \argu{nb} indique le nombre de couleurs souhaité. La fonction renvoie une liste de \argu{nb} couleurs régulièrement réparties dans \argu{colors}. L'argument (facultatif) \argu{table} est un booléen (\false par défaut) qui indique si les couleurs sont renvoyées sous forme de tables $\{r,g,b\}$ ou non (par défaut c'est sous forme de chaînes).  
    
    \item La méthode \cmd{g:Newcolor(name, color)} permet de définir dans l'export TikZ au format rgb une nouvelle couleur dont le nom sera \argu{name} (chaîne), l'argument \argu{color} peut être soit une table de trois composantes : rouge, vert, bleu (entre 0 et 1) définissant cette couleur, soit une chaîne de caractères représentant une couleur. Dans le premier cas la méthode utilise un \verb|\definecolor| et dans le second cas elle utilise un \verb|\colorlet|.
    
\end{itemize}
On peut également utiliser toutes les possibilités habituelles de TikZ pour la gestion des couleurs.

Par défaut, il y a cinq palettes de couleurs\footnote{Une palette est une table de couleurs, celles-ci sont elle-mêmes des tables de nombres entre $0$ et $1$ représentant les composantes rouge, vert, bleu.}.

\begin{demo}{Les cinq palettes par défaut}
\begin{luadraw}{name=palettes}
local ld = luadraw
local Z = ld.cpx.Z
local g = ld.graph:new{window={-5,5,-5,5},bbox=false, border=true}
g:Linewidth(1)
local Dpalette = function(pal,A,h,L,N,name)
    local dl = L/N
    for k = 1, N do
        local color = ld.palette(pal,(k-1)/(N-1))
        g:Drectangle(A,A+h,A+h+dl,"color="..color..",fill="..color)
        A = A+dl
    end
    g:Drectangle(A,A+h,A+h-L); g:Dlabel(name,A+h/2,{pos="E"})
end
local A, h, dh, L, N = Z(-5,4), Z(0,-1), Z(0,-1.1), 5, 100
Dpalette(ld.rainbow,A,h,L,N,"ld.rainbow"); A = A+dh
Dpalette(ld.palAutumn,A,h,L,N,"ld.palAutumn"); A = A+dh
Dpalette(ld.palGistGray,A,h,L,N,"ld.palGistGray"); A = A+dh
Dpalette(ld.palGasFlame,A,h,L,N,"ld.palGasFlame"); A = A+dh
Dpalette(ld.palRainbow,A,h,L,N,"ld.palRainbow")
g:Show()
\end{luadraw}
\end{demo}
%
\section{Constructions géométriques}

Dans cette section sont regroupées les fonctions construisant des figures géométriques sans méthode graphique dédiée correspondante.

\subsection{circumcircle(), incircle()}

\begin{itemize}
    \item La fonction \cmd{ld.circumcircle(a, b, c)} (ou \cmd{ld.circumcircle(\{a, b, c\})}), où \argu{a}, \argu{b} et \argu{c} sont trois points (trois nombres complexes), renvoie le cercle circonscrit au triangle formé par ces trois points, sous la forme d'une séquence: $C,r$, où $C$ est le centre du cercle (nombre complexe), $r$ son rayon.
    \item La fonction \cmd{ld.incircle3d(a, b, c)} (ou \cmd{ld.incircle3d(\{a, b, c\})}), où \argu{a}, \argu{b}, et \argu{c} sont trois points (trois nombres complexes), renvoie le cercle inscrit dans triangle formé par ces trois points, sous la forme d'une séquence: $C,r$, où $C$ est le centre du cercle (nombre complexe), $r$ son rayon.    
\end{itemize}

\subsection{cvx\_hull2d()}

La fonction \cmd{ld.cvx\_hull2d(L)} où \argu{L} est une liste de complexes, calcule et renvoie une liste de complexes représentant l'enveloppe convexe de $L$.

\subsection{delaunay()}

La fonction \cmd{ld.delaunay(L)} où \argu{L} est une liste de nombres complexes \textbf{distincts}, renvoie une liste de triangles (un triangle étant une liste de trois nombres complexes) obtenus par triangulation de Delaunay des points de \argu{L} (le cercle circonscrit de chacun des triangles ne contient aucun des autres points).

\begin{demo}{Triangulation de Delaunay}
\begin{luadraw}{name=delaunay}
local ld = luadraw
local g = ld.graph:new{bbox=false, pictureoptions="scale=2"}
local i = ld.cpx.I; g:Linewidth(6)
local L = {0.285+1.46*i,1.556-0.142*i,2.344+1.313*i,-2.38+1.218*i,1.548-0.624*i,
0.969+1.819*i, -0.086-2.191*i,-0.477+1.834*i,-0.904+1.322*i,-2.892+0.025*i}
local T = ld.delaunay(L) -- liste de triangles
local n = #T
local num = 7 --on choisit un triangle
local colors = ld.getpalette(ld.palGasFlame,n)
for k = 1, n do
    g:Dpolyline(T[k],true,'fill='..colors[k])
    end
g:Ddots(L)
g:Dcircle( {ld.circumcircle(T[num])}, "line width=0.4pt,gray,dashed" )
g:Ddots(T[num],"gray")
g:Show()
\end{luadraw}
\end{demo}

\subsection{voronoi()}

La fonction \cmd{ld.voronoi(L \fac{, window})} où \argu{L} est une liste de nombres complexes \textbf{distincts}, détermine le diagramme de Voronoï des points de la liste \argu{L}. Cette fonction renvoie une liste d'éléments de la forme \code{\{A,polygone\}}  où \code{A} est un point la liste \argu{L}, et \code{polygone} une liste de nombres complexes représentant les sommets de la cellule associée à \code{A}. Il y a ainsi une cellule par point de \argu{L}. La cellule du point \code{A} contient les points du plan qui sont plus proches de \code{A} que des autres points de \argu{L}. Cette fonction utilise la triangulation de Delaunay. L'argument optionnel \argu{window}, qui vaut par défaut \code{\{-5,5,-5,5\}}, est utilisée pour clipper les cellules de Voronoï qui sont non bornées, cette fenêtre est automatiquement agrandie si nécessaire, pour contenir tous les points de \argu{L} ainsi que tous les centes des cercles circonscrits au triangles de Delaunay (attention: cela ne change pas la fenêtre 2D du graphique en cours).

\begin{demo}{Diagramme de Voronoï}
\begin{luadraw}{name=voronoi}
local ld = luadraw
local g = ld.graph:new{ bbox=true, margin={0,0,0,0}, size={10,10}}
local i = ld.cpx.I
local S = {0.285+1.46*i,1.556-0.142*i,2.344+1.313*i,-2.38+1.218*i,1.548-0.624*i,
0.969+1.819*i,-0.086-2.191*i,-0.477+1.834*i,-0.904+1.322*i,-2.892+0.025*i}
local V = ld.voronoi(S)
local colors = ld.getpalette(ld.rainbow,#V)
for k,T in ipairs(V) do
    local A, polygon = table.unpack(T)
    g:Dpolyline(polygon,true,"color=white, line width=1.2pt,fill="..colors[k])
    g:Ddots(A,"mark=x,white,scale=2,line width=1.2pt")-- A est un des points de S
end
g:Dpolyline(ld.delaunay(S),true,"dotted,line width=0.6pt") -- triangles de Delaunay
g:Show()
\end{luadraw}
\end{demo}

\subsection{line2strip()}

La fonction \cmd{ld.line2strip(L, width \fac{, close, ends, mode})} où \argu{L} est une liste de nombres complexes, ou une liste de listes de nombres complexes, renvoie un chemin représentant une "bande" calculée sur \argu{L} et de largeur \argu{width}. L'argument optionnel \argu{close} est un booléen qui indique si \argu{L} doit être refermée (\false par défaut). 

L'argument optionnel \argu{ends} indique la façon dont les deux extrémités de la bande sont dessinées, les valeurs possibles sont : \val{"none"}, \val{"butt"} (valeur par défaut) ou \val{"round"}. Pour rester compatible avec l'ancienne version, cet argument peut également prendre les valeurs \true (équivalent à \val{"butt"}) ou \false (équivalent à \val{"none"}).

L'argument \argu{mode} peut valoir soit \val{"center"} (valeur par défaut), soit \val{"left"}, soit \val{"right"}, dans le premier cas la bande est centrée sur \argu{L}, dans le second cas elle est sur la gauche de \argu{L}, et dans le troisième cas elle est sur la droite de \argu{L}.

\begin{demo}{Exemple avec \emph{line2strip}}
\begin{luadraw}{name=line2strip}
local ld = luadraw
local g = ld.graph:new{bbox=false, bg="lightgray"}
local i = ld.cpx.I; g:Linewidth(8)
local p = {-3+3*i,-3,"l",0,3,3,1,"ca", 3+3*i,"l"}
local P = ld.path(p) -- p est converti en ligne polygonale
g:Setmatrix({-3+3*i,0.5,0.5*i})
local L = ld.line2strip(P,1,true)
g:Dpath(L,"Crimson,fill=Gold"); g:Dpath(p,"gray,dashed")
g:Dlabel("close=true",i,{})
g:Setmatrix({3+3*i,0.5,0.5*i})
L = ld.line2strip(P,1,false)
g:Dpath(L,"Crimson,fill=Gold"); g:Dpath(p,"gray,dashed")
g:Dlabel("close=false",i,{})
g:Setmatrix({-3-i,0.5,0.5*i})
L = ld.line2strip(P,1,false,"round","right")
g:Dpath(L,"Crimson,fill=Gold"); g:Dpath(p,"gray,dashed")
g:Dlabel('mode="right"',1.5*i,{}); g:Dlabel('ends="round"',0.5*i,{})
g:Setmatrix({3-i,0.5,0.5*i})
L = ld.line2strip(P,1,true,false,"left")
g:Dpath(L,"Crimson,fill=Gold"); g:Dpath(p,"gray,dashed")
g:Dlabel('mode="left"',1.5*i,{}); g:Dlabel("close=true",0.5*i,{})
g:Show()
\end{luadraw}
\end{demo}

\subsection{parallel\_polyline()}

La fonction \cmd{ld.parallel\_polyline(L, width \fac{, close})} où \argu{L} est une liste de nombres complexes, ou une liste de listes de nombres complexes, renvoie une ligne polygonale parallèle à \argu{L} et située à une "distance" égale à \argu{width}. L'argument \argu{width} peut être positif ou négatif pour être d'un coté ou de l'autre de \argu{L} (cela dépend du sens de parcours de \argu{L}). L'argument optionnel \argu{close} est un booléen qui indique si \argu{L} doit être refermée (\false par défaut).


\subsection{sss\_triangle()}

La fonction \cmd{ld.sss\_triangle(ab, bc, ca)} où \argu{ab}, \argu{bc} et \argu{ca} sont trois longueurs, calcule et renvoie une liste de trois points (3 complexes) $\{A,B,C\}$ formant les sommets d'un triangle direct dont les longueurs des côtés sont les arguments, c'est-à-dire $AB=ab$, $BC=bc$ et $CA=ca$, lorsque cela est possible. Le sommet $A$ est toujours le complexe $0$ et le sommet $B$ est toujours le complexe $ab$. Ce triangle peut être dessiné avec la méthode \cmd{g:Dpolyline}.

\subsection{sas\_triangle()}

La fonction \cmd{ld.sas\_triangle(ab, alpha, ca)} où \argu{ab} et \argu{ca} sont deux longueurs, \argu{alpha} un angle en degrés, calcule et renvoie une liste de trois points (3 complexes) $\{A,B,C\}$ formant les sommets d'un triangle tel que $AB=ab$, $CA=ca$, et tel que l'angle $(\vec{AB},\vec{AC})$ a pour mesure \argu{alpha}, lorsque cela est possible. Le sommet $A$ est toujours le complexe $0$ et le sommet $B$ est toujours le complexe $ab$. Ce triangle peut être dessiné avec la méthode \cmd{g:Dpolyline}.

\subsection{asa\_triangle()}

La fonction \cmd{ld.asa\_triangle(alpha, ab, beta)} où \argu{ab} est une longueur, \argu{alpha} et \argu{beta} deux angles en degrés, calcule et renvoie une liste de trois points (3 complexes) $\{A,B,C\}$ formant les sommets d'un triangle tel que $AB=ab$, tel que l'angle $(\vec{AB},\vec{AC})$ a pour mesure \argu{alpha}, et tel que l'angle $(\vec{BA},\vec{BC})$ a pour mesure \argu{beta}, lorsque cela est possible. Le sommet $A$ est toujours le complexe $0$ et le sommet $B$ est toujours le complexe $ab$. Ce triangle peut être dessiné avec la méthode \cmd{g:Dpolyline}.


\begin{demo}{sss\_triangle, sas\_triangle et asa\_triangle}
\begin{luadraw}{name=sss_triangles_and_co}
local ld = luadraw
local g = ld.graph:new{window={-5,5,-3,5},size={10,10}}
g:Labelsize("footnotesize"); g:Linewidth(8)
local Zp, i, deg = ld.cpx.Zp, ld.cpx.I, ld.deg
local T1 = ld.shift( ld.sss_triangle(4,5,3), 2*i-2)
local T2 = ld.shift( ld.sas_triangle(4,60,2), -4-2*i)
local T3 = ld.shift( ld.asa_triangle(30,4,50), 0.5-i)
g:Dpolyline({T1,T2,T3}, true)
g:Linewidth(4)
g:Darc(T2[2],T2[1],T2[3],0.5,1,"->")
g:Darc(T3[2],T3[1],T3[3],0.75,1,"->")
g:Darc(T3[1],T3[2],T3[3],0.75,-1,"->")
g:Dlabel(
  "$4$",(T1[1]+T1[2])/2,{pos="N"}, "$5$",(T1[2]+T1[3])/2,{pos="NE"},"$3$",(T1[1]+T1[3])/2,{pos="W"},
  "$4$",(T2[1]+T2[2])/2,{pos="N"},"$60^\\circ$", T2[1]+Zp(0.9,30*deg),{pos="center"}, 
  "$2$",(T2[1]+T2[3])/2,{pos="W"},
  "$4$",(T3[1]+T3[2])/2,{pos="N"}, "$30^\\circ$",T3[1]+Zp(1.15,15*deg),{pos="center"},
  "$50^\\circ$",T3[2]+Zp(1.15,155*deg),{pos="center"},
  "sss\\_triangle(4,5,3)",(T1[1]+T1[2])/2,{pos="S"}, "sas\\_triangle(4,60,2)",(T2[1]+T2[2])/2,{},
  "asa\\_triangle(30,4,50)",(T3[1]+T3[2])/2,{})
for _,T in ipairs({T1,T2,T3}) do
    g:Dlabel("$A$",T[1],{pos="SW"}, "$B$",T[2],{pos="SE"},"$C$",T[3],{pos="N"})
end
g:Show()
\end{luadraw}
\end{demo}
%
\section{Calculs sur les listes}

\subsection{concat}

La fonction \cmd{ld.concat(table1, table2, \ldots)} concatène toutes les tables passées en argument, et renvoie la table qui en résulte.

\begin{itemize}
 \item Chaque argument peut être un réel, un complexe ou une table.
\item Exemple : l'instruction \code{ld.concat(1,2,3,{4,5,6},7)} renvoie la table \val{\{1,2,3,4,5,6,7\}}.
\end{itemize}

\subsection{cut}

La fonction \cmd{ld.cut(L, A \fac{, before})} permet de couper \argu{L} au point \argu{A} qui est supposé être situé sur la ligne \argu{L} (qui est soit une liste de complexes, soit une ligne polygonale c'est-à-dire une liste de listes de complexes). Si l'argument \argu{before} vaut \false (valeur par défaut), alors la fonction renvoie la partie située avant \argu{A}, suivie de la partie située après \argu{A}, sinon c'est l'inverse.

\subsection{cutpolyline}

La fonction \cmd{ld.cutpolyline(L, D \fac{, close})} permet de couper la ligne polygonale \argu{L} avec la droite \argu{D}. L'argument \argu{L} doit être une liste de complexes ou une liste de listes de complexes, l'argument \argu{D} est une liste de la forme \code{\{A,u\}} où $A$ est un complexe (point de la droite) et $u$ un complexe non nul (vecteur directeur de la droite). L'argument \argu{close} indique si la ligne \argu{L} doit être refermée (\false par défaut). La fonction renvoie trois choses:
\begin{itemize}
    \item La partie de \argu{L} qui est dans le demi-plan défini par la droite à "gauche" de $u$ (c'est à dire contenant le point $A+iu$) (c'est une ligne polygonale),
    \item suivi de la partie de \argu{L} qui est dans l'autre demi-plan (ligne polygonale),
    \item suivi de la liste des points d'intersection entre \argu{L} et la droite \argu{D}.
\end{itemize}

\begin{demo}{Illustrer un exercice de programmation linéaire}
\begin{luadraw}{name=cutpolyline}
local ld = luadraw
local g = ld.graph:new{window={-5,5,-5,5}, size={10,10},margin={0,0,0,0}}
g:Linewidth(6)
local i = ld.cpx.I
local P = g:Box2d() -- polygone représentant la fenêtre 2D
local D1, D2, D3 = {0,1+i}, {2.5,-i}, {-3*i,-1-i/4} -- trois droites
local P1 = ld.cutpolyline(P,D1,true)
local P2 = ld.cutpolyline(P,D2,true)
local P3 = ld.cutpolyline(P,D3,true)
g:Daxes({0,1,1},{grid=true,gridcolor="LightGray",arrows="->",legend={"$x$","$y$"}})
g:Filloptions("horizontal","blue"); g:Dpolyline(P1,true,"draw=none")
g:Filloptions("fdiag","red"); g:Dpolyline(P2,true,"draw=none")
g:Filloptions("bdiag","green"); g:Dpolyline(P3,true,"draw=none")
g:Filloptions("none","black",1)
g:Linewidth(8)
g:Dline(D1,"blue"); g:Dline(D2,"red"); g:Dline(D3,"green")
g:Dlabel(
  "$x-y\\leqslant 0$",-3-3*i,{pos="N",dir={1+i,-1+i},dist=0.1,node_options="fill=white,fill opacity=0.8"},
  "$x-2.5\\geqslant0$", 2.5+i,{dir={-i,1}},
  "$-\\frac{x}{4}+y+3\\leqslant0$", -3-15/4*i,{pos="S",dir={1+i/4,i-1/4}})
g:Show()
\end{luadraw}
\end{demo}

\subsection{cutpolyline2}

La fonction \cmd{ld.cutpolyline2(P, f, sg, \fac{, close})} permet de couper le polygone \argu{P} avec la courbe représentative de la fonction \argu{f}. L'argument \argu{P} doit être une liste de nombres complexes, l'argument \argu{f} est une fonction \argu{f}$\colon x \mapsto f(x)\in \mathbf R$. L'argument \argu{sg} est soit \val{">"}, soit \val{"<"}, dans le premier cas, la fonction renvoie la partie du polygone \argu{P} dont les points vérifient $y>f(x)$, dans le second cas ce sont les points vérifiant $y<f(x)$. L'argument facultatif \argu{close} indique si \argu{P} doit être refermé. 

\begin{demo}{cutpolyline2}
\begin{luadraw}{name=cupolyline2}
local ld = luadraw
local g = ld.graph:new{window = {-4, 4, -4, 4}, size ={10,10}}
g:Daxes({0,1,1}, {grid=true, arrows="-Stealth", legend={"$x$","$y$"}})
local f = function(x) return 1.5*math.sin(x) end
local P = ld.polyreg(0,3,5) -- pentagone
local L = ld.cutpolyline2(P,f,'<',true) -- partie du polygone située sous la courbe de f
g:Dpolyline(L,true,"draw=none,fill=pink,fill opacity=0.6")
g:Linewidth(12)
g:Dcartesian(f, {draw_options="blue"})
g:Dpolyline(P,true,"red")
g:Show()
\end{luadraw}
\end{demo}

\subsection{getbounds}

\begin{itemize}
    \item La fonction \cmd{ld.getbounds(L)} renvoie les bornes \val{xmin,xmax,ymin,ymax} de la ligne polygonale \argu{L}.
    \item Exemple : \code{local xmin,xmax,ymin,ymax=ld.getbounds(L)} (où \argu{L} désigne une ligne polygonale).
\end{itemize}

\subsection{getdot}
La fonction \cmd{ld.getdot(x, L)} renvoie le point d'abscisse \argu{x} (réel entre $0$ et $1$) le long de la composante connexe \argu{L} (liste de complexes). L'abscisse $0$ correspond au premier point et l'abscisse $1$ au dernier, plus généralement, \argu{x} correspond à un pourcentage de la longueur de \argu{L}.

\subsection{insert}
La fonction \cmd{ld.insert(table1, table2 \fac{, pos})} insère les éléments de \argu{table2} dans \argu{table1} à la position \argu{pos}.

\begin{itemize}
    \item L'argument \argu{table2} peut être un réel, un complexe ou une table.
    \item L'argument \argu{table1} doit être une variable qui désigne une table, celle-ci sera modifiée par la fonction.
    \item Si l'argument \argu{pos} vaut \nil, l'insertion se fait à la fin de \argu{table1}.
    \item Exemple : si une variable $L$ vaut \val{\{1,2,6\}}, alors après l'instruction \code{ld.insert(L,{3,4,5},3)}, la variable $L$ sera égale à \val{\{1,2,3,4,5,6\}}.
\end{itemize}

\subsection{interCC}
La fonction \cmd{ld.interCC(C1, C2)} renvoie l'intersection cercle \argu{C1} avec le cercle \argu{C2}, où \code{C1=\{O1,r1\}} (cercle de centre $O_1$ et de rayon $r_1$), et \code{C2=\{O2,r2\}} (cercle de centre $O_2$ et de rayon $r_2$). La fonction renvoie une liste contenant $1$ ou $2$ points, ou bien le cercle en entier si l'intersection n'est pas vide, sinon elle renvoie \nil.

\begin{demo}{Tangentes à un cercle \{O,2\} et à une ellipse \{O,3,2\} issues d'un point}
\begin{luadraw}{name=interCC}
local ld = luadraw
local cpx = ld.cpx
local g = ld.graph:new{window={-10,10,-5,5}, margin={0,0,0,0},size={16,8}}
local i = cpx.I
-- pour le cercle {O,2}
g:Saveattr(); g:Viewport(-10,0,-5,5); g:Coordsystem(-4,6,-5,5)
local O = -1
local C1, I = {O, 2}, 4-i
local C2 = {(O+I)/2,cpx.abs(I-O)/2}
local rep = ld.interCC(C1,C2) -- points de tangence
g:Dcircle(C1,"blue"); g:Dcircle(C2,"dashed")
g:Dhline(I,rep[1],"red"); g:Dhline(I,rep[2],"red") --demi-tangentes
g:Ddots(rep); g:Ddots({O,I}); g:Dlabel("$I$",I,{pos="SE"},"$O$",O,{pos="W"},
  "tangentes au cercle issues de $I$",1-5*i,{pos="N"})
g:Restoreattr()
-- pour l'ellipse (E) : {O,3,2}
g:Saveattr(); g:Viewport(0,10,-5,5); g:Coordsystem(-4,6,-5,5)
local mat = {0,1.5,i} -- cette matrice transforme un cercle {01,2} en l'ellipse (E)
local inv_mat = ld.invmatrix(mat) -- matrice inverse
local O1, I1 = table.unpack( ld.mtransform({O,I},inv_mat) ) -- antécédents de O et de I
C1 = {O1, 2}
C2 = {(O1+I1)/2,cpx.abs(I1-O1)/2}
rep = ld.interCC(C1,C2) -- points de tangence (tangentes issues de I1)
g:Composematrix(mat) -- on applique la matrice pour retrouver l'ellipse, la tangence est conservée
g:Dcircle(C1,"blue"); g:Dcircle(C2,"dashed")
g:Dhline(I1,rep[1],'red'); g:Dhline(I1,rep[2],"red")
g:Ddots(rep); g:Ddots({O1,I1}); g:Dlabel("$I$",I1,{pos="SE"},"$O$",O1,{pos="W"},
  "tangentes à l'ellipse issues de $I$",1-5*i,{pos="N"})
g:Restoreattr()
g:Show()
\end{luadraw}
\end{demo}

\subsection{interD}
La fonction \cmd{ld.interD(d1, d2)} renvoie le point d'intersection des droites \argu{d1} et \argu{d2}, une droite est une liste de deux complexes : un point de la droite et un vecteur directeur.

\subsection{interDC}
La fonction \cmd{ld.interDC(d, C)} renvoie l'intersection de la droite \argu{d} avec le cercle \argu{C}, où \code{d=\{A,u\}} (droite passant par $A$ et dirigée par $u$), et \code{C=\{O,r\}} (cercle de centre $O$ et de rayon $r$). La fonction renvoie une liste contenant $1$ ou $2$ points  si l'intersection n'est pas vide, elle renvoie \nil sinon.


\subsection{interDL}
La fonction \cmd{ld.interDL(d, L)} renvoie la liste des points d'intersection entre la droite \argu{d} et la ligne polygonale \argu{L}.

\subsection{interL}
La fonction \cmd{ld.interL(L1, L2)} renvoie la liste des points d'intersection des lignes polygonales définies par \emph{L1} et \emph{L2}, ces deux arguments sont deux listes de complexes ou deux listes de listes de complexes).

\subsection{interP}
La fonction \cmd{ld.interP(P1, P2)} renvoie la liste des points d'intersection des chemins définis par \argu{P1} et \argu{P2}, ces deux arguments sont deux listes de complexes et d'instructions (voir \emph{Dpath} page \pageref{Dpath}).


\subsection{linspace}
La fonction \cmd{ld.linspace(a, b \fac{, nbdots})} renvoie une liste de \argu{nbdots} nombres équirépartis de \argu{a} jusqu'à \argu{b}. Par défaut \argu{nbdots} vaut $50$.

\subsection{map}
La fonction \cmd{ld.map(f, list)} applique la fonction \argu{f} à chaque élément de la \argu{list} et renvoie la table des résultats. Lorsqu'un résultat vaut \nil, c'est le complexe \val{cpx.Jump} qui est inséré dans la liste.

\subsection{merge}
La fonction \cmd{ld.merge(L \fac{, epsilon})} recolle si c'est possible, les composantes connexes de \argu{L} qui doit être une liste de listes de complexes, la fonction renvoie une nouvelle ligne polygonale. Les comparaisons se font à \argu{epsilon} près, par défaut on a \argu{epsilon}$=10^{-10}$.

\subsection{polyline2path}
La fonction \cmd{ld.polyline2path(L)} où \argu{L} est une liste de nombres complexes ou une liste de listes de nombres complexes, renvoie \argu{L} sous la forme d'un chemin (que l'on peut dessiner avec la méthode \cmd{g:Dpath()}).

\subsection{range}
La fonction \cmd{ld.range(a, b \fac{, step})} renvoie la liste des nombres de \argu{a} jusqu'à \argu{b} avec un pas égal à \argu{step}, celui-ci vaut $1$ par défaut.

\subsection{read\_csv\_file}
La fonction \cmd{ld.read\_csv\_file(file, options)} permet la lecture d'un fichier \emph{csv}. L'argument \argu{file} est une chaîne de caractères représentant le fichier avec extension: \code{"<name>.csv"}. L'argument \argu{options} est une table dont les champs représentent les options, celles-ci sont (avec leur valeur par défaut):
\begin{itemize}
    \item \opt{header=true}, avec la valeur \true cela signifie que la première ligne du fichier contient les noms des colonnes.
    
    \item \opt{dic=false}, avec la valeur \true cela signifie que le résultat renvoyé sera une liste de dictionnaires (un par ligne), les clés de ces dictionnaires étant les noms des colonnes. Avec la valeur \false (valeur par défaut) le résultat est une liste de listes de valeurs (une liste de valeurs par ligne). Lorsque l'option \opt{header} a la valeur \false, l'option \opt{dic} reste automatiquement à \false.
    
    \item \opt{sep=","}, chaîne de caractères indiquant le séparateur des valeurs sur chaque ligne.
    
    \item \opt{num=true}, booléen indiquant si les valeurs doivent automatiquement être converties en valeurs numériques (lorsque cette conversion échoue, la valeur reste une chaîne de caractères).
    
    \item \opt{comment="\%\%"}, chaîne de caractères annonçant un commentaire (une ligne commençant par ces caractères est ignorée). \textbf{Attention} : la recherche utilise la méthode \emph{match} de Lua, méthode pour laquelle le caractère \% est un caractère spécial, il doit être doublé ("\%\%") pour être considéré comme un caractère.
\end{itemize}
Le résultat renvoyé par cette fonction est une séquence constituée dans cet ordre par:
    \begin{enumerate}
        \item Une liste de listes de résultats (une liste de résultats par ligne).
        \item La liste des valeurs de la première ligne lorsque l'option \opt{header} a la valeur \true.
    \end{enumerate}


\subsection{Fonctions de clipping}

\begin{itemize}
    \item La fonction \cmd{ld.clipseg(A, B, xmin, xmax, ymin, ymax)} clippe le segment $[A;B]$ (nombres complexes) avec la fenêtre $[x_{\mathrm{min}};x_{\mathrm{max}}]\times [y_{\mathrm{min}};y_{\mathrm{max}}]$ et renvoie le résultat.
    
    \item La fonction \cmd{ld.clipline(d, xmin, xmax, ymin, ymax)} clippe la droite \argu{d} avec la fenêtre $[x_{\mathrm{min}};x_{\mathrm{max}}]\times [y_{\mathrm{min}};y_{\mathrm{max}}]$ et renvoie le résultat. La droite \argu{d} est une liste de deux complexes : un point et un vecteur directeur.
    
    \item La fonction \cmd{ld.clippolyline(L, xmin, xmax, ymin, ymax \fac{, close})} clippe la ligne polygonale \argu{L} avec la fenêtre $[x_{\mathrm{min}};x_{\mathrm{max}}]\times [y_{\mathrm{min}};y_{\mathrm{max}}]$ et renvoie le résultat. L'argument \argu{L} est une liste de complexes ou une liste de listes de complexes. L'argument facultatif \argu{close} (\false par défaut) indique si la ligne polygonale doit être refermée.
    
    \item La fonction \cmd{ld.clipdots(L, xmin, xmax, ymin, ymax)} clippe la liste de points (nombres complexes) \argu{L} avec la fenêtre $[x_{\mathrm{min}};x_{\mathrm{max}}]\times [y_{\mathrm{min}};y_{\mathrm{max}}]$ et renvoie le résultat (les points extérieurs sont simplement exclus). L'argument \argu{L} est une liste de complexes ou une liste de listes de complexes.
\end{itemize}

\subsection{Ajout de fonctions mathématiques}

Outre les fonctions associées aux méthodes graphiques qui font des calculs et renvoient une ligne polygonale (comme \cmd{ld.cartesian}, \cmd{ld.periodic}, \cmd{ld.implicit}, \cmd{ld.odesolve}, etc), le paquet \luadrawenv ajoute quelques fonctions mathématiques qui ne sont pas proposées nativement dans le module \emph{math}.

\subsubsection{Évaluation protégée : evalf}
La fonction \cmd{ld.evalf(f, \ldots)} permet d'évaluer \code{f(\ldots)} et de renvoyer le résultat s'il n'y a pas d'erreur d'exécution par Lua, dans le cas contraire, la fonction renvoie \nil. Exemple, l'exécution de :
\begin{Luacode}
local Z = cpx.Z
local f = function(a,b)
    return 2*Z(a,1/b)
end
print(f(1,0))
\end{Luacode}
provoque l'erreur d'exécution \verb|attempt to perform arithmetic on a nil value| (dans la console), car ici \code{Z(1,1/0)} renvoie \nil, et Lua n'accepte pas un argument égal à \nil dans un calcul. Par contre, l'exécution de :
\begin{Luacode}
local Z = cpx.Z
local f = function(a,b)
    return 2*Z(a,1/b)
end
print(ld.evalf(f,1,0))
\end{Luacode}
ne provoque pas d'erreur de la part de Lua, et il n'y a pas d'affichage non plus dans la console puisque la valeur à afficher est \nil.

\subsubsection{int}
La fonction \cmd{ld.int(f, a, b)} renvoie une valeur approchée de l'intégrale de la fonction \argu{f} sur l'intervalle $[a;b]$ (deux nombres complexes). La fonction \argu{f} est à variable réelle et à valeurs réelles ou complexes. La méthode utilisée est la méthode de Simpson accélérée deux fois avec la méthode de Romberg.

\paragraph{Exemple :}
\begin{TeXcode}
$\int_0^1 e^{t^2}\mathrm d t \approx \directlua{tex.sprint(int(function(t) return math.exp(t^2) end, 0, 1))}$
\end{TeXcode}
\paragraph{Résultat :} $\int_0^1 e^{t^2}\mathrm d t \approx \directlua{tex.sprint(luadraw.int(function(t) return math.exp(t^2) end, 0, 1))}$.

\subsubsection{gcd}
La fonction \cmd{ld.gcd(a, b)} renvoie le plus grand diviseur commun entre \argu{a} et \argu{b} deux entiers.

\subsubsection{lcm}
La fonction \cmd{ld.lcm(a, b)} renvoie le plus petit diviseur commun strictement positif entre \argu{a} et \argu{b} deux entiers.

\subsubsection{nth\_root}
La fonction \cmd{ld.nth\_root(n, x)} renvoie la racine énième du réel\argu{x}, l'argument \argu{n} doit être un entier.

\subsubsection{solve}
La fonction \cmd{ld.solve(f, a, b \fac{, n, df})} fait une résolution numérique de l'équation $f(x)=0$ dans l'intervalle $[a;b]$, l'argument \argu{f} est la fonction, \argu{a}, \argu{b} sont deux nombres complexes. L'intervalle est subdivisé en $n$ morceaux ($25$ par défaut). L'argument facultatif \argu{df} et une fonction qui représente la dérivée de \argu{f}, lorsque cet argument est absent, une approximation de la dérivée est utilisée. La méthode mise en œuvre est une variante de la méthode de Newton. La fonction renvoie une liste de nombres ou bien \nil.

\paragraph{Exemple 1 :}
\begin{TeXcode}
\begin{luacode}
local ld = luadraw
resol = function(f,a,b)
    local y = ld.solve(f,a,b)
    if y == nil then tex.sprint("\\emptyset")
    else
        local str = y[1]
        for k = 2, #y do
            str = str..", ".. y[k]
        end
        tex.sprint(str)
    end
end
\end{luacode}
\def\solve#1#2#3{\directlua{resol(#1,#2,#3)}}%
\begin{luacode}
f1 = function(x) return math.cos(x)-x end
f2 = function(x) return x^3-2*x^2+1/2 end
\end{luacode}
La résolution de l'équation $\cos(x)=x$ dans $[0;\frac{\pi}2]$ donne $\solve{f1}{0}{math.pi/2}$.\par
La résolution de l'équation $\cos(x)=x$ dans $[\frac{\pi}2;\pi]$ donne $\solve{f1}{math.pi/2}{math.pi}$.\par
La résolution de l'équation $x^3-2x^2+\frac12=0$ dans $[-1;2]$ donne : $\{\solve{f2}{-1}{2}\}$.
\end{TeXcode}
\paragraph{Résultat :}\ \par

\begin{luacode}
local ld = luadraw
resol = function(f,a,b)
    local y = ld.solve(f,a,b)
    if y == nil then tex.sprint("\\emptyset")
    else
        local str = y[1]
        for k = 2, #y do
            str = str..", ".. y[k]
        end
        tex.sprint(str)
    end
end
\end{luacode}
\def\solve#1#2#3{\directlua{resol(#1,#2,#3)}}%
\begin{luacode}
f1 = function(x) return math.cos(x)-x end
f2 = function(x) return x^3-2*x^2+1/2 end
\end{luacode}

La résolution de l'équation $\cos(x)=x$ dans $[0;\frac{\pi}2]$ donne $\solve{f1}{0}{math.pi/2}$.\par
La résolution de l'équation $\cos(x)=x$ dans $[\frac{\pi}2;\pi]$ donne $\solve{f1}{math.pi/2}{math.pi}$.\par
La résolution de l'équation $x^3-2x^2+\frac 12=0$ dans $[-1;2]$ donne : $\{\solve{f2}{-1}{2}\}$.

\paragraph{Exemple 2 :} on souhaite tracer la courbe de la fonction $f$ définie par la condition :
\[\forall x\in \mathbf R,\ \int_x^{f(x)} \exp(t^2)\mathrm d t = 1.\]
On a deux méthodes possibles :
\begin{enumerate}
    \item On considère la fonction $G\colon (x,y) \mapsto \int_x^y \exp(t^2)\mathrm d t-1$, et on dessine la courbe implicite d'équation $G(x,y)=0$.
    \item On détermine un réel $y_0$ tel que $\int_0^{y_0}\exp(t^2)\mathrm d t = 1$ et on dessine la solution de l'équation différentielle $y'=e^{x^2-y^2}$ vérifiant la  condition initiale $y(0)=y_0$.
\end{enumerate}
Dessinons les deux :
\begin{demo}{Fonction $f$ définie par $\int_x^{f(x)} \exp(t^2)\mathrm d t = 1$.}
\begin{luadraw}{name=int_solve}
local ld = luadraw
local Z = ld.cpx.Z
local g = ld.graph:new{window={-3,3,-3,3},size={10,10}}
local h = function(t) return math.exp(t^2) end
local G = function(x,y) return ld.int(h,x,y)-1 end
local H = function(y) return G(0,y) end
local F = function(x,y) return math.exp(x^2-y^2) end
local y0 = ld.solve(H,0,1)[1] -- solution de H(x)=0
g:Daxes({0,1,1}, {arrows="->"})
g:Dimplicit(G, {draw_options="line width=4.8pt,Pink"})
g:Dodesolve(F,0,y0,{draw_options="line width=0.8pt"})
g:Lineoptions("dashed","gray",4); g:DlineEq(1,-1,0); g:DlineEq(1,1,0) -- bissectrices
g:Dlabel("${\\mathcal C}_f$",Z(2.15,2),{pos="S"})
g:Show()
\end{luadraw}
\end{demo}

On voit que les deux courbes se superposent bien, cependant la première méthode (courbe implicite) est beaucoup plus gourmande en calculs, la méthode 2 est donc préférable.
%
\section{Transformations}
Dans ce qui suit :
\begin{itemize}
    \item l'argument \argu{L} est soit un complexe, soit une liste de complexes soit une liste de listes de complexes,
    \item la droite \argu{d} est une liste de deux complexes : un point de la droite et un vecteur directeur.
  \end{itemize}
  
\subsection{affin}
La fonction \cmd{ld.affin(L, d, v, k)} renvoie l'image de \argu{L} par l'affinité de base la droite \argu{d}, parallèlement au vecteur \argu{v} et de rapport \argu{k}.

\subsection{ftransform}
La fonction \cmd{ld.ftransform(L, f)} renvoie l'image de \argu{L} par la fonction \argu{f} qui doit être une fonction de la variable complexe. Si un des éléments de \argu{L} est le complexe \val{cpx.Jump} alors celui-ci est renvoyé tel quel dans le résultat.

\subsection{hom}
La fonction \cmd{ld.hom(L, factor \fac{, center})} renvoie l'image de \argu{L} par l'homothétie de centre \argu{center} et de rapport \argu{factor}. Par défaut, l'argument \argu{center} vaut $0$.

\subsection{inv}
La fonction \cmd{ld.inv(L, radius \fac{, center})} renvoie l'image de \argu{L} par l'inversion par rapport au cercle de centre \argu{center} et de rayon \argu{radius}. Par défaut, l'argument \argu{center} vaut $0$.

\subsection{proj}
La fonction \cmd{ld.proj(L, d)} renvoie l'image de \argu{L} par la projection orthogonale sur la droite \argu{d}.

\subsection{projO}
La fonction \cmd{ld.projO(L, d, v)} renvoie l'image de \argu{L} par la projection sur la droite \argu{d} parallèlement au vecteur \argu{v}.

\subsection{rotate}
La fonction \cmd{ld.rotate(L, angle \fac{, center})} renvoie l'image de \argu{L} par la rotation de centre \argu{center} et d'angle \argu{angle} (en degrés). Par défaut, l'argument \argu{center} vaut $0$.  

\subsection{shift}
La fonction \cmd{ld.shift(L, u)} renvoie l'image de \argu{L} par la translation de vecteur \argu{u}.

\subsection{simil}
La fonction \cmd{ld.simil(L, factor, angle \fac{, center})} renvoie l'image de \argu{L} par la similitude de centre \argu{center}, de rapport \emph{factor} et d'angle \emph{angle} (en degrés). Par défaut, l'argument \emph{center} vaut arguargu$0$.

\subsection{sym}
La fonction \cmd{ld.sym(L, d)} renvoie l'image de \argu{L} par la symétrie orthogonale d'axe la droite \argu{d}.

\subsection{symG}
La fonction \cmd{ld.symG(L, d, v)} renvoie l'image de \argu{L} par la symétrie par rapport à la droite \argu{d} suivie de la translation de vecteur \argu{v} (symétrie glissée).

\subsection{symO}
La fonction \cmd{ld.symO(L, d)} renvoie l'image de \argu{L} par la symétrie par rapport à la droite \argu{d} et parallèlement au vecteur \argu{v} (symétrie oblique).

\begin{demo}{Utilisation de transformations}
\begin{luadraw}{name=Sierpinski}
local ld = luadraw
local g = ld.graph:new{window={-5,5,-5,5},size={10,10}}
local i, hom = ld.cpx.I, ld.hom
local rand = math.random
local A, B, C = 5*i, -5-5*i, 5-5*i -- triangle initial
local T, niv = {{A,B,C}}, 5
for k = 1, niv do
    T = ld.concat( hom(T,0.5,A), hom(T,0.5,B), hom(T,0.5,C) )
end
for _,cp in ipairs(T) do
    g:Filloptions("full", ld.rgb(rand(),rand(),rand()))
    g:Dpolyline(cp,true)
end
g:Show()
\end{luadraw}
\end{demo}
%
\section{Calcul matriciel}

Si $f$ est une application affine du plan complexe, on appellera matrice de $f$ la liste (table) :
\codeln{\{f(0),Lf(1),Lf(i)\}}
où $Lf$ désigne la partie linéaire de $f$ (on a $Lf(1) = f(1)-f(0)$ et $Lf(i)=f(i)-f(0)$). La matrice identité est la variable globale notée \varglob{ID} dans le paquet \luadrawenv, elle correspond simplement à la liste $\{0,1,i\}$.

\subsection{Calculs sur les matrices}

\subsubsection{applymatrix et applyLmatrix}
\begin{itemize}
    \item La fonction \cmd{ld.applymatrix(z, M)} applique la matrice \argu{M} au complexe \argu{z} et renvoie le résultat (ce qui revient à calculer $f(z)$ si \argu{M} est la matrice de $f$). Lorsque \argu{z} est le complexe \val{cpx.Jump} alors le résultat est \val{cpx.Jump}. Lorsque \argu{z} est une chaîne de caractères alors la fonction renvoie \argu{z}.
    
    \item La fonction \textbf{ld.applyLmatrix(z, M)} applique la partie linéaire la matrice \argu{M} au complexe \argu{z} et renvoie le résultat (ce qui revient à calculer $Lf(z)$ si \argu{M} est la matrice de $f$). Lorsque \argu{z} est le complexe \val{cpx.Jump} alors le résultat est \val{cpx.Jump}.
\end{itemize}

\subsubsection{composematrix}
La fonction \cmd{ld.composematrix(M1, M2)} effectue le produit matriciel \argu{M1}$\times$\argu{M2} et renvoie le résultat.

\subsubsection{invmatrix}
La fonction \cmd{ld.invmatrix(M)} calcule et renvoie l'inverse de la matrice \argu{M} lorsque cela est possible.

\subsubsection{matrixof}
\begin{itemize}
    \item La fonction \cmd{ld.matrixof(f)} calcule et renvoie la matrice de \argu{f} (qui doit être une application affine du plan complexe).
    \item Exemple : \code{ld.matrixof(function(z) return ld.proj(z,{0,Z(1,-1)}) end)} renvoie \par
     \val{\{0,Z(0.5,-0.5),Z(-0.5,0.5)\}} (matrice de la projection orthogonale sur la deuxième bissectrice).
\end{itemize}

\subsubsection{mtransform et mLtransform}
\begin{itemize}
    \item La fonction \cmd{ld.mtransform(L, M)} applique la matrice \argu{M} à la liste \argu{L} et renvoie le résultat. \argu{L} doit être une liste de complexes ou une liste de listes de complexes, si l'un d'eux est le complexe \val{cpx.Jump} ou une chaîne de caractères alors il est renvoyé tel quel.
    
    \item La fonction \cmd{ld.mLtransform(L, M)} applique la partie linéaire la matrice \argu{M} à la liste \argu{L} et renvoie le résultat. \argu{L} doit être une liste de complexes ou une liste de listes de complexes, si l'un d'eux est le complexe \val{cpx.Jump} alors il est renvoyé tel quel.
\end{itemize}

\subsection{Matrice associée au graphe}

Lorsque l'on crée un graphe dans l'environnement \luadrawenv, par exemple :
\begin{Luacode}
local ld = luadraw
local g = ld.graph:new{window={-5,5,-5,5},size={10,10}}
\end{Luacode}
l'objet \emph{g} créé possède une matrice de transformation qui est initialement l'identité. Toutes les méthodes graphiques utilisées appliquent automatiquement la matrice de transformation du graphe. Cette matrice est désignée par \code{g.matrix}, mais pour manipuler celle-ci, on dispose des méthodes qui suivent.

\subsubsection{g:Composematrix()}
La méthode \textbf{g:Composematrix(M)} multiplie la matrice du graphe \emph{g} par la matrice \argu{M} (avec \argu{M} à droite) et le résultat est affecté à la matrice du graphe. L'argument \argu{M} doit donc être une matrice.

\subsubsection{g:Det2d()}
La méthode \cmd{g:Det2d()} envoie $1$ lorsque la matrice de transformation a un déterminant positif, et $-1$ dans le cas contraire. Cette information est utile lorsqu'on a besoin de savoir si l'orientation du plan a été changée ou non.

\subsubsection{g:IDmatrix()}
La méthode \cmd{g:IDmatrix()} réaffecte l'identité à la matrice du graphe \emph g.

\subsubsection{g:Mtransform()}
La méthode \cmd{g:Mtransform(L)} applique la matrice du graphe \emph g à \argu{L} et renvoie le résultat, l'argument \argu{L} doit être une liste de complexes, ou une liste de listes de complexes.

\subsubsection{g:MLtransform()}
La méthode \cmd{g:MLtransform(L)} applique la partie linéaire de la matrice du graphe \emph g à \argu{L} et renvoie le résultat, l'argument \argu{L} doit être une liste de complexes, ou une liste de listes de complexes.

\begin{demo}{Utilisation de la matrice du graphe}
\begin{luadraw}{name=Pythagore}
local ld = luadraw
local g = ld.graph:new{window={-15,15,0,22},size={10,10}}
local a, b, c = 3, 4, 5 -- un triplet de Pythagore
local i, arccos, exp = ld.cpx.I, math.acos, ld.cpx.exp
local f1 = function(z)
    return (z-c)*a/c*exp(-i*arccos(a/c))+c+i*c end
local M1 = ld.matrixof(f1)
local f2 = function(z)
    return z*b/c*exp(i*arccos(b/c))+i*c end
local M2 = ld.matrixof(f2)
local arbre
arbre = function(n)
    local color = ld.mixcolor(ld.ForestGreen,1,ld.Brown,n)
    g:Linecolor(color); g:Dsquare(0,c,1,"fill="..color)
    if n > 0 then
        g:Savematrix(); g:Composematrix(M1); arbre(n-1)
        g:Restorematrix(); g:Savematrix(); g:Composematrix(M2)
        arbre(n-1); g:Restorematrix()
    end
end
arbre(8)
g:Show()
\end{luadraw}
\end{demo}


\subsubsection{g:Rotate()}
La méthode \cmd{g:Rotate(angle \fac{, center})} modifie la matrice de transformation du graphe \emph g en la composant avec la matrice de la rotation d'angle \argu{angle} (en degrés) et de centre \argu{center}. L'argument \argu{center} est un complexe qui vaut $0$ par défaut.

\subsubsection{g:Scale()}
La méthode \cmd{g:Scale(factor \fac{, center})} modifie la matrice de transformation du graphe \emph g en la composant avec la matrice de l'homothétie de rapport \argu{factor} et de centre \argu{center}. L'argument \argu{center} est un complexe qui vaut $0$ par défaut.

\subsubsection{g:Savematrix() et g:Restorematrix()}
\begin{itemize}
    \item La méthode \cmd{g:Savematrix()} permet de sauvegarder dans une pile la matrice de transformation du graphe \emph g.
    \item La méthode \cmd{g:Restorematrix()} permet de restaurer la matrice de transformation du graphe \emph g à sa dernière valeur sauvegardée.
\end{itemize}

\subsubsection{g:Setmatrix()}
La méthode \cmd{g:Setmatrix(M)} permet d'affecter la matrice \argu{M} à la matrice de transformation du graphe \emph g.

\subsubsection{g:Shift()}
La méthode \cmd{g:Shift(v)} modifie la matrice de transformation du graphe \emph g en la composant avec la matrice de la translation de vecteur \argu{v} qui doit être un complexe.

\begin{demo}{Utilisation de Shift, Rotate et Scale}
\begin{luadraw}{name=free_art}
local ld = luadraw
local du = math.sqrt(2)/2
local g = ld.graph:new{window={1-du,4+du,1-du,4+du},
    margin={0,0,0,0},size={7,7}}
local i = ld.cpx.I
g:Linestyle("noline")
g:Filloptions("full","Navy",0.1)
for X = 1, 4 do
    for Y = 1, 4 do
        g:Savematrix()
        g:Shift(X+i*Y); g:Rotate(45)
        for k = 1, 25 do
            g:Dsquare((1-i)/2,(1+i)/2,1)
            g:Rotate(7); g:Scale(0.9)
        end
        g:Restorematrix()
    end
end
g:Show()
\end{luadraw}
\end{demo}

\subsection{Changement de vue. Changement de repère}

\paragraph{Changement de vue : } lors de la création d'un nouveau graphique, par exemple :
\begin{Luacode}
local ld = luadraw
local g = ld.graph:new{window={-5,5,-5,5},size={10,10}}
\end{Luacode}
L'option \opt{window=\{xmin,xmax,ymin,ymax\}} fixe la vue pour le graphique \emph{g}, ce sera le pavé $[x_{\mathrm{min}}; x_{\mathrm{max}}]\times [y_{\mathrm{min}}; y_{\mathrm{max}}]$ de $\mathbf R^2$, et tous les tracés vont être clippés par cette fenêtre (sauf les labels qui peuvent débordés dans les marges, mais pas au-delà).
Il est possible, à l'intérieur de ce pavé, de définir un autre pavé pour faire une nouvelle vue, avec la méthode \cmd{g:Viewport(x1, x2, y1, y2)}. Les valeurs de \argu{x1}, \argu{x2}, \argu{y1}, \argu{y2} se réfèrent la fenêtre initiale définie par l'option \opt{window}. À partir de là, tout ce qui sort de cette nouvelle zone va être clippé, et la matrice du graphe est réinitialisée à l'identité, par conséquent il faut sauvegarder auparavant les paramètres graphiques courants :
\begin{Luacode}
g:Saveattr()
g:Viewport(x1,x2,y1,y2)
\end{Luacode}
Pour revenir à la vue précédente avec la matrice précédente, il suffit d'effectuer une restauration des paramètres graphiques avec la méthode \cmd{g:Restoreattr()}.

\paragraph{Attention : } à chaque instruction \cmd{g:Saveattr()} doit correspondre une instruction \cmd{g:Restoreattr()}, sinon il y aura une erreur à la compilation.

\paragraph{Changement de repère : } on peut changer le système de coordonnées de la vue courante avec la méthode \cmdln{g:Coordsystem(x1, x2, y1, y2 \fac{, ortho}).}
Cette méthode va modifier la matrice du graphe de sorte que tout se passe comme si la vue courante correspondait au pavé $[x_1;x_2]\times[y_1;y_2]$, l'argument booléen facultatif \argu{ortho} indique si le nouveau repère doit être orthonormé ou non (\false par défaut). Comme la matrice du graphe est modifiée il est préférable de sauvegarder les paramètres graphiques avant, et de les restaurer ensuite. Cela peut servir par exemple à faire plusieurs figures dans le graphique en cours.

\begin{demo}{Classification des points d'une courbe paramétrée}
\begin{luadraw}{name=viewport_changewin}
local ld = luadraw
local g = ld.graph:new{window={-5,5,-5,5},size={10,10}}
local i, Z = ld.cpx.I, ld.cpx.Z
g:Labelsize("tiny")
local styleA = "\\tikzset{->-/.style={decoration={markings, mark="
local styleB = "at position #1 with {\\arrow{>}}}, postaction={decorate}}}"
g:Writeln(styleA..styleB)
g:Dline({0,1},"dashed,gray"); g:Dline({0,i},"dashed,gray")
local legende = {"Point ordinaire", "Point d'inflexion", "Rebroussement première espèce",
"Rebroussement seconde espèce"}
local A, B, C =(1+i)*0.75, 0.75, 0
local A2, B2 ={-1.25+i*0.5,-0.75-i*0.5,1.25-0.5*i, 0.5+i}, {-0.75,-0.75,0.75,0.75}
local u = {Z(-5,0),Z(0,0),-5-5*i,-5*i}
for k = 1, 4 do
    g:Saveattr(); g:Viewport(u[k].re,u[k].re+5,u[k].im,u[k].im+5)
    g:Coordsystem(-1.4,2.25,-1,1.25)
    g:Composematrix({0,1,1+i}) -- pour pencher l'axe Oy
    g:Dpolyline({{-1,1},{-i*0.5,i}}) -- axes
    g:Lineoptions(nil,"blue",8)
    g:Dpath({A2[k],(B2[k]+2*A2[k])/3,(C+5*B2[k])/6, C,"b"},"->-=0.5")
    g:Dpath({C,(C+5*B)/6,(B+2*A)/3,A,"b"},"->-=0.75")
    g:Dpolyline({{0,0.75},{0,0.75*i}},false,"->,red")
    g:Dlabel(
    legende[k],0.75-0.5*i, {pos="S"},
        "$f^{(p)}(t_0)$",1,{pos="E",node_options="red"},
        "$f^{(q)}(t_0)$",0.75*i,{pos="W",dist=0.05})
    g:Restoreattr()
end
g:Show()
\end{luadraw}
\end{demo}
%
\section{Ajouter ses propres méthodes à la classe ld.graph}

Sans avoir à modifier les fichiers sources Lua associés au paquet \luadrawenv, on peut ajouter ses propres méthodes à la classe \emph{ld.graph}, ou modifier une méthode existante. Ceci n'a d'intérêt que si ces modifications doivent être utilisées dans différents graphiques et/ou différents documents (sinon il suffit d'écrire localement une fonction dans le graphique où on en a besoin).

\subsection{Un exemple}
Dans le graphique de la page \pageref{champ}, nous avons dessiné un champ de vecteurs, pour cela on a écrit une fonction qui calcule les vecteurs avant de faire le dessin, mais cette fonction est locale. On pourrait en faire une fonction globale dans l'espace de noms \emph{luadraw}, elle serait alors utilisable dans tout le document, mais pas dans un autre document !

Pour généraliser cette fonction, on va devoir créer un fichier Lua qui pourra ensuite être importé dans des documents en cas de besoin. Pour rendre l'exemple un peu consistant, on va créer un fichier qui va définir une fonction qui calcule les vecteurs d'un champ, et qui va ajouter à la classe \emph{ld.graph} deux nouvelles méthodes : une pour dessiner un champ de vecteurs d'une fonction $f\colon(x,y)\to f(x,y)\in \mathbf R^2$, on la nommera \cmd{ld.graph:Dvectorfield}, et une autre pour dessiner un champ de gradient d'une fonction $f\colon(x,y)\to f(x,y) \in \mathbf R$, on la nommera \cmd{ld.graph:Dgradientfield}. Du coup nous appellerons ce fichier : \emph{luadraw\_fields.lua}.

\paragraph{Contenu du fichier :}
\begin{Luacode}
local ld = luadraw
local graph = ld.graph
local cpx = ld.cpx
local Z = cpx.Z

function ld.field(f,x1,x2,y1,y2,grid,long)
-- champ de vecteurs dans le pavé [x1,x2]x[y1,y2]
-- f fonction de deux variables à valeurs dans R^2
-- grid = {nbx, nby} : nombre de vecteurs suivant x et suivant y
-- long = longueur d'un vecteur
    if grid == nil then grid = {25,25} end
    local deltax, deltay = (x2-x1)/(grid[1]-1), (y2-y1)/(grid[2]-1) -- pas suivant x et y
    if long == nil then long = math.min(deltax,deltay) end -- longueur par défaut
    local vectors = {} -- contiendra la liste des vecteurs
    local x, y, v = x1 
    for _ = 1, grid[1] do -- parcours suivant x
        y = y1
        for _ = 1, grid[2] do -- parcours suivant y
            v = f(x,y) -- on suppose que v est bien défini
            v = Z(v[1],v[2]) -- passage en complexe
            v = cpx.normalize(v)
            if v ~= nil then
                table.insert(vectors, {Z(x,y)-long/2*v, Z(x,y)+long/2*v} ) -- on ajoute le vecteur
            end
            y = y+deltay
        end
        x = x+deltax
    end
    return vectors -- on renvoie le résultat (ligne polygonale)
end

function graph:Dvectorfield(f,args)
-- dessine un champ de vecteurs
-- f fonction de deux variables à valeurs dans R^2
-- args table à 4 champs :
-- { view={x1,x2,y1,y2}, grid={nbx,nby}, length=, draw_options=""}
    args = args or {}
    local view = args.view or {self:Xinf(),self:Xsup(),self:Yinf(),self:Ysup()} -- repère utilisateur par défaut
    local vectors = ld.field(f,view[1],view[2],view[3],view[4],args.grid,args.length) -- calcul du champ
    self:Dpolyline(vectors,false,args.draw_options,view) -- le dessin (ligne polygonale non fermée)
end

function graph:Dgradientfield(f,args)
-- dessine un champ de gradient
-- f fonction de deux variables à valeurs dans R
-- args table à 4 champs :
-- { view={x1,x2,y1,y2}, grid={nbx,nby}, length=, draw_options=""}
    local h = 1e-6
    local grad_f = function(x,y) -- fonction gradient de f
        return { (f(x+h,y)-f(x-h,y))/(2*h), (f(x,y+h)-f(x,y-h))/(2*h) }
    end
    self:Dvectorfield(grad_f,args) -- on utilise la méthode précédente
end
\end{Luacode}

\subsection{Comment importer le fichier}

Il y a deux méthodes pour cela :

\begin{enumerate}
    \item Avec l'instruction Lua \emph{dofile}. On peut l'écrire par exemple dans le préambule après la déclaration du paquet :
    \begin{TeXcode}
    \usepackage[]{luadraw}
    \directlua{dofile("<chemin>/luadraw_fields.lua")}
    \end{TeXcode}
    Bien entendu, il faudra remplacer \verb|<chemin>| par le chemin d'accès à ce fichier. 
    
    L'instruction \verb|\directlua{dofile("<chemin>/luadraw_fields.lua")}| peut être placée ailleurs dans le document pourvu que ce soit après le chargement du paquet (sinon la classe \emph{ld.graph} ne sera pas reconnue lors de la lecture du fichier). On peut aussi placer l'instruction \verb|dofile("<chemin>/luadraw_fields.lua")| dans un environnement \emph{luacode}, et donc en particulier dans un environnement \luadrawenv.
    
    Dès que le fichier est importé, les nouvelles méthodes sont disponibles pour la suite du document.
    
    Cette façon de procéder a au moins deux inconvénients : il faut se souvenir du chemin à chaque utilisation, et d'autre part l'instruction \emph{dofile} ne vérifie par si le fichier a déjà été lu. Pour ces raisons, on préférera la méthode suivante.
    
    \item Avec l'instruction Lua \emph{require}. On peut l'écrire par exemple dans le préambule après la déclaration du paquet :
    \begin{TeXcode}
    \usepackage[]{luadraw}
    \directlua{require "luadraw_fields"}
    \end{TeXcode}
    On remarquera l'absence du chemin (et l'extension lua est inutile).
    
    L'instruction \verb|\directlua{require "luadraw_fields"}| peut être placée ailleurs dans le document pourvu que ce soit après le chargement du paquet (sinon la classe \emph{graph} ne sera pas reconnue lors de la lecture du fichier). On peut aussi placer l'instruction \verb|require "luadraw_fields"| dans un environnement \emph{luacode}, et donc en particulier dans un environnement \luadrawenv.
    
    L'instruction \emph{require} vérifie si le fichier a déjà été chargé ou non, ce qui est préférable. Mais il faut cependant que Lua soit capable de trouver ce fichier, et le plus simple pour cela est qu'il soit quelque part dans une arborescence connue de TeX. On peut par exemple créer dans son \emph{texmf} local le chemin suivant :
    \begin{TeXcode}
    texmf/tex/lualatex/myluafiles/
    \end{TeXcode}
    puis copier le fichier \emph{luadraw\_fields.lua} dans le dossier \emph{myluafiles}.
\end{enumerate}

\begin{demo}[fields]{Utilisation des nouvelles méthodes}
\begin{luadraw}{name=fields}
require "luadraw_fields" -- import des nouvelles méthodes
local ld = luadraw
local g = ld.graph:new{window={0,21,0,10},size={16,10}}
local i = ld.cpx.I
g:Labelsize("footnotesize")
local f = function(x,y) return {x-x*y,-y+x*y} end -- Volterra
local F = function(x,y) return x^2+y^2+x*y-6 end
local H = function(t,Y) return f(Y[1],Y[2]) end
-- graphique du haut
g:Saveattr();g:Viewport(0,10,0,10);g:Coordsystem(-5,5,-5,5)
g:Dgradbox({-4.5-4.5*i,4.5+4.5*i,1,1}, {originloc=0,originnum={0,0},grid=true,
    title="gradient field, $f(x,y)=x^2+y^2+xy-6$"})
g:Arrows("->"); g:Lineoptions(nil,"blue",6)
g:Dgradientfield(F,{view={-4,4,-4,4},grid={15,15},long=0.5})
g:Arrows("-"); g:Lineoptions(nil,"Crimson",12); g:Dimplicit(F, {view={-4,4,-4,4}})
g:Restoreattr()
-- graphique du bas
g:Saveattr();g:Viewport(11,21,0,10);g:Coordsystem(-5,5,-5,5)
g:Dgradbox({-4.5-4.5*i,4.5+4.5*i,1,1}, {originloc=0,originnum={0,0},grid=true,
    title="vector field, $f(x,y)=(x-xy,-y+xy)$"})
g:Arrows("->"); g:Lineoptions(nil,"blue",6); g:Dvectorfield(f,{view={-4,4,-4,4}})
g:Arrows("-");g:Lineoptions(nil,"Crimson",12)
g:Dodesolve(H,0,{2,3},{t={0,50},out={2,3},nbdots=250})
g:Restoreattr()
g:Show()
\end{luadraw}
\end{demo}

\subsection{Modifier une méthode existante}

Prenons par exemple la méthode \cmd{g:DplotXY(X, Y \fac{, draw\_options})} qui prend comme arguments deux listes (tables) de réels et dessine la ligne polygonale formée par les points de coordonnées $(X[k],Y[k])$. Nous allons la modifier afin qu'elle prenne en compte le cas où \argu{X} est une liste de noms (chaînes), dans ce cas, on affichera les noms sous l'axe des abscisses (avec l'abscisse $k$ pour le k\ieme{} nom) et on dessinera la ligne polygonale formée par les points de coordonnées $(k,Y[k])$, sinon on fera comme l'ancienne méthode. Il suffit pour cela de réécrire la méthode (dans un fichier Lua pour pouvoir ensuite l'importer) :

\begin{Luacode}
local ld = luadraw
local cpx = ld.cpx
local Z = cpx.Z

function ld.graph:DplotXY(X,Y,draw_options)
-- X est une liste de réels ou de chaînes
-- Y est une liste de réels de même longueur que X 
    local L = {} -- liste des points à dessiner
    if type(X[1]) == "number" then -- liste de réels
        for k,x in ipairs(X) do
            table.insert(L,Z(x,Y[k]))
        end
    else
        local noms = {} -- liste des labels à placer
        for k = 1, #X do
            table.insert(L,Z(k,Y[k]))
            insert(noms,{X[k],k,{pos="E",node_options="rotate=-90"}})
        end
        self:Dlabel(table.unpack(noms)) --dessin des labels
    end
    self:Dpolyline(L,draw_options) -- dessin de la courbe
end
\end{Luacode}

Dès que le fichier sera importé, cette nouvelle définition va écraser l'ancienne (pour toute la suite du document). Bien entendu on pourrait imaginer ajouter d'autres options sur le style de tracé par exemple (ligne, bâtons, points ...).

\begin{demo}{Modification d'une méthode existante}
\begin{luadraw}{name=newDplotXY}
require "luadraw_fields" -- import de la méthode modifiée
local ld = luadraw
local g = ld.graph:new{window={-0.5,11,-1,20}, margin={0.5,0.5,0.5,1}, size={10,10,0}}
g:Labelsize("scriptsize")
local X, Y = {}, {} -- on définit deux listes X et Y, on pourrait aussi les lire dans un fichier
for k = 1, 10 do
    table.insert(X,"nom"..k)
    table.insert(Y,math.random(1,20))
end
ld.defaultlabelshift = 0
g:Daxes({0,1,2},{limits={{0,10},{0,20}}, labelpos={"none","left"},arrows="->", grid=true})
g:DplotXY(X,Y,"line width=0.8pt, blue")
g:Show()
\end{luadraw}
\end{demo}
%

\chapter{Dessin 3D}

\begin{center}
\captionof{figure}{Point col en $M(0,0,0)$ ($z=x^2-y^2$)}\label{pointcol}\par
\begin{luadraw}{name=point_col}
local ld = luadraw
local pt3d = ld.pt3d
local Origin, vecI, vecJ, vecK, M = pt3d.Origin, pt3d.vecI, pt3d.vecJ, pt3d.vecK, pt3d.M

local g = ld.graph3d:new{window3d={-2,2,-2,2,-4,4}, window={-3.5,3,-5,5}, size={8,9,0}, viewdir={120,60}}
local S = ld.cartesian3d(function(u,v) return u^2-v^2 end, -2,2,-2,2,{20,20}) -- surface d'équation z=x^2-y^2
local Tx = g:Intersection3d(S, {Origin,vecI}) --intersection entre S et le plan yOz
local Ty = g:Intersection3d(S, {Origin,vecJ}) --intersection entre S et le plan xOz
g:Dboxaxes3d({grid=true,gridcolor="gray",fillcolor="LightGray",drawbox=true})
g:Dfacet(S,{mode=ld.mShadedOnly,color="ForestGreen"}) -- dessin de la surface
g:Dedges(Tx, {color="Crimson", hidden=true, width=8}) -- dessin de l'intersection avec yOz
g:Dedges(Ty, {color="Navy",hidden=true, width=8}) -- dessin de l'intersection avec xOz
g:Dpolyline3d( {M(2,0,4),M(-2,0,4),M(-2,0,-4)}, "Navy,line width=.8pt")
g:Dpolyline3d( {M(0,-2,4),M(0,2,4),M(0,2,-4)}, "Crimson,line width=.8pt")
g:Show()
\end{luadraw}
\end{center}



\input{body-fr/luadraw-doc3d-section-1-Introduction.tex}%
\section{La classe pt3d}

\subsection{Représentation des points et vecteurs}

L'espace usuel est $\mathbf R^3$, les points et les vecteurs sont donc des triplets de réels, on les appellera: points 3D. La classe \emph{pt3d} (qui est automatiquement chargée) permet de gérer les points 3D, les opérations possibles, et un certain nombre de méthodes et de constantes.

\begin{itemize}
    \item  Quatre triplets portent un nom spécifique (variables globales), il s'agit de :
    \begin{itemize}
        \item \varglob{pt3d.Origin}, qui représente le triplet $(0,0,0)$.
        \item \varglob{pt3d.vecI}, qui représente le triplet $(1,0,0)$.
        \item \varglob{pt3d.vecJ}, qui représente le triplet $(0,1,0)$.
        \item \varglob{pt3d.vecK}, qui représente le triplet $(0,0,1)$.
    \end{itemize}
    \item  Pour créer un point 3D, il y a trois méthodes :
        \begin{itemize}
            \item Définition en cartésien : la fonction \cmd{pt3d.M(x, y, z)} renvoie le triplet $(x,y,z)$. On peut également obtenir ce triplet en faisant : \code{x*vecI+y*vecJ+z*vecK}.
            \item Définition en cylindrique : la fonction \cmd{pt3d.Mc(r, theta, z)} (angle exprimé en radians) renvoie le triplet $(r\cos(\theta),r\sin(\theta),z)$.
            \item Définition en sphérique : la fonction \cmd{pt3d.Ms(r, theta, phi)} renvoie le triplet $(r\cos(\theta)\sin(\varphi), r\sin(\theta)\sin(\varphi),r\cos(\varphi))$ (angles exprimés en radians).
        \end{itemize}
    Accès aux composantes d'un point 3D : si une variable $A$ désigne un point 3D, alors ses trois composantes sont $A.x$, $A.y$ et $A.z$.
    
    Pour tester si une variable $A$ désigne un point 3D, on dispose de la fonction \cmd{pt3d.isPoint3d()} qui renvoie un booléen.
    
    Conversion : pour convertir un réel ou un complexe en point 3D, on dispose de la fonction \cmd{pt3d.toPoint3d()}.
\end{itemize}

\subsection{Opérations sur les points 3D}

Ces opérations sont les opérations usuelles avec les symboles usuels :
\begin{itemize}
    \item L'addition (+), la différence (-), l'opposé (-).
    \item Le produit par un scalaire, si $k$ et un réel, \code{k*M(x,y,z)} renvoie \val{M(k*a,k*y,k*z)}.
    \item On peut diviser un point 3D par un scalaire, par exemple, si $A$ et $B$ sont deux points 3D, alors le milieu s'écrit simplement $(A+B)/2$.
    \item On peut tester l'égalité de deux points 3D avec le symbole =.
\end{itemize}

\subsection{Méthodes de la classe \emph{pt3d}}

Celles-ci sont :
\begin{itemize}
    \item \cmd{pt3d.abs(u)} : renvoie la norme euclidienne du point 3D \argu{u}.
    
    \item \cmd{pt3d.abs2(u)} : renvoie la norme euclidienne au carré du point 3D \argu{u}.
    
    \item \cmd{pt3d.N1(u)} : renvoie la norme $1$ du point 3D $u$. Si $u=M(x,y,z)$, alors \code{pt3d.N1(u)} renvoie \val{$|x|+|y|+|z|$}.
    
    \item \cmd{pt3d.dot(u, v)} : renvoie le produit scalaire entre les vecteurs (points 3D) \argu{u} et \argu{v}.
    
    \item \cmd{pt3d.det(u, v, w)} : renvoie le déterminant entre les vecteurs (points 3D) \argu{u}, \argu{v} et \argu{w}.
    
    \item \cmd{pt3d.prod(u, v)} : renvoie le produit vectoriel entre les vecteurs (points 3D) \argu{u} et \argu{v}.
    
    \item \cmd{pt3d.angle3d(u, v \fac{, epsilon})} : renvoie l'écart angulaire (en radians) entre les vecteurs (points 3D) \argu{u} et \argu{v}. supposés non nuls. L'argument (facultatif) \argu{epsilon} vaut $0$ par défaut, il indique à combien près se fait un certain test d'égalité sur un flottant.
    \item \cmd{pt3d.normalize(u)} : renvoie le vecteur (point 3D) \argu{u} normalisé (renvoie \nil si \argu{u} est nul).
    
    \item \cmd{pt3d.round(u, nbDeci)} : renvoie un point 3D dont les composantes sont celles du point 3D \argu{u} arrondies avec \argu{nbDeci} décimales.
    
    \item \cmd{pt3d.isobar3d(L)} : renvoie l'isobarycentre des points 3D de la liste (table) \argu{L} (les éléments de \argu{L} qui ne sont pas des points 3D sont ignorés).
    
    \item \cmd{pt3d.insert3d(L, A \fac{, epsilon})} : cette fonction insère le point 3D \argu{A} dans la liste \argu{Lu} qui doit être une \textbf{variable} (et qui sera donc modifiée). Le point \argu{A} est inséré \textbf{sans doublon} et la fonction renvoie sa position (indice) dans la liste \argu{L} après insertion. L'argument (facultatif) \argu{epsilon} vaut $0$ par défaut, il indique à combien près se font les comparaisons.
\end{itemize}

\subsection{Afficher une variable dans le terminal}

L'instruction \cmd{ld.whatis(variable \fac{, msg})} affiche dans le terminal lors de la compilation, le type de la \argu{variable} ainsi que son contenu. Les types reconnus sont, les types prédéfinis plus : \emph{complex number}, \emph{list of (complex) numbers}, \emph{list of lists of (complex) numbers}, \emph{3D point}, \emph{list of 3D points}, \emph{list of lists of 3D points}. L'argument \argu{msg} est une chaîne optionnelle (vide par défaut) qui est affichée avec le type pour repérer la variable dans le terminal.
%
\section{Méthodes graphiques élémentaires}

Toutes les méthodes graphiques 2D s'appliquent. À cela s'ajoute la possibilité de dessiner dans l'espace des lignes polygonales, des segments, droites, courbes, chemins, points, labels, plans, solides. Avec les solides vient également la notion de facettes que l'on ne trouvait pas en 2D.

Les méthodes graphiques 3D vont calculer automatiquement la projection sur le plan de l'écran, après avoir appliquer aux objets la matrice de transformation 3D associée au graphique (qui est l'identité par défaut), ce sont ensuite les méthodes graphiques 2D qui prendront le relai.

La méthode qui applique la matrice 3D et fait la projection sur l'écran (plan passant par l'origine et normal au vecteur unitaire dirigé vers l'observateur et défini par les angles de vue), est : \cmd{g:Proj3d(L)} où \argu{L} est soit un point 3D, soit une liste de points 3D, soit une liste de listes de points 3D. Cette fonction renvoie des complexes (affixes des projetés sur l'écran).

\textbf{Attention} : lorsque la matrice 3D du graphe est une transformation affine non linéaire, le projeté sur l'écran d'un vecteur $u$ de l'espace n'est pas \code{g:Proj3d(u)}, mais \code{g:Proj3d(A+u)-g:Proj3d(A)} où $A$ désigne un point quelconque de l'espace. Pour éviter ces calculs, la méthode \cmd{g:Proj3dV()} a été  introduite, elle fait la projection des \textbf{vecteurs} sur l'écran, et renvoie des complexes (affixes des projetés sur l'écran).

\subsection{Dessin aux traits}

\subsubsection{Ligne polygonale : Dpolyline3d}

La méthode \cmd{g:Dpolyline3d(L \fac{, close, draw\_options, clip})} (où \emph{g} désigne le graphique en cours de création), \argu{L} est une ligne polygonale 3D (liste de listes de points 3D), \argu{close} un argument facultatif qui vaut \true ou \false, indiquant si la ligne doit être refermée ou non (\false par défaut), et \argu{draw\_options} est une chaîne de caractères qui sera passée directement à l'instruction \drawcmd dans l'export. L'argument \argu{clip} vaut \false par défaut, il indique si la ligne \argu{L} doit être clippée avec la fenêtre 3D courante.
    
\subsubsection{Angle droit : Dangle3d}

La méthode \cmd{g:Dangle3d(B, A, C \fac{, r, draw\_options, clip})} dessine l'angle \(BAC\) avec un parallélogramme (deux côtés seulement sont dessinés), l'argument facultatif \argu{r} précise la longueur d'un côté ($0.25$ par défaut). Le parallélogramme est dans le plan défini par les points \argu{A}, \argu{B} et \argu{C}, ceux-ci ne doivent donc pas être alignés. L'argument \argu{draw\_options} est une chaîne (vide par défaut) qui sera passée telle quelle à l'instruction  \drawcmd . L'argument \argu{clip} vaut \false par défaut, il indique si le tracé doit être clippé avec la fenêtre 3D courante.
    
\subsubsection{Segment : Dseg3d}

La méthode \cmd{g:Dseg3d(seg \fac{, scale, draw\_options, clip})} dessine le segment défini par l'argument \argu{seg} qui doit être une liste de deux points 3D. L'argument facultatif \argu{scale} ($1$ par défaut) est un nombre qui permet d'augmenter ou réduire la longueur du segment (la longueur naturelle est multipliée par \argu{scale}). L'argument \argu{draw\_options} est une chaîne (vide par défaut) qui sera passée telle quelle à l'instruction \drawcmd . L'argument \argu{clip} vaut \false par défaut, il indique si le tracé doit être clippé avec la fenêtre 3D courante.
    
\subsubsection{Droite : Dline3d}

La méthode \cmd{g:Dline3d(d \fac{, draw\_options, clip})} trace la droite \argu{d}, celle-ci est une liste du type \argu{d}=$\{A,u\}$ où $A$ représente un point de la droite (point 3D) et $u$ un vecteur directeur (un point 3D non nul). 

Variante : la méthode \cmd{g:Dline3d(A, B \fac{, draw\_options, clip})} trace la droite passant par les points \argu{A} et \argu{B} (deux points 3D). L'argument \argu{draw\_options} est une chaîne (vide par défaut) qui sera passée telle quelle à l'instruction \drawcmd . L'argument \argu{clip} vaut \false par défaut, il indique si le tracé doit être clippé avec la fenêtre 3D courante.

La méthode \cmd{g:Line3d2seg(d \fac{, scale})} renvoie une table constituée de deux points 3D représentant un segment, ce segment est la partie de la droite \argu{d} à l'intérieur la fenêtre 3D courante. L'argument \argu{scale} ($1$ par défaut) permet de faire varier la taille de ce segment. Lorsque la fenêtre est trop petite l'intersection peut être vide.

 \subsubsection{Arc de cercle : Darc3d}
 
\begin{itemize}
    \item La méthode \cmd{g:Darc3d(B, A, C, r, sens \fac{, normal, draw\_options, clip})} dessine un arc de cercle de centre \argu{A} (point 3D), de rayon \argu{r}, allant de \argu{B} (point 3D) vers \argu{C} (point 3D) dans le sens direct si l'argument \argu{sens} vaut $1$, le sens inverse sinon. Cet arc est tracé dans le plan contenant les trois points \argu{A}, \argu{B} et \argu{C}, lorsque ces trois points sont alignés il faut préciser l'argument \argu{normal} (point 3D non nul) qui représente un vecteur normal au plan. Ce plan est orienté par le produit vectoriel $\vec{AB}\wedge\vec{AC}$ ou bien par le vecteur \argu{normal} si celui-ci est précisé. L'argument \argu{draw\_options} est une chaîne (vide par défaut) qui sera passée telle quelle à l'instruction \drawcmd . L'argument \argu{clip} vaut \false par défaut, il indique si le tracé doit être clippé avec la fenêtre 3D courante.
    
    \item La fonction \cmd{ld.arc3d(B, A, C, r, sens \fac{, normal})} renvoie la liste des points de cet arc (ligne polygonale 3D). 
    
    \item La fonction \cmd{ld.arc3db(B, A, C, r, sens \fac{, normal})} renvoie cet arc sous forme d'un chemin 3D (voir \emph{Dpath3d} page \pageref{Dpath3d}) utilisant des courbes de Bézier.
\end{itemize}

\subsubsection{Cercle : Dcircle3d}

\begin{itemize}
    \item La méthode \cmd{g:Dcircle3d(I, R, normal \fac{, draw\_options, clip})} trace le cercle de centre \argu{I} (point 3D) et de rayon \argu{R}, dans le plan contenant \argu{I} et normal au vecteur défini par l'argument \argu{normal} (point 3D non nul). L'argument \argu{draw\_options} est une chaîne (vide par défaut) qui sera passée telle quelle à l'instruction \drawcmd . L'argument \argu{clip} vaut \false par défaut, il indique si le tracé doit être clippé avec la fenêtre 3D courante. 
    
    Autre syntaxe possible :  \cmd{g:Dcircle3d(C \fac{, draw\_options, clip})} où \argu{C}=\{\argu{I},\argu{R},\argu{normal}\}.
    
    \item La fonction \cmd{ld.circle3d(I, R, normal)} renvoie la liste des points de ce cercle (ligne polygonale 3D). 
    
    \item La fonction \cmd{ld.circle3db(I, R, normal)} renvoie ce cercle sous forme d'un chemin 3D (voir \emph{Dpath3d} page \pageref{Dpath3d}) utilisant des courbes de Bézier.
\end{itemize}
    
\subsubsection{Chemin 3D : Dpath3d}\label{Dpath3d}

La méthode \cmd{g:Dpath3d(chemin \fac{, draw\_options, clip})} fait le dessin du \argu{chemin}. L'argument \argu{draw\_options} est une chaîne (vide par défaut) qui sera passée telle quelle à l'instruction \drawcmd . L'argument \argu{clip} vaut \false par défaut, il indique si le tracé doit être clippé avec la fenêtre 3D courante. L'argument \argu{chemin} est une liste de points 3D suivis d'instructions (chaînes) fonctionnant \textbf{sur le même principe qu'en 2D}. Instructions disponibles et leur syntaxe, le mot \emph{last} représente le dernier point du morceau précédent:
      \begin{itemize}
      \item \code{p1,"m"} (moveto), permet de commencer une nouvelle composante du chemin au point 3D $p_1$.
      \item \code{p1,...,pn,"l"} (lineto), dessine la ligne polygonale 3D \code{\{last,p1,...,pn\}}.
      \item \code{c1,c2,p2,"b"} (bézier) dessine la courbe de Bézier \code{\{last,c1,c2,p2\}}, où $c_1$ et $c_2$ sont les deux points 3D de contrôle.
      \item \code{p1,n,"c"} (cercle), dessine le cercle de centre $p_1$ et passant par le point \emph{last}, et normal au vecteur 3D $n$.
      \item \code{p1,p2,r,sens,n,"ca"} (arc de cercle), dessine un arc de cercle de centre $p_1$, de rayon $r$, allant de \emph{last} vers $p_2$, dans le sens direct lorsque \emph{sens}=$1$ (et donc dans le sens inverse si \emph{sens}=$-1$). Le vecteur 3D $n$ est optionnel, il indique un vecteur normal au plan du cercle lorsque les points \emph{last}, $p_1$ et $p_2$ sont alignés (et dans ce cas, le vecteur $n$ est obligatoire).
      \item \code{"cl"} (closepath), cette instruction s'utilise seule, elle permet de refermer la composante courante en traçant un segment reliant le dernier point au premier point (de la composante courante).
      \end{itemize}      

Voici par exemple le code de la figure \ref{viewdir}.

\begin{Luacode}
\begin{luadraw}{name=viewdir}
local ld = luadraw
local pt3d = ld.pt3d
local Origin, vecI, vecJ, vecK = pt3d.Origin, pt3d.vecI, pt3d.vecJ, pt3d.vecK 
local g = ld.graph3d:new{ size={8,8} }
local i, M = ld.cpx.I, ld.pt3d.M
local O, A = Origin, M(4,4,4)
local B, C, D, E = ld.pxy(A), ld.px(A), ld.py(A), ld.pz(A) --projeté de A sur le plan xOy et sur les axes
g:Dpolyline3d( {{O,A},{-5*vecI,5*vecI},{-5*vecJ,5*vecJ},{-5*vecK,5*vecK}}, "->") -- axes
g:Dpolyline3d( {{E,A,B,O}, {C,B,D}}, "dashed")
g:Dpath3d( {C,O,B,2.5,1,"ca",O,"l","cl"}, "draw=none,fill=cyan,fill opacity=0.8") --secteur angulaire
g:Darc3d(C,O,B,2.5,1,"->") -- arc de cercle pour theta
g:Dpath3d( {E,O,A,2.5,1,"ca",O,"l","cl"}, "draw=none,fill=cyan,fill opacity=0.8") --secteur angulaire
g:Darc3d(E,O,A,2.5,1,"->") -- arc de cercle pour phi
g:Dballdots3d(O) -- le point origine sous forme d'une petite sphère
g:Labelsize("footnotesize")
g:Dlabel3d(
  "$x$", 5.25*vecI,{}, "$y$", 5.25*vecJ,{}, "$z$", 5.25*vecK,{},
  "vers observateur", A, {pos="E"},
  "$O$", O, {pos="NW"},
  "$\\theta$", (B+C)/2, {pos="N", dist=0.15},
  "$\\varphi$", (A+E)/2, {pos="S",dist=0.25})
g:Dlabel("viewdir=\\{"ortho",$\\theta,\\varphi$\\} (en degrés)",-5*i,{pos="N"}) -- label 2D
g:Show()
\end{luadraw}
\end{Luacode}

\textbf{Conversion} : la fonction \cmd{ld.polyline2path3d(L)} renvoie \argu{L} qui est une liste de points 3D ou une liste de listes de points 3D, sous la forme d'un chemin.

\subsubsection{Plan : Dplane}

\begin{itemize}
    \item La méthode \cmd{g:Dplane(P, V, L1, L2 \fac{, mode, draw\_options})} permet de dessiner les bords du plan \argu{P}=$\{A,u\}$ où $A$ est un point du plan et $u$ un vecteur normal au plan (\argu{P} est donc une table de deux points 3D). L'argument \argu{V} doit être un vecteur non nul du plan \argu{P}, Les arguments \argu{L1} et \argu{L2} sont deux longueurs. La méthode construit un parallélogramme centré sur $A$, dont un côté est $L_1\frac{V}{\|V\|}$ et l'autre $L_2\frac{W}{\|W\|}$ où $W = u\wedge V$. L'argument \argu{mode} est un entier naturel qui indique les bords à tracer. Pour calculer cet entier on utilise les variables prédéfinies : \varglob{ld.top} (=8), \varglob{ld.right} (=4), \varglob{ld.bottom} (=2), \varglob{ld.left} (=1) et \varglob{ld.all} (=15), que l'on peut ajouter entre elles, par exemple :
    \begin{itemize}
        \item \code{mode=ld.bottom+ld.left} : pour les côtés bas et gauche
        \item \code{mode=ld.top+ld.right+ld.bottom} : pour les côtés haut, droit et bas
        \item etc.
    \end{itemize}
    Par défaut le mode vaut \varglob{ld.all} ce qui correspond à \code{mode=ld.top+ld.right+ld.bottom+ld.left}.

    \item La fonction \cmd{ld.plane2rectangle(P \fac{, V, L1, L2})} calcule et renvoie le rectangle (liste de points 3D) dessiné par la méthode \cmd{g:Dplane()}. Le vecteur \argu{V} est facultatif, il permet d'imposer un bord du rectangle (il doit appartenir au plan), s'il est omis la fonction choisira elle-même un vecteur $V$. Les longueurs \argu{L1} et \argu{L2} sont facultatives et valent $5$ par défaut. Lorsque l'argument \argu{L2} et omis, il a implicitement la même valeur que \argu{L1}. Le résultat peut être dessiné avec la méthode \cmd{g:Dpolyline3d()}, ou bien dessiné en tant que facette.
\end{itemize}

\begin{demo}{Dplane, exemple avec mode = left+bottom}
\begin{luadraw}{name=Dplane}
local ld = luadraw
local cpx, pt3d = ld.cpx, ld.pt3d
local Origin, vecI, vecJ, vecK, M = pt3d.Origin, pt3d.vecI, pt3d.vecJ, pt3d.vecK, pt3d.M

local g = ld.graph3d:new{size={8,8},window={-5.25,3,-2.5,2.5},margin={0,0,0,0},border=true}
local i = cpx.I
g:Labelsize("footnotesize")
local A = Origin
local P = {A, vecK}
g:Dplane(P, vecJ, 6, 6, ld.left + ld.bottom)
g:Dcrossdots3d({A,vecK},nil,0.75)
g:Dseg3d({A,A+2*vecK},"->")
g:Dangle3d(-vecJ,A,vecK,0.25)
g:Dpolyline3d({{M(3.5,-3,0),M(3.5,3,0)},{M(3,-3.5,0), M(-3,-3.5,0)}}, "->,line width=0.8pt")
g:Dlabel3d("$A$",A,{pos="E"},
    "$u$",2*vecK,{},
    "$P$", M(3,-3,0),{pos="NE", dir={vecJ,-vecI}},
    "$L_1\\frac{V}{\\|V\\|}$ (bottom)", M(3.5,0,0), {pos="S"},
    "$L_2\\frac{W}{\\|W\\|}$ (left)", M(0,-3.5,0), {pos="N",dir={-vecI,-vecJ}}
)
g:Show()
\end{luadraw}
\end{demo}

\paragraph{Attention} : les notions de haut, droite, bas et gauche sont relatives ! Elles dépendent du sens des vecteurs $u$ (vecteur normal au plan) et $V$ (vecteur donné dans le plan). Le troisième vecteur $W$ est le produit vectoriel $u\wedge V$.

\subsubsection{Courbe paramétrique : Dparametric3d}

\begin{itemize}
\item La fonction \cmd{ld.parametric3d(p, t1, t2 \fac{, nbdots, discont, nbdiv})} fait le calcul des points de la courbe et renvoie une ligne polygonale 3D (pas de dessin).
  \begin{itemize}
    \item L'argument \argu{p} est le paramétrage, ce doit être une fonction d'une variable réelle $t$ et à valeurs dans $\mathbf R^3$ (les images sont des points 3D), par exemple :
    \code{local p=function(t) return Mc(3,t,t/3) end}.
    
    \item  Les arguments \argu{t1} et \argu{t2} sont obligatoires avec \(t_1 < t_2\), ils forment les bornes de l'intervalle pour le paramètre.
    
    \item L'argument \argu{nbdots} est facultatif, c'est le nombre de points (minimal) à calculer, il vaut $40$ par défaut.
    
    \item L'argument \argu{discont} est un booléen facultatif qui indique s'il y a des discontinuités ou non, c'est \false par défaut.
    
    \item L'argument \argu{nbdiv} est un entier positif qui vaut $5$ par défaut et indique le nombre de fois que l'intervalle entre deux valeurs consécutives du paramètre peut être coupé en deux (par dichotomie) lorsque les points correspondants sont trop éloignés.
  \end{itemize}
  
\item La méthode \cmd{g:Dparametric3d(p, options)} fait le calcul des points et le dessin de la courbe paramétrée par \argu{p}. L'argument \argu{options} est une table dont les champs sont les options possibles. Celles-ci sont, avec leur valeur par défaut:

  \begin{itemize}
      \item \opt{t=\{g:Xinf(),g:Xsup()\}},
      \item \opt{nbdots=40}, 
      \item \opt{discont=false},
      \item \opt{nbdiv=5},
      \item \opt{clip=false}, il indique si la courbe doit être clippée avec la fenêtre 3D courante.
      \item \opt{draw\_options=""}, chaîne qui sera transmise telle quelle à l'instruction \drawcmd .
  \end{itemize}
\end{itemize} 

\begin{demo}{Une courbe et ses projections sur trois plans} 
\begin{luadraw}{name=Dparametric3d}
local ld = luadraw
local pt3d = ld.pt3d
local vecI, vecJ, vecK, M, Mc = pt3d.vecI, pt3d.vecJ, pt3d.vecK, pt3d.M, pt3d.Mc

local g = ld.graph3d:new{window3d={-4,4,-4,4,-3,3}, window={-7.5,6.5,-7,6}, size={8,8}}
local pi = math.pi
g:Labelsize("footnotesize")
local p = function(t) return Mc(3,t,t/3) end
local L = ld.parametric3d(p,-2*pi,2*pi,25,false,2)
g:Dboxaxes3d({grid=true,gridcolor="gray",fillcolor="LightGray"})
g:Lineoptions("dashed","red",2)
-- projection sur le plan y=-4
g:Dpolyline3d(ld.proj3d(L,{M(0,-4,0),vecJ}))
-- projection sur le plan x=-4
g:Dpolyline3d(ld.proj3d(L,{M(-4,0,0),vecI}))
-- projection sur le plan z=-3
g:Dpolyline3d(ld.proj3d(L,{M(0,0,-3),vecK}))
-- dessin de la courbe
g:Lineoptions("solid","Navy",8)
g:Dparametric3d(p,{t={-2*pi,2*pi}})
g:Show()
\end{luadraw}
\end{demo}

\subsubsection{Paramétrisation d'une ligne polygonale: \emph{curvilinear\_param3d}}
Soit $L$ une liste de points 3D représentant une ligne continue, il est possible d'obtenir une paramétrisation de cette ligne en fonction d'un paramètre $t$ entre $0$ et $1$ ($t$ est l'abscisse curviligne divisée par la longueur totale de $L$).

La fonction \cmd{ld.curvilinear\_param3d(L \fac{, close})} renvoie une fonction d'une variable $t\in[0;1]$ et à valeurs sur la ligne \argu{L} (points 3D), la valeur en $t=0$ est le premier point de \argu{L}, et la valeur en $t=1$ est le dernier point; cette fonction est suivie d'un nombre qui représente la longueur total de \argu{L}. L'argument optionnel \argu{close} indique si la ligne doit être refermée (\false par défaut).


\subsubsection{Le repère : Dboxaxes3d}

La méthode \cmd{g:Dboxaxes3d(options)} permet de dessiner les trois axes, avec un certain nombre d'options définies dans la table \argu{options}. Ces options sont :
%\def\opt#1{\textcolor{blue}{\texttt{#1}}}%
\begin{itemize}
    \item \opt{xaxe=true}, \opt{yaxe=true} et \opt{zaxe=true} : indique si les axes correspondant doivent être dessinés ou non.

    \item \opt{drawbox=false} : indique si une boite doit être dessinée avec les axes.

    \item \opt{grid=false} : indique si une grille doit être dessinée (une pour $x$, une pour $y$ et une pour $z$). Lorsque cette option vaut \true, on peut utiliser aussi les options suivantes :
        \begin{itemize}
            \item \opt{gridwidth=1}: indique l'épaisseur de trait de la grille en dixième de point,
            \item \opt{gridcolor="black"}: indique la couleur de la grille,
            \item \opt{fillcolor=""}: couleur pour peindre le fond des grilles.
        \end{itemize}
    
    \item \opt{xlimits=\{x1,x2\}}, \opt{ylimits=\{y1,y2\}}, \opt{zlimits=\{z1,z2\}} : permet de définir les trois intervalles utilisés pour les axes. Par défaut ce sont les valeurs fournies à l'argument \opt{window3d} à la création du graphe.

    \item \opt{xgradlimits=\{x1,x2\}}, \opt{ygradlimits=\{y1,y2\}}, \opt{zgradlimits=\{z1,z2\}} : permet de définir les trois intervalles de graduation sur les axes. Par défaut ces options ont la valeur \val{"auto"}, ce qui veut dire qu'elles prennent les mêmes valeurs que \opt{xlimits}, \opt{ylimits} et \opt{zlimits}.
    
    \item \opt{xyzstep=1} : indique le pas des graduations sur les trois axes.
    
    \item \opt{xstep=xyzstep}, \opt{ystep=xyzstep}, \opt{zstep=xyzstep} : indique le pas des graduations sur chaque axe (valeur de \opt{xyzstep} par défaut).

    \item \opt{xyzticks=0.2} : indique la longueur des graduations.

    \item \opt{labels=true} : indique si la valeur des graduations doit être affichée ou non.
    
    \item \opt{xlabelsep=0.25}, \opt{ylabelsep=0.25}, \opt{zlabelsep=0.25} : indique la distance entre les labels et les graduations.
    
    \item \opt{xlabelstyle=<style courant>}, \opt{ylabelstyle=<style courant>}, \opt{zlabelstyle=<style courant>} : indique le style des labels, c'est à dire la position par rapport au point d'ancrage. Par défaut c'est le style en cours qui s'applique.

    \item \opt{xlegend="\$x\$"}, \opt{ylegend="\$y\$"}, \opt{zlegend="\$z\$"} : permet de définir une légende pour les axes.
    
    \item \opt{xlegendsep=0.5}, \opt{ylegendsep=0.5}, \opt{zlegendsep=0.5} : indique la distance entre les legendes et les graduations.     
\end{itemize}

\subsection{Points et labels}

\subsubsection{Points 3D : Ddots3d, Dballdots3d, Dcrossdots3d}

Il y a trois possibilités de dessiner des points 3D. Pour les deux premières, l'argument \argu{L} peut être soit un seul point 3D, soit une liste (une table) de points 3D, soit une liste de listes de points 3D :

\begin{itemize}
    \item La méthode \cmd{g:Ddots3d(L \fac{, mark\_options, clip})}. Le principe est le même que dans la version 2D, les points sont dessinés dans la couleur courante du tracé de lignes avec le style courant. L'argument \argu{mark\_options} est une chaîne de caractères facultative qui sera passée telle quelle à l'instruction \drawcmd (modifications locales). L'argument \argu{clip} vaut \false par défaut, il indique si le tracé doit être clippé avec la fenêtre 3D courante.
    
    \item La méthode \cmd{g:Dballdots3d(L \fac{, color, scale, clip})} dessine les points de \argu{L} sous forme d'une sphère. L'argument facultatif \argu{color} précise la couleur de la sphère (\val{"black"} par défaut), et l'argument facultatif \argu{scale} permet de jouer sur la taille de la sphère ($1$ par défaut).
    
    \item La méthode \cmd{g:Dcrossdots3d(L \fac{, color, scale, clip})} dessine les points de \argu{L} sous forme d'une croix plane. L'argument \argu{L} est une liste de la forme \{point 3D, vecteur normal\} ou \{ \{point3d, vecteur normal\}, \{point3d, vecteur normal\}, ...\}. Pour chaque point 3D, le vecteur normal associé permet de déterminer le plan contenant la croix. L'argument facultatif \argu{color} précise la couleur de la sphère (\val{"black"} par défaut), et l'argument facultatif \argu{scale} permet de jouer sur la taille de la sphère ($1$ par défaut).
\end{itemize}

\begin{demo}{Un tétraèdre et les centres de gravité de chaque face}
\begin{luadraw}{name=Ddots3d}
local ld = luadraw
local pt3d = ld.pt3d
local vecI, vecJ, vecK, M = pt3d.vecI, pt3d.vecJ, pt3d.vecK, pt3d.M

local g = ld.graph3d:new{viewdir={15,60},bbox=false,size={8,8}}
local A, B, C, D = 4*M(1,0,-0.5), 4*M(-1/2,math.sqrt(3)/2,-0.5), 4*M(-1/2,-math.sqrt(3)/2,-0.5), 4*M(0,0,1)
local u, v, w = B-A, C-A, D-A
-- centres de gravité faces cachées
for _, F in ipairs({{A,B,C},{B,C,D}}) do
local G, u = pt3d.isobar3d(F), pt3d.prod(F[2]-F[1],F[3]-F[1])
g:Dcrossdots3d({G,u}, "blue",0.75)
g:Dpolyline3d({{F[1],G,F[2]},{G,F[3]}},"dotted")
end
-- dessin du tétraèdre construit sur A, B, C et D
g:Dpoly( ld.tetra(A,u,v,w),{mode=mShaded,opacity=0.7,color="Crimson"})
-- centres de gravité faces visibles
for _, F in ipairs({{A,B,D},{A,C,D}}) do
    local G, u = pt3d.isobar3d(F), pt3d.prod(F[2]-F[1],F[3]-F[1])
    g:Dcrossdots3d({G,u}, "blue",0.75)
    g:Dpolyline3d({{F[1],G,F[2]},{G,F[3]}},"dotted")
end
g:Dballdots3d({A,B,C,D}, "orange") --sommets
g:Show()
\end{luadraw}
\end{demo}

\subsubsection{Labels 3D : Dlabel3d}

La méthode pour placer un label dans l'espace est :

cmdln{g:Dlabel3d(text1, anchor1, options1, text2, anchor2, options2, \ldots).}

    \begin{itemize}
    \item  Les arguments \argu{text1}, \argu{text2}, \ldots, sont des chaînes de caractères, ce sont les labels.
    \item  Les arguments \argu{anchor1}, \argu{anchor2}, \ldots, sont des points 3D représentant les points d'ancrage des labels.
    \item  Les arguments \argu{options1}, \argu{options2}, \ldots,  permettent de définir localement les options des labels, ces options sont (avec leur valeur par défaut) :

        \begin{itemize}
            \item \opt{pos="center"} : indique la position du label dans le plan de l'écran par rapport au point d'ancrage, il peut valoir \val{"N"} pour nord, \val{"NE"} pour nord-est, \val{"NW"} pour nord-ouest, ou encore \val{"S"}, \val{"SE"}, \val{"SW"}. Par défaut, il vaut \val{"center"}, et dans ce cas le label est centré sur le point  d'ancrage.
            
            \item \opt{dist=0} : c'est la distance (en cm) entre le label et son point d'ancrage lorsque \opt{pos} n'est pas égal a \val{"center"}.
            
            \item \opt{dir=\{\}} : indique la direction de l'écriture dans l'espace, la liste vide indique le sens par défaut. Plus généralement, \code{dir=\{dirX,dirY,dep\}}, les 3 valeurs \emph{dirX}, \emph{dirY} et \emph{dep} sont trois points 3D représentant 3 vecteurs, les deux premiers indiquent le sens de l'écriture, le troisième un déplacement (translation) du label par rapport au point d'ancrage.
            
            \item \opt{node\_options=""} : chaîne destinée à recevoir des options qui seront directement passées à TikZ dans l'instruction \emph{\textbackslash node[]}.
            
            \item Les labels sont dessinés dans la couleur courante du texte du document, mais on peut changer de couleur avec l'option \opt{node\_options} en mettant par exemple : \opt{node\_options="color=blue"}.
            
            \textbf{Attention} : les options choisies pour un label s'appliquent aussi aux labels suivants si elles sont inchangées.
        \end{itemize}
  \end{itemize}

\subsection{Solides de base (sans facette)}

Les quatre méthodes décrites ci-dessous sont sensibles à la variable globale \varglob{Hiddenlines} qui vaut \false par défaut.

\subsubsection{Cylindre : Dcylinder}

Dessiner un cylindre à base circulaire (droit ou penché). Plusieurs syntaxes possibles :
\begin{itemize}
    \item Ancienne syntaxe : \cmd{g:Dcylinder(A, V, r \fac{, options})} dessine un cylindre droit, où \argu{A} est un point 3D représentant le centre d'une des faces circulaires, \argu{V} est un point 3D, c'est un vecteur représentant l'axe du cône, le centre de la face circulaire opposée est le point $A+V$ (cette face est orthogonale à \argu{V}), et \argu{r} est le rayon de la base circulaire.
    
    \item La syntaxe : \cmd{g:Dcylinder(A, r, B \fac{, optons})} dessine un cylindre droit, où \argu{A} est un point 3D représentant le centre d'une des faces circulaires, \argu{B} est le centre de la face opposée, et \argu{r} est le rayon. Le cylindre est droit, c'est à dire que les faces circulaires sont orthogonales à l'axe $(AB)$.
    
    \item Pour un cylindre penché :  \cmd{g:Dcylinder(A, r, V, B \fac{, options})}, où \argu{A} est un point 3D représentant le centre d'une des faces circulaires, \argu{B} est le centre de la face circulaire opposée, \argu{r} est le rayon, et \argu{V} est un vecteur 3D non nul orthogonal au plan des faces circulaires.
\end{itemize}

Pour les trois syntaxes, \argu{options} est une table dont les champs définissent les options. Ces options sont (avec leur valeur par défaut) :
\begin{itemize}
    \item \opt{mode=ld.mWireframe} : deux valeurs possibles \val{ld.mWireframe} ou \val{ld.mGrid}. En mode \val{ld.mWireframe} c'est un dessin en fil de fer, en mode \val{ld.mGrid} c'est un dessin en grille (comme s'il y avait des facettes).
    
    \item \opt{hiddenstyle=ld.Hiddenlinestyle} : définit le style de ligne pour les parties cachées (mettre \val{"noline"} pour ne pas les afficher). Par défaut cette option a la valeur de la variable globale \varglob{ld.Hiddenlinestyle} qui est elle même initialisée avec la valeur \val{"dotted"}.
    
    \item \opt{hiddencolor=edgecolor}: définit la couleur des lignes cachées, par défaut c'est la même valeur que l'option \opt{edgecolor}.
    
    \item \opt{edgecolor=<couleur courante>}, définit la couleur des lignes.
    
    \item \opt{edgestyle=<style courant>}, définit le style de ligne pour les arêtes visibles.

    \item \opt{edgewidth=<épaisseur courante>}, définit l'épaisseur des des arêtes visible en dixième de point.    
    
    \item \opt{color=""}, lorsque cette option est une chaîne vide (valeur par défaut) il n'y a pas de remplissage,  lorsque c'est une couleur (sous forme de chaîne) il y a un remplissage avec un gradient linéaire.
    
    \item \opt{opacity=1}, définit la transparence du dessin.
    
    \item Paramètres pour le gradient : lorsqu'il y a un remplissage avec la couleur, un gradient linéaire est utilisé pour le côté du cylindre et un autre pour la section, dans les deux cas le gradient est défini comme suit:\par
    \codeln{left color=<color>!<l>, right color=<color>!<r>, middlecolor=<color>!<m>}
    où \emph{<color>} désigne la couleur utlisée, et \emph{<l>}, \emph{<m>}, \emph{<r>} sont trois nombres entre $0$ et $100$, ce sont ces trois nombres qui constituent les paramètres du gradient sous forme d'une table \emph{\{<l>,<m>,<r>\}}:
        \begin{itemize}
            \item Pour le côté du cylindre, c'est l'option \opt{gradside}, par défaut : \opt{gradside=\{50,10,100\}}.
            \item Pour la section du cylindre, c'est l'option \opt{gradsection}, par défaut : \opt{gradsection=\{25,18,50\}}.
        \end{itemize}
\end{itemize}

\subsubsection{Cône : Dcone}

Dessiner un cône à base circulaire (droit ou penché). Plusieurs syntaxes possibles :
\begin{itemize}
    \item Ancienne syntaxe : \cmd{g:Dcone(A, V, r \fac{, options})} dessine un cône droit, où \argu{A} est un point 3D représentant le sommet du cône, \argu{V} est un point 3D, c'est un vecteur représentant l'axe du cône, le centre de la face circulaire est le point $A+V$ (cette face est orthogonale à \argu{V}), et \argu{r} est le rayon de la base circulaire.
    
    \item La syntaxe : \cmd{g:Dcone(C, r, A \fac{, options})} dessine un cône droit, où \argu{A} est un point 3D représentant le sommet du cône, \argu{C} est le centre de la face circulaire, et \argu{r} est le rayon. Le cône est droit, c'est à dire que la face circulaire est orthogonale à l'axe $(AC)$.
    
    \item Pour un cône penché : \cmd{g:Dcone(C, r, V, A \fac{, options})}, où \argu{A} est un point 3D représentant le sommet du cône, \argu{C} est le centre de la face circulaire, \argu{r} est le rayon, et \argu{V} est un vecteur 3D non nul orthogonal au plan de la face circulaire.
\end{itemize}

Pour les trois syntaxes, \argu{options} est une table dont les champs définissent les options. Ces options sont (avec leur valeur par défaut) :
\begin{itemize}
    \item \opt{mode=ld.mWireframe} : deux valeurs possibles \val{ld.mWireframe} ou \val{ld.mGrid}. En mode \val{ld.mWireframe} c'est un dessin en fil de fer, en mode \val{ld.mGrid} c'est un dessin en grille (comme s'il y avait des facettes).
    
    \item \opt{hiddenstyle=ld.Hiddenlinestyle} : définit le style de ligne pour les parties cachées (mettre \val{"noline"} pour ne pas les afficher). Par défaut cette option a la valeur de la variable globale \varglob{ld.Hiddenlinestyle} qui est elle même initialisée avec la valeur \val{"dotted"}.
    
    \item \opt{hiddencolor=edgecolor}: définit la couleur des lignes cachées, par défaut c'est la même valeur que l'option \opt{edgecolor}.
    
    \item \opt{edgecolor=<couleur courante>}, définit la couleur des lignes.
    
    \item \opt{edgestyle=<style courant>}, définit le style de ligne pour les arêtes visibles.

    \item \opt{edgewidth=<épaisseur courante>}, définit l'épaisseur des des arêtes visible en dixième de point.    
    
    \item \opt{color=""}, lorsque cette option est une chaîne vide (valeur par défaut) il n'y a pas de remplissage,  lorsque c'est une couleur (sous forme de chaîne) il y a un remplissage avec un gradient linéaire.
    
    \item \opt{opacity=1}, définit la transparence du dessin.
    
    \item Paramètres pour le gradient : lorsqu'il y a un remplissage avec la couleur, un gradient linéaire est utilisé pour le côté du cône et un autre pour la section, dans les deux cas le gradient est défini comme suit:\par
    \codeln{left color=<color>!<l>, right color=<color>!<r>, middlecolor=<color>!<m>}
    où \emph{<color>} désigne la couleur utlisée, et \emph{<l>}, \emph{<m>}, \emph{<r>} sont trois nombres entre $0$ et $100$, ce sont ces trois nombres qui constituent les paramètres du gradient sous forme d'une table \emph{\{<l>,<m>,<r>\}}:
        \begin{itemize}
            \item Pour le côté du cône, c'est l'option \opt{gradside}, par défaut : \opt{gradside=\{50,10,100\}}.
            \item Pour la section du cône, c'est l'option \opt{gradsection}, par défaut : \opt{gradsection=\{25,18,50\}}.
        \end{itemize}
\end{itemize}

\subsubsection{Tronc de cône : Dfrustum}

Dessiner un tronc de cône à base circulaire (droit ou penché). Deux syntaxes possibles :

\begin{itemize}
    \item La syntaxe : \cmd{g:Dfrustum(A, R, r, V \fac{, options})} pour un tronc de cône droit, \argu{A} est un point 3D représentant le centre de la face de rayon \argu{R}, \argu{V} est un vecteur 3D représentant l'axe du tronc de cône, le centre de la deuxième face circulaire est le point $A+V$, et son rayon est \argu{r},  (les faces sont orthogonales à \argu{V}). Lorsque $R=r$ on a simplement un cylindre.
    
    \item La syntaxe : \cmd{g:Dfrustum(A, R, r, V, B \fac{, options})} pour un tronc de cône penché, \argu{A} est un point 3D représentant le centre de la face de rayon \argu{R}, \argu{V} est un vecteur 3D représentant un vecteur normal aux faces circulaires, le centre de la deuxième face circulaire est le point \argu{B}, et son rayon est \argu{r}. Lorsque $R=r$ on a un cylindre penché.
\end{itemize}

Dans les deux cas, \argu{options} est une table dont les champs définissent les options. Ces options sont (avec leur valeur par défaut) :
\begin{itemize}
    \item \opt{mode=ld.mWireframe} : deux valeurs possibles \val{ld.mWireframe} ou \val{ld.mGrid}. En mode \val{ld.mWireframe} c'est un dessin en fil de fer, en mode \val{ld.mGrid} c'est un dessin en grille (comme s'il y avait des facettes).
    
    \item \opt{hiddenstyle=ld.Hiddenlinestyle} : définit le style de ligne pour les parties cachées (mettre \val{"noline"} pour ne pas les afficher). Par défaut cette option a la valeur de la variable globale \varglob{ld.Hiddenlinestyle} qui est elle même initialisée avec la valeur \val{"dotted"}.
    
    \item \opt{hiddencolor=edgecolor}: définit la couleur des lignes cachées, par défaut c'est la même valeur que l'option \opt{edgecolor}.
    
    \item \opt{edgecolor=<couleur courante>}, définit la couleur des lignes.
    
    \item \opt{edgestyle=<style courant>}, définit le style de ligne pour les arêtes visibles.

    \item \opt{edgewidth=<épaisseur courante>}, définit l'épaisseur des des arêtes visible en dixième de point.
    
    \item \opt{color=""}, lorsque cette option est une chaîne vide (valeur par défaut) il n'y a pas de remplissage,  lorsque c'est une couleur (sous forme de chaîne) il y a un remplissage avec un gradient linéaire.
    
    \item \opt{opacity=1}, définit la transparence du dessin.
    
    \item Paramètres pour le gradient : lorsqu'il y a un remplissage avec la couleur, un gradient linéaire est utilisé pour le côté du cône et un autre pour la section, dans les deux cas le gradient est défini comme suit:\par
    \codeln{left color=<color>!<l>, right color=<color>!<r>, middlecolor=<color>!<m>}
    où \emph{<color>} désigne la couleur utlisée, et \emph{<l>}, \emph{<m>}, \emph{<r>} sont trois nombres entre $0$ et $100$, ce sont ces trois nombres qui constituent les paramètres du gradient sous forme d'une table \emph{\{<l>,<m>,<r>\}}:
        \begin{itemize}
            \item Pour le côté du cône, c'est l'option \opt{gradside}, par défaut : \opt{gradside=\{50,10,100\}}.
            \item Pour la section du cône, c'est l'option \opt{gradsection}, par défaut : \opt{gradsection=\{25,18,50\}}.
        \end{itemize}
\end{itemize}

\subsubsection{ Sphère : Dsphere}

La méthode \cmd{g:Dsphere(A, r \fac{, options})} dessine une sphère.

\begin{itemize}
    \item \argu{A} est un point 3D représentant le centre de la sphère.
    \item \argu{r} est le rayon de la pshère
    \item \argu{options} est une table dont les champs définissent les options. Ces options sont (avec leur valeur par défaut) :
        \begin{itemize}
            \item \opt{mode=ld.mWireframe} : trois valeurs possibles \val{ld.mWireframe} ou \val{ld.mGrid} ou \val{ld.mBorder}. En mode \val{ld.mWireframe} on dessine le contour (cercle) et l'équateur, en mode \val{ld.mGrid} c'est le contour avec méridiens et fuseaux (grille), et en mode \val{ld.mBorder} c'est le contour uniquement.
            
            \item \opt{hiddenstyle=ld.Hiddenlinestyle} : définit le style de ligne pour les parties cachées (mettre \val{"noline"} pour ne pas les afficher). Par défaut cette option a la valeur de la variable globale \varglob{ld.Hiddenlinestyle} qui est elle même initialisée avec la valeur \val{"dotted"}.
            
            \item \opt{hiddencolor=edgecolor}: définit la couleur des lignes cachées, par défaut c'est la même valeur que l'option \opt{edgecolor}.
            
            \item \opt{edgecolor=<couleur courante>}, définit la couleur des lignes.

            \item \opt{edgestyle=<style courant>}, définit le style de ligne pour les arêtes visibles.

            \item \opt{edgewidth=<épaisseur courante>}, définit l'épaisseur des des arêtes visible en dixième de point.            
    
            \item \opt{color=""}: lorsque cette option est une chaîne vide (valeur par défaut) il n'y a pas de remplissage, lorsque c'est une couleur (sous forme de chaîne) il y a un remplissage avec "ball color".
            
            \item \opt{opacity=1}, définit la transparence du dessin.
        \end{itemize}
\end{itemize}

\begin{demo}{Cylindres, cônes et sphères}
\begin{luadraw}{name=cylindre_cone_sphere}
local ld = luadraw
local pt3d = ld.pt3d
local vecI, vecJ, vecK, M = pt3d.vecI, pt3d.vecJ, pt3d.vecK, pt3d.M

local g = ld.graph3d:new{ size={10,10} }
local dessin = function(args)
    g:Dsphere(M(-1,-2.5,1),2.5, args)
    g:Dcone(M(-1,2.5,5),-5*vecK,2, args)
    g:Dcylinder(M(3,-2,0),6*vecJ,1.5, args)
end
-- en haut à gauche, options par défaut
g:Saveattr(); g:Viewport(-5,0,0,5); g:Coordsystem(-5,5,-5,5,true); dessin(); g:Restoreattr()
g:Saveattr(); g:Viewport(0,5,0,5); g:Coordsystem(-5,5,-5,5,true) -- en haut à droite
dessin({mode=ld.mGrid, hiddenstyle="solid", hiddencolor="LightGray"}); g:Restoreattr()
g:Saveattr(); g:Viewport(-5,0,-5,0); g:Coordsystem(-5,5,-5,5,true) -- en bas à gauche
dessin({mode=ld.mBorder, color="orange"}); g:Restoreattr()
g:Saveattr(); g:Viewport(0,5,-5,0); g:Coordsystem(-5,5,-5,5,true) -- en bas à droite
dessin({mode=ld.mGrid,opacity=0.8,hiddenstyle="noline",color="LightBlue"}); g:Restoreattr()
g:Show()
\end{luadraw}
\end{demo}
%
\section{Solides à facettes}

\subsection{Définition d'un solide}

Il y a deux façons de définir un solide :
\begin{enumerate}
    \item Sous forme d'une liste (table) de facettes. Une facette est elle-même une liste de points 3D (au moins 3) coplanaires et non alignés, qui sont les sommets. Les facettes sont supposées convexes et elles sont orientées par l'ordre d'apparition des sommets. C'est à dire, si $A$, $B$ et $C$ sont les trois premiers sommets d'une facette $F$, alors la facette est orientée avec le vecteur normal $\vec{AB}\wedge\vec{AC}$. Si ce vecteur normal est dirigé vers l'observateur, alors la facette est considérée comme visible. La méthode \cmd{g:Cosine\_incidence(N \fac{, A})} renvoie le cosinus de l'angle au point \argu{A} entre le vecteur \argu{N} (qui doit être unitaire) et le vecteur unitaire dirigé vers l'observateur, si ce cosinus est positif alors \argu{N} est dirigé vers l'observateur (le point \argu{A} est facultatif SAUF en projection centrale).
    
    Dans la définition d'un solide, les vecteurs normaux aux facettes doivent être dirigés vers \textbf{l'extérieur} du solide pour que l'orientation soit correcte.
    
    \item Sous forme de \textbf{polyèdre}, c'est à dire une table à deux champs, un premier champ appelé \emph{vertices} qui est la liste des sommets du polyèdre (points 3D), et un deuxième champ appelé \emph{facets} qui la liste des facettes, mais ici, dans la définition des facettes, les sommets sont remplacés par leur indice dans la liste \emph{vertices}. Les facettes sont orientées de la même façon que précédemment. Cette définition correspond au format \emph{obj} mais sans vecteurs normaux.
\end{enumerate}

Par exemple, considérons les quatre points $A=M(-2,-2,0)$, $B=M(3,0,0)$, $C=M(-2,2,0)$ et $D=M(0,0,4)$, alors on peut définir le tétraèdre construit sur ces quatre points :
\begin{itemize}
    \item soit sous forme d'une liste de facettes : \code{T=\{\{A,B,D\},\{B,C,D\},\{C,A,D\},\{A,C,B\}\}} (attention à l'orientation),
    \item soit sous forme de polyèdre : 
    \code{T=\{vertices=\{A,B,C,D\},facets=\{\{1,2,4\},\{2,3,4\},\{3,1,4\},\{1,3,2\}\}\}}.
\end{itemize}

\paragraph{Fonctions de conversion entre les deux définitions}
\begin{itemize}
    \item La fonction \cmd{ld.poly2facet(P)} où \argu{P} est un polyèdre, renvoie ce solide sous forme d'une liste de facettes.
    \item La fonction \cmd{ld.facet2poly(L \fac{, epsilon})} renvoie la liste de facettes \argu{L} sous forme de polyèdre. L'argument facultatif \argu{epsilon} vaut $10^{-8}$ par défaut, il précise à combien près sont faites les comparaisons entre points 3D. S'il y a beaucoup de facettes, le temps de calcul peut devenir important.
\end{itemize}

\subsection{Dessin d'un polyèdre : Dpoly}

La fonction \cmd{g:Dpoly(P, options)} permet de représenter le polyèdre \argu{P} (par l'algorithme naïf du peintre). L'argument \argu{options} est une table contenant les options, celles-ci sont (avec leur valeur par défaut) :
\begin{itemize}
    \item \opt{mode=ld.mShaded} : définit le mode de représentation. Il y a six valeurs possibles :
        \begin{itemize}
            \item \val{ld.mWireframe} : mode fil de fer, on dessine les arêtes visibles et cachées.
            \item \val{ld.mFlat} : on dessine les faces de couleur unie, ainsi que les arêtes visibles.
            \item \val{ld.mFlatHidden} : on dessine les faces de couleur unie, les arêtes visibles, et les arêtes cachées.
            \item \val{ld.mShaded} : on dessine les faces de couleur nuancée en fonction de leur inclinaison, ainsi que les arêtes visibles. C'est le mode par défaut.
            \item \val{ld.mShadedHidden} : on dessine les faces de couleur nuancée en fonction de leur inclinaison, les arêtes visibles et cachées.
            \item \val{ld.mShadedOnly} :  on dessine les faces de couleur nuancée en fonction de leur inclinaison, mais pas les arêtes.
        \end{itemize}
        \item \opt{contrast=1} : nombre qui permet d'accentuer ou diminuer la nuance des couleurs des facettes dans les modes \val{ld.mShaded}, \val{ld.mShadedHidden}, \val{ld.mShadedOnly}.
        \item \opt{edgestyle=<style courant>} : chaîne qui définit le style de ligne des arêtes.
        \item \opt{edgecolor=<couleur courante>} : chaîne qui définit la couleur des arêtes.
        \item \opt{hiddenstyle=ld.Hiddenlinestyle} : chaîne qui définit le style de ligne des arêtes cachées. Par défaut c'est la valeur contenue dans la variable globale \varglob{ld.Hiddenlinestyle} (qui vaut elle-même \val{"dotted"} par défaut).
        \item \opt{hiddencolor=<couleur courante>} : chaîne qui définit la couleur des arêtes cachées. 
        \item \opt{edgewidth=<épaisseur courante>} : épaisseur de trait des arêtes en dixième de point. 
        \item \opt{opacity=1} : nombre entre 0 et 1 qui permet de mettre une transparence ou non sur les facettes. La valeur par défaut est 1, ce qui signifie pas de transparence.
        \item \opt{backcull=false} : avec la valeur \true, les facettes considérées comme non visibles (vecteur normal non dirigé vers l'observateur) ne sont pas affichées. Cette option est intéressante pour les polyèdres convexes car elle permet de diminuer le nombre de facettes à dessiner.
        \item \opt{twoside=true} : avec la valeur \true on distingue les deux côtés des facettes (intérieur et extérieur), les deux côtés n'auront pas exactement la même couleur.
        \item \opt{color="white"} : chaîne définissant la couleur de remplissage des facettes.
        \item \opt{usepalette=nil} : cette option permet éventuellement de préciser une palette de couleurs pour peindre les facettes ainsi qu'un mode de calcul, la syntaxe est : \opt{usepalette=\{palette,mode\}}, où \argu{palette} désigne une table de couleurs qui sont elles-mêmes des tables de la forme $\{r,g,b\}$ où $r$, $g$ et $b$ sont des nombres entre $0$ et $1$. L'argument  \argu{mode} peut être :
            \begin{itemize}
                \item soit une des chaînes : \val{"x"}, \val{"y"}, \val{"z"}. Dans le premier cas par exemple, les facettes au centre de gravité d'abscisse minimale ont la première couleur de la palette, les facettes au centre de gravité d'abscisse maximale ont la dernière couleur de la palette, pour les autres, la couleur est calculée en fonction de l'abscisse du centre de gravité par interpolation linéaire.
                \item soit une fonction : \argu{mode}$\colon f \mapsto \mathrm{mode}(f)\in\mathbb R$, où $f$ désigne une facette (liste de points 3D). Les facettes ayant la valeur minimale ont la première couleur de la palette, celles ayant la valeur minimale ont la dernière couleur de la palette, pur les autres la couleur est calculée par interpolation linéaire.
            \end{itemize}
\end{itemize}

\begin{demo}{Section d'un tétraèdre par un plan}
\begin{luadraw}{name=tetra_coupe}
local ld = luadraw
local pt3d = ld.pt3d
local Origin, vecI, vecJ, vecK, M = pt3d.Origin, pt3d.vecI, pt3d.vecJ, pt3d.vecK, pt3d.M

local g = ld.graph3d:new{viewdir={10,60},bbox=false, size={10,10}, bg="gray!30"}
local A,B,C,D = M(-2,-4,-2),M(4,0,-2),M(-2,4,-2),M(0,0,2)
local T = ld.tetra(A,B-A,C-A,D-A) -- tétraèdre de sommets A, B, C, D
local plan = {Origin, -vecK} -- plan de coupe
local T1, T2, section = ld.cutpoly(T,plan) -- on coupe le tétraèdre
-- T1 est le polyèdre résultant dans le demi espace contenant -vecK
-- T2 est le polyèdre résultant dans l'autre demi espace
-- section est une facette (c'est la coupe)
g:Dpoly(T1,{color="Crimson", edgecolor="white", opacity=0.8, edgewidth=8})
g:Filloptions("bdiag","Navy"); g:Dpolyline3d(section,true,"draw=none")
g:Dpoly(ld.shift3d(T2,2*vecK), {color="Crimson", edgecolor="white", opacity=0.8, edgewidth=8})
g:Dballdots3d({A,B,C,D+2*vecK}) -- on a dessiné T2 translaté avec le vecteur 2*vecK
g:Show()
\end{luadraw}
\end{demo}

\subsection{Visualiser les numéros des facettes et/ou ceux des sommets d'un polyèdre} 

La méthode \cmd{g:Dpolynames(P \fac{, option, opacity})} affiche le polyèdre \argu{P} avec la possibilité d'ajouter sur chaque face son numéro (qui est son rang d'apparition dans la liste \emph{P.facets}) précédé de la lettre \emph{F}, et éventuellement le numéro de chaque sommet (qui est son rang d'apparition dans la liste \emph{P.vertices}) précédé de la lettre \emph{V}. L'argument \argu{option} peut prendre les valeurs: \val{"facet"} ou bien \val{"vertex"} ou bien \val{"both"} (qui est la valeur par défaut). L'argument \argu{opacity} est un nombre entre $0$ et $1$ qui vaut $0.6$ par défaut.

\begin{demo}{Visualiser les faces et sommets d'un polyèdre}
\begin{luadraw}{name=show_facet_number}
local ld = luadraw
local pt3d = ld.pt3d
local Origin, vecI, vecJ, vecK, M = pt3d.Origin, pt3d.vecI, pt3d.vecJ, pt3d.vecK, pt3d.M

local g = ld.graph3d:new{viewdir={30,60}, window={-2.5,5,-3,3},size={10,10}}
P = ld.parallelep(Origin, 4*vecI,5*vecJ,3*vecK)
local A, B, C = M(4,2.5,3), M(2,5,3), M(4,5,1.5)
P = ld.cutpoly(P, ld.plane(A,B,C), true) -- P est coupé avec le plan, la section est ajoutée en tant que facette
g:Dpolynames(P) -- on visualise les numéros des facettes et des sommets de P
g:Show()
\end{luadraw}
\end{demo}


\subsection{Fonctions de construction de polyèdres}

Les fonctions suivantes renvoient un polyèdre, c'est à dire une table à deux champs, un premier champ appelé \emph{vertices} qui est la liste des sommets du polyèdre (points 3D), et un deuxième champ appelé \emph{facets} qui la liste des facettes, mais dans la définition des facettes, les sommets sont remplacés par leur indice dans la liste \emph{vertices}.

\begin{itemize}
    \item \cmd{ld.tetra(S, v1, v2, v3)} renvoie le tétraèdre  construit à partir du sommet \argu{S} (point 3D) et des 3 vecteurs \argu{v1}, \argu{v2}, \argu{v3} (points 3D) supposés dans le sens direct. Les sommets de ce tétraèdre sont $S$, $S+v_1$, $S+v_2$, $S+v_3$.
    
    \item \cmd{ld.parallelep(A, v1, v2, v3)} renvoie le parallélépipède construit à partir du sommet \argu{A} (point 3D) et de 3 vecteurs \argu{v1}, \argu{v2}, \argu{v3} (points 3D) supposés dans le sens direct.
    
    \item \cmd{ld.prism(base, vector \fac{, open})} renvoie un prisme, l'argument \argu{base} est une liste de points 3D (une des deux bases du prisme), \argu{vector} est le vecteur de translation (point 3D) qui permet d'obtenir la seconde base. L'argument facultatif \argu{open} est un booléen indiquant si le prisme est ouvert ou non (\false par défaut). Dans le cas où il est ouvert, seules les facettes latérales sont renvoyées. La \argu{base} doit être orientée par le \argu{vector}.
    
    \item \cmd{ld.pyramid(base, vertex \fac{, open})} renvoie une pyramide, l'argument \argu{base} est une liste de points 3D, \argu{vertex} est le sommet de la pyramide (point 3D). L'argument facultatif \argu{open} est un booléen indiquant si la pyramide est ouverte ou non (\false par défaut). Dans le cas où elle est ouverte, seules les facettes latérales sont renvoyées. La \argu{base} doit être orientée par le sommet.
    
    \item \cmd{ld.regular\_pyramid(n, side, height \fac{, open, center, axe})} renvoie une pyramide régulière, \argu{n} est le nombre de côtés de la base, l'argument \argu{side} est la longueur d'un côté, et \argu{height} est la hauteur de la pyramide. L'argument facultatif \argu{open} est un booléen indiquant si la pyramide est ouverte ou non (\false par défaut). Dans le cas où elle est ouverte, seules les facettes latérales sont renvoyées. L'argument facultatif \argu{center} est le centre de la base (\varglob{Origin} par défaut), et l'argument facultatif \argu{axe} est un vecteur directeur de l'axe de la pyramide (\varglob{vecK} par défaut).
    
    \item \cmd{ld.truncated\_pyramid(base, vertex, height \fac{, open})} renvoie une pyramide tronquée, l'argument \argu{base} est une liste de points 3D, \argu{vertex} est le sommet de la pyramide (point 3D). L'argument \argu{height} est un nombre indiquant la hauteur par rapport à la base, où s'effectue la troncature, celle-ci est parallèle au plan de la base. L'argument facultatif \argu{open} est un booléen indiquant si la pyramide est ouverte ou non (\false par défaut). Dans le cas où elle est ouverte, seules les facettes latérales sont renvoyées. La \argu{base} doit être orientée par le sommet.
    
    \item \cmd{ld.cylinder(A, V, R \fac{, nbfacet, open})} renvoie un cylindre droit de rayon \argu{R}, \argu{A} (point 3D) est le centre d'une des bases circulaires et \argu{V} est vecteur 3D non nul tel que le centre de la seconde base est le point $A+V$. L'argument facultatif \argu{nbfacet} vaut $35$ par défaut (nombre de facettes latérales). L'argument facultatif \argu{open} est un booléen indiquant si le cylindre est ouvert ou non (\false par défaut). Dans le cas où il est ouvert, seules les facettes latérales sont renvoyées.
    
    \item \cmd{ld.cylinder(A, R, B \fac{, nbfacet, open})} renvoie un cylindre droit de rayon \argu{R}, \argu{A} (point 3D) est le centre d'une des bases circulaires et \argu{B} est le centre de la seconde base. L'argument facultatif \argu{nbfacet} vaut $35$ par défaut (nombre de facettes latérales). L'argument facultatif \argu{open} est un booléen indiquant si le cylindre est ouvert ou non (\false par défaut). Dans le cas où il est ouvert, seules les facettes latérales sont renvoyées.
    
   \item \cmd{ld.cylinder(A, R, V, B \fac{, nbfacet, open})} renvoie un cylindre de rayon \argu{R}, \argu{A} (point 3D) est le centre d'une des bases circulaires et \argu{B} est le centre de la seconde base, et \argu{V} est un vecteur 3D normal au plan des bases circulaires (le cylindre peut donc être penché). L'argument facultatif \argu{nbfacet} vaut $35$ par défaut (nombre de facettes latérales). L'argument facultatif \argu{open} est un booléen indiquant si le cylindre est ouvert ou non (\false par défaut). Dans le cas où il est ouvert, seules les facettes latérales sont renvoyées.
    
    \item \cmd{ld.cone(A, V, R \fac{, nbfacet, open})} renvoie un cône de sommet \argu{A} (point 3D), d'axe dirigé par \argu{V} (point 3D), de base le cercle de centre le point $A+V$ de rayon \argu{R} (dans un plan orthogonal à \argu{V}). L'argument facultatif \argu{nbfacet} vaut $35$ par défaut (nombre de facettes latérales). L'argument facultatif \argu{open} est un booléen indiquant si le cône est ouvert ou non (\false par défaut). Dans le cas où il est ouvert, seules les facettes latérales sont renvoyées.

    \item \cmd{ld.cone(C, R, A \fac{, nbfacet, open})} renvoie un cône de sommet \argu{A} (point 3D), \argu{C} est le centre de base circulaire et \argu{R} son rayon (dans un plan orthogonal à l'axe $(AC)$). L'argument facultatif \argu{nbfacet} vaut $35$ par défaut (nombre de facettes latérales). L'argument facultatif \argu{open} est un booléen indiquant si le cône est ouvert ou non (\false par défaut). Dans le cas où il est ouvert, seules les facettes latérales sont renvoyées.    
    
    \item \cmd{ld.cone(C, R, V, A \fac{, nbfacet, open})} renvoie un cône de sommet \argu{A} (point 3D), \argu{C} est le centre de base circulaire, \argu{R} son rayon, la base est dans un plan orthogonal à \argu{V} (vecteur 3D). L'axe $(AC)$ n'est donc pas forcément orthogonal à la face circulaire (cône penché). L'argument facultatif \argu{nbfacet} vaut $35$ par défaut (nombre de facettes latérales). L'argument facultatif \argu{open} est un booléen indiquant si le cône est ouvert ou non (\false par défaut). Dans le cas où il est ouvert, seules les facettes latérales sont renvoyées.        

    \item \cmd{ld.frustum(C, R, r, V \fac{, nbfacet, open})} renvoie un tronc de cône droit. Le point \argu{C} (point 3D) est le centre de la base circulaire de rayon \argu{R}, le vecteur \argu{V} dirige l'axe du tronc de cône. Le centre de l'autre base circulaire est le point $C+V$, et son rayon est \argu{r} (les bases sont orthogonales à \argu{V}). L'argument facultatif \argu{nbfacet} vaut $35$ par défaut (nombre de facettes latérales). L'argument facultatif \argu{open} est un booléen indiquant si le tronc de cône est ouvert ou non (\false par défaut). Dans le cas où il est ouvert, seules les facettes latérales sont renvoyées.    
    
    \item \cmd{ld.frustum(C, R, r, V, A \fac{, nbfacet, open})} renvoie un tronc de cône. Le point \argu{C} (point 3D) est le centre de la base circulaire de rayon \argu{R}, le centre de l'autre base circulaire est le point \argu{A}, et son rayon est \argu{r}, les bases sont orthogonales au vecteur \argu{V}, mais pas forcément orthogonales à l'axe $(AC)$. L'argument facultatif \argu{nbfacet} vaut $35$ par défaut (nombre de facettes latérales). L'argument facultatif \\argu{C}{open} est un booléen indiquant si le tronc de cône est ouvert ou non (\false par défaut). Dans le cas où il est ouvert, seules les facettes latérales sont renvoyées.

    \item \cmd{ld.sphere(A, R \fac{, nbu, nbv})} renvoie la sphère de centre \argu{A} (point 3D) et de rayon \argu{R}. L'argument facultatif \argu{nbu} représente le nombre de fuseaux ($36$ par défaut) et l'argument facultatif \argu{nbv} le nombre de parallèles ($20$ par défaut).
\end{itemize}

\begin{demo}{Cône tronqué, pyramide tronquée, cylindre oblique}
\begin{luadraw}{name=frustum_pyramid}
local ld = luadraw
local pt3d = ld.pt3d
local Origin, vecI, vecJ, vecK, M = pt3d.Origin, pt3d.vecI, pt3d.vecJ, pt3d.vecK, pt3d.M

local g = ld.graph3d:new{adjust2d=true,bbox=false, size={10,10} }
g:Dfrustum(M(-1,-4,0),3,1,5*vecK, {color="cyan"})
g:Dcylinder(M(-4,4,0),2,vecK,M(-4,2,5), {color="orange"})
local base = ld.map(pt3d.toPoint3d, ld.polyreg(0,3,5))
g:Dpoly(ld.truncated_pyramid( ld.shift3d(base,8*vecI-vecJ-2*vecK), M(5,0,5),4), {mode=4,color="Crimson"})
g:Dcone(M(6,7,-2),3,vecK,M(6,8,5),{color="Pink"})
g:Show()
\end{luadraw}
\end{demo}

\paragraph{Remarque} : nous avons déjà des primitives pour dessiner des cylindres, cônes, et sphères sans passer par des facettes. Un des intérêts de donner une définition de ces objets sous forme de polyèdres est que l'on va pouvoir faire certains calculs sur ces objets comme par exemple des sections planes.

\begin{demo}{Hyperbole : intersection cône - plan}
\begin{luadraw}{name=hyperbole}
local ld = luadraw
local pt3d = ld.pt3d
local Origin, vecI, vecJ, vecK, M = pt3d.Origin, pt3d.vecI, pt3d.vecJ, pt3d.vecK, pt3d.M

local g = ld.graph3d:new{window={-8,6,-9,9},bbox=false, 
    viewdir={"central",45,65}, size={10,10}}
ld.Hiddenlinestyle = "dashed"; ld.Hiddenlines = true
local C1 = ld.cone(Origin,4*vecK,3,35,true)
local C2 = ld.cone(Origin, -4*vecK,3,35,true)
local P = {M(1,-1,-2),vecI} -- plan de la section
local I1 = g:Intersection3d(C1,P) -- intersection entre le cône C1 et le plan P
local I2 = g:Intersection3d(C2,P) -- intersection entre le cône C2 et le plan P
-- I1 et I2 sont de type Edges (arêtes)
g:Dcone(Origin,4*vecK,3,{color="orange"}); g:Dcone(Origin,-4*vecK,3,{color="orange"})
g:Lineoptions("solid","Navy",8)
g:Dedges(I1); g:Dedges(I2) -- dessin des arêtes I1 et I2
g:Dplane(P, vecK,14,9)
g:Show()
\end{luadraw}
\end{demo}

Dans cet exemple, les cônes $C_1$ et $C_2$ sont définis sous forme de polyèdres pour déterminer leur intersection avec le plan $P$, mais pas pour les dessiner. La méthode \cmd{g:Intersection3d(C1, P)} renvoie l'intersection du polyèdre $C_1$ avec le plan $P$ sous la forme d'une table à deux champs, un champ nommé \emph{visible} qui contient une ligne polygonale 3D représentant les arêtes (segments) visibles de l'intersection (c'est à dire qui sont sur une facette visible de $C_1$), et un autre champ nommé \emph{hidden} qui contient une ligne polygonale 3D représentant les arêtes cachées de l'intersection (c'est à dire qui sont sur une facette non visible de $C_1$). La méthode \cmd{g:Dedges()} permet de dessiner ce type d'objets.

\begin{demo}{Section de cône avec plusieurs vues}
\begin{luadraw}{name=several_views}
local ld = luadraw
local pt3d = ld.pt3d
local Origin, vecI, vecJ, vecK, M = pt3d.Origin, pt3d.vecI, pt3d.vecJ, pt3d.vecK, pt3d.M

local g = ld.graph3d:new{window3d={-3,3,-3,3,-3,3}, size={10,10}, margin={0,0,0,0}}
g:Labelsize("footnotesize")
local y0, R = 1, 2.5
local C = ld.cone(M(0,0,3),-6*vecK,R,35,true) -- cone ouvert
local P1 = {M(0,0,0),vecK+vecJ} -- 1er plan de coupe
local P2 = {M(0,y0,0),vecJ} -- 2ieme plan de coupe
local I, I2
local dessin = function() -- un dessin par vue
    g:Dboxaxes3d({grid=true,gridcolor="gray",fillcolor="LightGray"})
    I1 = g:Intersection3d(C,P1) -- intersection entre le cône C et les plans P1 et P2
    I2 = g:Intersection3d(C,P2) -- I1 et I2 sont de type Edges
    g:Dpolyline3d( {{M(0,-3,3),M(0,0,3),M(0,0,-3),M(3,0,-3)}, {M(0,0,-3),M(0,3,-3)}},"red,line width=0.4pt" )
    g:Dcone( M(0,0,3),-6*vecK,R, {color="cyan"})
    g:Dedges(I1, {hidden=true,color="Navy", width=8})
    g:Dedges(I2, {hidden=true,color="DarkGreen", width=8})
end
-- en haut à gauche, vue dans l'espace, on ajoute les plans au dessin
g:Saveattr(); g:Viewport(-5,0,0,5); g:Coordsystem(-7,6,-6,5,1); g:Setviewdir("central"); dessin()
g:Dpolyline3d( {M(-3,-3,3),M(3,-3,3),M(3,3,-3),M(-3,3,-3)},"Navy,line width=0.8pt")
g:Dpolyline3d( {M(-3,y0,3),M(3,y0,3),M(3,y0,-3)},"DarkGreen,line width=0.8pt")
g:Dlabel3d( "$P_1$",M(3,-3,3),{pos="SE",dir={-vecI,-vecJ+vecK},node_options="Navy, draw"})
g:Dlabel3d( "$P_2$",M(-3,y0,3),{pos="SW",dir={-vecI,vecK},node_options="DarkGreen,draw"})
g:Restoreattr()
-- en haut à droite, projection sur le plan xOy
g:Saveattr(); g:Viewport(0,5,0,5); g:Coordsystem(-6,6,-6,5,1); g:Setviewdir("xOy"); dessin()
g:Restoreattr()
-- en bas à gauche, projection sur le plan xOz
g:Saveattr(); g:Viewport(-5,0,-5,0); g:Coordsystem(-6,6,-6,5,1); g:Setviewdir("xOz"); dessin()
g:Restoreattr()
-- en bas à droite, projection sur le plan yOz
g:Saveattr(); g:Viewport(0,5,-5,0); g:Coordsystem(-6,6,-6,5,1); g:Setviewdir("yOz"); dessin()
g:Restoreattr()
g:Show()
\end{luadraw}
\end{demo}

\subsection{Lecture dans un fichier obj}

La fonction \cmd{ld.read\_obj\_file(file)}\footnote{Cette fonction est une contribution de Christophe BAL.} permet de lire le contenu du fichier \emph{obj} désigné par la chaîne de caractères \argu{file}. La fonction lit les définitions des sommets (lignes commençant par \verb|v |), et les lignes définissant les facettes (lignes commençant par \verb|f |). Les autres lignes sont ignorées. La fonction renvoie une séquence constituée du polyèdre, suivi d'une liste de quatre réels \code{\{x1,x2,y1,y2,z1,z2\}} représentant la boîte 3D englobante (bounding box) du polyèdre.

\begin{demo}{Masque de Nefertiti}
\begin{luadraw}{name=lecture_obj}
local ld = luadraw
local P,bbox = ld.read_obj_file("obj/nefertiti.obj")
local g = ld.graph3d:new{window3d=bbox,window={-6,5,-7,7},viewdir={"central",35,65,20},
margin={0,0,0,0}, size={10,10}, bg="LightGray"}
g:Dpoly(P, {usepalette={ld.palAutumn,"z"},mode=ld.mShadedOnly})
g:Show()
\end{luadraw}
\end{demo}

\subsection{Surface au format obj}

Deux syntaxes possibles:
\begin{enumerate}
    \item \cmd{ld.obj\_surface(f, u1, u2, v1, v2 \fac{, grid})} renvoie au format \emph{obj} la surface paramétrée par la fonction \argu{f}$\colon(u,v) \mapsto f(u,v)\in \mathbf R^3$. L'intervalle pour le paramètre $u$ est donné par \argu{u1} et \argu{u2}. L'intervalle pour le paramètre $v$ est donné par \argu{v1} et \argu{v2}. Le paramètre facultatif \argu{grid} vaut $\{25,25\}$ par défaut, il définit le nombre de points à calculer pour le paramètre $u$ suivi du nombre de points à calculer pour le paramètre $v$ (les valeurs de $u$ et $v$ sont équiréparties).
    
    \item \cmd{ld.obj\_surface(f, mesh)} avec \argu{mesh}$=\{\{u_1,\ldots,u_n\}, \{v_1,\ldots,v_m\}\}$, renvoie la surface paramétrée par la fonction \argu{f}$\colon(u,v) \mapsto f(u,v)\in \mathbf R^3$. Les valeurs des paramètres $u$ et $v$ sont données par l'argument \argu{mesh}, elles doivent être dans l'ordre strictement croissant, mais elles ne sont pas forcément équiréparties.
\end{enumerate}
Dans les deux cas, le résultat n'est pas une liste de facettes mais une table à trois champs :
\codeln{\{vertices=\{\ldots\}, normals=\{\ldots\}, facets=\{\{\ldots\},\{\ldots\},\ldots\} \}.}
    \begin{itemize}
        \item Les champs \emph{vertices} et \emph{facets} sont identiques au cas des polyèdres : listes des sommets (points 3D) et liste des facettes avec numéros des sommets, sauf qu'ici les facettes sont toutes \textbf{triangulaires}.
        \item Le champ \emph{normals} est une liste de vecteurs 3D unitaires représentant des vecteurs normaux à la surface à raison de un par sommet.
    \end{itemize}



\subsection{Dessin d'une liste de facettes : Dfacet et Dmixfacet}

Il y a deux méthodes possibles :
\begin{enumerate}
    \item Pour un solide $S$ sous forme d'une liste de facettes (avec points 3D), la méthode est :
    cmdln{g:Dfacet(S, options)}
    où \argu{S} est la liste de facettes et \argu{options} une table définissant les options. Celles-ci sont :
    \begin{itemize}
        \item \opt{mode=ld.mShaded} : définit le mode de représentation. Les valeurs possibles sont:
            \begin{itemize}
                \item \val{ld.mWireframe} : mode fil de fer, on dessine les arêtes seulement.
                \item \val{ld.mFlat} ou \val{ld.mFlatHidden} : on dessine les faces de couleur unie, ainsi que les arêtes.
                \item \val{ld.mShaded} ou \val{ld.mShadedHidden} : on dessine les faces de couleur nuancée en fonction de leur inclinaison, ainsi que les arêtes.
                \item \val{ld.mShadedOnly} : on dessine les faces de couleur nuancée en fonction de leur inclinaison, mais pas les arêtes.
            \end{itemize}
        \item \opt{contrast=1} : ce nombre permet d'accentuer ou diminuer la nuance des couleurs des facettes dans les modes \val{ld.mShaded}, \val{ld.mShadedHidden}, \val{ld.mShadedOnly}.
        \item \opt{edgestyle=<style courant>} : chaîne qui définit le style de ligne des arêtes.
        \item \opt{edgecolor=<couleur courante>} : chaîne qui définit la couleur des arêtes.
        \item \opt{hiddenstyle=ld.Hiddenlinestyle} : chaîne qui définit le style de ligne des arêtes cachées. Par défaut c'est la valeur contenue dans la variable globale \varglob{ld.Hiddenlinestyle} (qui vaut elle-même \val{"dotted"} par défaut).
        \item \opt{hiddencolor=<couleur par défaut>} : chaîne qui définit la couleur des arêtes cachées.
        \item \opt{edgewidth=<épaisseur courante>} : épaisseur de trait des arêtes en dixième de point.
        \item \opt{opacity=1} : nombre entre 0 et 1 qui permet de mettre une transparence ou non sur les facettes.
        \item \opt{backcull=false} : avec la valeur \true, les facettes considérées comme non visibles (vecteur normal non dirigé vers l'observateur) ne sont pas affichées. Cette option est intéressante pour les polyèdres convexes car elle permet de diminuer le nombre de facettes à dessiner.
        \item \opt{clip=false} : avec la valeur \true, les facettes sont clippées par la fenêtre 3D.
        \item \opt{twoside} : booléen qui vaut true par défaut, ce qui signifie qu'on distingue les deux côtés des facettes (intérieur et extérieur), les deux côtés n'auront pas exactement la même couleur.
        \item \opt{color="white"} : chaîne définissant la couleur de remplissage des facettes.
        \item \opt{twoside=true} : avec la valeur \true on distingue les deux côtés des facettes (intérieur et extérieur), les deux côtés n'auront pas exactement la même couleur.
        \item \opt{color="white"} : chaîne définissant la couleur de remplissage des facettes.
        \item \opt{usepalette=nil} : cette option permet éventuellement de préciser une palette de couleurs pour peindre les facettes ainsi qu'un mode de calcul, la syntaxe est : \opt{usepalette=\{palette,mode\}}, où \argu{palette} désigne une table de couleurs qui sont elles-mêmes des tables de la forme $\{r,g,b\}$ où $r$, $g$ et $b$ sont des nombres entre $0$ et $1$. L'argument  \argu{mode} peut être :
            \begin{itemize}
                \item soit une des chaînes : \val{"x"}, \val{"y"}, \val{"z"}. Dans le premier cas par exemple, les facettes au centre de gravité d'abscisse minimale ont la première couleur de la palette, les facettes au centre de gravité d'abscisse maximale ont la dernière couleur de la palette, pour les autres, la couleur est calculée en fonction de l'abscisse du centre de gravité par interpolation linéaire.
                \item soit une fonction : \argu{mode}$\colon f \mapsto \mathrm{mode}(f)\in\mathbb R$, où $f$ désigne une facette (liste de points 3D). Les facettes ayant la valeur minimale ont la première couleur de la palette, celles ayant la valeur minimale ont la dernière couleur de la palette, pur les autres la couleur est calculée par interpolation linéaire.
            \end{itemize}    
        \end{itemize}
    
    \item Pour plusieurs listes de facettes dans un même dessin, la méthode est :
        \cmdln{g:Dmixfacet(S1, options1, S2, options2, \ldots)}
    où \argu{S1}, \argu{S2}, \ldots, sont des listes de facettes, et \argu{options1}, \argu{options2}, \ldots, sont les options correspondantes. Les options d'une liste de facettes s'appliquent aussi aux suivantes si elles ne sont pas changées. Ces options sont identiques à la méthode précédente.
    
    Cette méthode est utile pour dessiner plusieurs solides ensemble, à condition qu'il n'y ait pas d'intersections entre les objets, car celles-ci ne sont pas gérées ici.
\end{enumerate}

\begin{demo}[courbeniv]{Exemple de courbes de niveaux sur une surface}
\begin{luadraw}{name=courbes_niv}
local ld = luadraw
local pt3d = ld.pt3d
local Origin, vecI, vecJ, vecK, M, Z = pt3d.Origin, pt3d.vecI, pt3d.vecJ, pt3d.vecK, pt3d.M, ld.cpx.Z
local cos, sin = math.cos, math.sin
local g = ld.graph3d:new{window3d={0,5,0,10,0,11}, adjust2d=true, size={10,10},
    viewdir={"central",220,60,15,M(2.5,5,5.5)}}
g:Labelsize("footnotesize")
local S = ld.cartesian3d(function(u,v) return (u+v)/(2+cos(u)*sin(v)) end,0,5,0,10,{30,30})
local n = 10 -- nombre de niveaux
local Colors = ld.getpalette(ld.palGasFlame,n,true) -- liste de 10 couleurs au format table
local niv, S1 = {}
for k = 1, n do
    S1, S = ld.cutfacet(S,{M(0,0,k),-vecK}) -- section de S avec le plan z=k
    ld.insert(niv,{S1, {color=Colors[k],mode=ld.mShaded,edgewidth=0.5}}) -- S1 est la partie sous le plan et S au dessus
end
ld.insert(niv,{S, {color=Colors[n+1]}}) -- insertion du dernier niveau
-- niv est une liste du type {facettes1, options1, facettes2, options2, ...}
g:Dboxaxes3d({grid=true, gridcolor="gray",fillcolor="lightgray"})
g:Dmixfacet(table.unpack(niv))
for k = 1, n do
    g:Dballdots3d( M(5,0,k), ld.rgb(Colors[k]))
end
g:Dlabel("$z=\\frac{x+y}{2+\\cos(x)\\sin(y)}$", Z((g:Xinf()+g:Xsup())/2, g:Yinf()), {pos="N"})
g:Show()
\end{luadraw}
\end{demo}

\subsection{Fonctions de construction de listes de facettes}

Les fonctions suivantes renvoient un solide sous forme d'une liste de facettes (avec points 3D).

\subsubsection{surface()}

Deux syntaxes possibles :
\begin{enumerate}
    \item \cmd{ld.surface(f, u1, u2, v1, v2 \fac{, grid})} renvoie la surface paramétrée par la fonction \argu{f}$\colon(u,v) \mapsto f(u,v)\in \mathbf R^3$. L'intervalle pour le paramètre $u$ est donné par \argu{u1} et \argu{u2}. L'intervalle pour le paramètre $v$ est donné par \argu{v1} et \argu{v2}. Le paramètre facultatif \argu{grid} vaut $\{25,25\}$ par défaut, il définit le nombre de points à calculer pour le paramètre $u$ suivi du nombre de points à calculer pour le paramètre $v$ (les valeurs de $u$ et $v$ sont équiréparties).
    
    \item \cmd{ld.surface(f, mesh)} avec \argu{mesh}$=\{\{u_1,\ldots,u_n\}, \{v_1,\ldots,v_m\}\}$, renvoie la surface paramétrée par la fonction \argu{f}$\colon(u,v) \mapsto f(u,v)\in \mathbf R^3$. Les valeurs des paramètres $u$ et $v$ sont données par l'argument \argu{mesh}, elles doivent être dans l'ordre strictement croissant, mais elles ne sont pas forcément équiréparties.
\end{enumerate}

    
Il y a deux variantes pour les surfaces :

\subsubsection{cartesian3d()}

La fonction \cmd{ld.cartesian3d(f, x1, x2, y1, y2 \fac{, grid, addwall})} renvoie la surface cartésienne d'équation $z=f(x,y)$ où \argu{f}$\colon(x,y)\mapsto f(x,y)\in\mathbb R$. L'intervalle pour $x$ est donné par \argu{x1} et \argu{x2}. L'intervalle pour $y$ est donné par \argu{y1} et \argu{y2}. Le paramètre facultatif \argu{grid} vaut $\{25,25\}$ par défaut, il définit le nombre de points à calculer pour $x$ suivi du nombre de points à calculer pour $y$. Le paramètre \argu{addwall} vaut \val{0} ou \val{"x"}, ou \val{"y"}, ou \val{"xy"} (\val{0} par défaut). Lorsque cette option vaut \val{"x"} (ou \val{"xy"}), la fonction renvoie, après la liste des facettes composant la surface, une liste de facettes séparatrices (murs ou cloisons) entre chaque "couche" de facettes, une couche correspond à deux valeurs consécutives du paramètre $x$\footnote{Ces cloisons sont en fait des plans d'équation $x=$constante}. Avec la valeur \val{"y"} (ou \val{"xy"}) c'est une liste de facettes séparatrices (murs) entre chaque "couche" correspond à deux valeurs consécutives du paramètre "y"\footnote{Ces cloisons sont en fait des plans d'équation $y=$constante}. Cette option peut être utile avec la méthode \cmd{g:Dscene3d()} (uniquement), car les cloisons séparatrices forment une partition de l'espace isolant les facettes de la surface, ce qui permet d'éviter des calculs d'intersection inutiles entre elles. C'est notamment le cas avec des surfaces non convexes.

Par exemple, voici le code de la figure \ref{pointcol}:
\begin{Luacode}
\begin{luadraw}{name=point_col}
local ld = luadraw
local pt3d = ld.pt3d
local Origin, vecI, vecJ, vecK, M = pt3d.Origin, pt3d.vecI, pt3d.vecJ, pt3d.vecK, pt3d.M

local g = ld.graph3d:new{window3d={-2,2,-2,2,-4,4}, window={-3.5,3,-5,5}, size={8,9,0}, viewdir={120,60}}
local S = ld.cartesian3d(function(u,v) return u^2-v^2 end, -2,2,-2,2,{20,20}) -- surface d'équation z=x^2-y^2
local Tx = g:Intersection3d(S, {Origin,vecI}) --intersection entre S et le plan yOz
local Ty = g:Intersection3d(S, {Origin,vecJ}) --intersection entre S et le plan xOz
g:Dboxaxes3d({grid=true,gridcolor="gray",fillcolor="LightGray",drawbox=true})
g:Dfacet(S,{mode=ld.mShadedOnly,color="ForestGreen"}) -- dessin de la surface
g:Dedges(Tx, {color="Crimson", hidden=true, width=8}) -- dessin de l'intersection avec yOz
g:Dedges(Ty, {color="Navy",hidden=true, width=8}) -- dessin de l'intersection avec xOz
g:Dpolyline3d( {M(2,0,4),M(-2,0,4),M(-2,0,-4)}, "Navy,line width=.8pt")
g:Dpolyline3d( {M(0,-2,4),M(0,2,4),M(0,2,-4)}, "Crimson,line width=.8pt")
g:Show()
\end{luadraw}
\end{Luacode}

\subsubsection{cylindrical\_surface()}

La fonction \cmd{ld.cylindrical\_surface(r, z, u1, u2, theta1, theta2 \fac{, grid, addwall})} renvoie la surface paramétrée en cylindrique par \emph{r(u,theta), theta, z(u,theta)}. Les arguments \argu{r} et \argu{z} sont donc deux fonctions de $u$ et $\theta$, à valeurs réelles. L'intervalle pour $u$ est donné par \argu{u1} et \argu{u2}. L'intervalle pour $\theta$ est donné par \argu{theta1} et \argu{theta2} (en radians). Le paramètre facultatif \argu{grid} vaut $\{25,25\}$ par défaut, il définit le nombre de points à calculer pour $u$ suivi du nombre de points à calculer pour $v$. Le paramètre \argu{addwall} vaut $0$ ou \val{"v"}  ou \val{"z"} ou \val{"vz"} ($0$ par défaut). Lorsque cette option vaut \val{"v"}  ou \val{"vz"} , la fonction renvoie, après la liste des facettes composant la surface, une liste de facettes séparatrices (murs ou cloisons) entre chaque "couche" de facettes, une couche correspond à deux valeurs consécutives de l'angle $\theta$\footnote{Ces cloisons sont en fait des plans d'équation $\theta=$constante}. Lorsque cette option vaut \val{"z"}  ou \val{"vz"}, la fonction renvoie, après la liste des facettes composant la surface, une liste de facettes séparatrices (murs ou cloisons) entre chaque "couche" de facettes, une couche correspond à deux valeurs consécutives de la cote $z$\footnote{Ces cloisons sont en fait des plans d'équation $z=$constante}, les valeurs de $z$ sont calculées à  partir des valeurs du paramètres $u$ et avec la valeur \argu{theta1}, ceci est utile lorsque $z$ ne dépend que $u$ (et donc pas de $\theta$). Cette option peut être utile avec la méthode \cmd{g:Dscene3d()} (uniquement), car les cloisons séparatrices forment une partition de l'espace isolant les facettes de la surface, ce qui permet d'éviter des calculs d'intersection inutiles entre elles. C'est notamment le cas avec des surfaces non convexes.

\begin{demo}{Surfaces utilisant l'option \emph{addwall}}
\begin{luadraw}{name=surface_with_addWall}
local ld = luadraw
local pi, ch, sh = math.pi, math.cosh, math.sinh
local O = ld.pt3d.Origin
local g = ld.graph3d:new{window3d={-4,4,-4,4,-5,5}, window={-10,10,-4,4}, size={10,10}, viewdir={60,60}}
g:Labelsize("footnotesize")
local S,wall = ld.cartesian3d(function(x,y) return x^2-y^2 end,-2,2,-2,2,nil,"xy")
g:Saveattr(); g:Viewport(-10,0,-4,4); g:Coordsystem(-4.5,4.5,-4.5,4.75)
g:Dscene3d(
    g:addWall(wall), -- 2 intersections de facettes avec cette instruction, et 529 sans
    g:addFacet(S,{color="SteelBlue"}),
    g:addAxes(O,{arrows=1}) )
g:Restoreattr()
g:Saveattr(); g:Viewport(0,10,-4,4); g:Coordsystem(-5,5,-5,5)
local r = function(u,v) return ch(u) end
local z = function(u,v) return sh(u) end
S,wall = ld.cylindrical_surface(r,z,-2,2,-pi,pi,{25,51},"zv")
g:Dscene3d(
    g:addWall(wall), -- 13 intersections de facettes avec cette instruction, et plus de 17000 sans ...
    g:addFacet(S,{color="Crimson"}),
    g:addAxes(O,{arrows=1}) )
g:Restoreattr()
g:Show()
\end{luadraw}
\end{demo}

\subsubsection{curve2cone()}
La fonction \cmd{ld.curve2cone(f, t1, t2, S, options)} construit un cône de sommet \argu{S} (point 3D) et de base la courbe paramétrée par \argu{f}$\colon t\mapsto f(t)\in\mathbf R^3$ sur l'intervalle défini par \argu{t1} et \argu{t2}. L'argument \argu{options} est une table dont les champs définissent les options, qui sont (avec leur valeur  par défaut) :
    \begin{itemize}
        \item \opt{nbdots=15} : nombre minimal de points de la courbe à calculer.
        
        \item \opt{ratio=0} :  nombre représentant le rapport d'homothétie (de centre le sommet \argu{S}) pour construire l'autre partie du cône, avec la valeur $0$ il n'y a pas de deuxième partie.
        
        \item \opt{nbdiv=0} : entier positif indiquant le nombre de fois que l'intervalle entre deux valeurs consécutives du paramètre $t$ peut être coupé en deux (dichotomie) lorsque les points correspondants sont trop éloignés.
        
        \item \opt{obj=\false} : avec la valeur \false cette fonction construit une liste de facettes, avec les valeur \true elle construit une table à trois champs : \code{\{vertices=\{points 3D\}, facets=\{\{index1,\ldots\},\ldots\}, normals=\{vecteurs 3D\}\}}.
    \end{itemize}
La fonction renvoie une séquence : 
\begin{enumerate}
    \item la liste de facettes ou bien la table au format \emph{obj}, suivie 
    \item d'une ligne polygonale 3D qui représente les bords du cône.
\end{enumerate}

 
\begin{demo}{Exemple de cône elliptique}
\begin{luadraw}{name=curve2cone}
local ld = luadraw
local O, M = ld.pt3d.Origin, ld.pt3d.M
local cos, sin, pi = math.cos, math.sin, math.pi
local g = ld.graph3d:new{ window3d={-2,2,-4,4,-3,3}, window={-5.5,5.5,-5.5,5},
    size={10,10}, viewdir="central"}
local f = function(t) return M(2*cos(t),4*sin(t),-3) end -- ellipse dans le plan z=-3
local C, bord = ld.curve2cone(f,-pi,pi,O,{nbdiv=2, ratio=-1})
g:Dboxaxes3d({grid=true,gridcolor="gray",fillcolor="LightGray"})
g:Dpolyline3d(bord[1],"red,line width=2.4pt") -- bord inférieur
g:Dfacet(C, {mode=ld.mShadedOnly,color="LightBlue"}) -- cône
g:Dpolyline3d(bord[2],"red,line width=0.8pt") -- bord supérieur
g:Show()
\end{luadraw}
\end{demo}

\subsubsection{curve2cylinder()}
La fonction \cmd{ld.curve2cylinder(f, t1, t2, V \fac{, options})} construit un cylindre d'axe dirigé par le vecteur \argu{V} (point 3D) et de base la courbe paramétrée par \argu{f}$\colon t\mapsto f(t)\in\mathbf R^3$ sur l'intervalle défini par \argu{t1} et \argu{t2}. La seconde base est la translatée de la première avec le vecteur \argu{V}. L'argument \argu{options} est une table dont les champs définissent les options, qui sont (avec leur valeur  par défaut) :
    \begin{itemize}
        \item \opt{nbdots=15} : nombre minimal de points de la courbe à calculer.
        
        \item \opt{nbdiv=0} : entier positif indiquant le nombre de fois que l'intervalle entre deux valeurs consécutives du paramètre $t$ peut être coupé en deux (dichotomie) lorsque les points correspondants sont trop éloignés
        
        \item \opt{obj=\false} : avec la valeur \false cette fonction construit une liste de facettes, avec les valeur \true elle construit une table à trois champs : \code{\{vertices=\{points 3D\}, facets=\{\{index1,\ldots\},\ldots\}, normals=\{vecteurs 3D\}\}}.
    \end{itemize}
La fonction renvoie une séquence : 
\begin{enumerate}
    \item la liste de facettes ou bien la table au format \emph{obj}, suivie 
    \item d'une ligne polygonale 3D qui représente les bords du cylindre.
\end{enumerate}
 
\begin{demo}{Section d'un cylindre non circulaire}
\begin{luadraw}{name=curve2cylinder}
local ld = luadraw
local pt3d = ld.pt3d
local Origin, vecI, vecJ, vecK, M = pt3d.Origin, pt3d.vecI, pt3d.vecJ, pt3d.vecK, pt3d.M
local cos, sin, pi = math.cos, math.sin, math.pi

local g = ld.graph3d:new{ window3d={-5,5,-5,5,-4,4},window={-9,8,-10.5,5.5},
    viewdir={"central",39,64}, size={10,10}}
local f = function(t) return M(4*cos(t)-cos(4*t),4*sin(t)-sin(4*t),-4) end -- courbe dans le plan z=-3
local V = 8*vecK
local C = ld.curve2cylinder(f,-pi,pi,V,{nbdots=25,nbdiv=2})
local plan = {M(0,0,2), -vecK} -- plan de coupe z=2
local C1, C2, section = ld.cutfacet(C,plan)
g:Dboxaxes3d({grid=true,gridcolor="gray",fillcolor="LightGray"})
g:Dfacet(C1, {mode=ld.mShaded,color="LightBlue"}) -- partie sous le plan
g:Dfacet(g:Plane2facet(plan), {opacity=0.3,color="Chocolate"}) -- dessin du plan sous forme d'une facette
g:Filloptions("fdiag","red"); g:Dpolyline3d(section) -- dessin de la section
g:Dfacet(C2, {mode=3,color="LightBlue"}) -- partie du cylindre au dessus du plan
g:Show()
\end{luadraw}
\end{demo}

\subsubsection{line2tube(); section2tube()}

La fonction \cmd{ld.line2tube(L, r \fac{, options})} construit un tube centré sur \argu{L} qui doit être une liste de points 3D, l'argument \argu{r} représente le rayon de ce tube. L'argument \argu{options} est une table dont les champs définissent les options, qui sont (avec leur valeur  par défaut) :
    \begin{itemize}
        \item \opt{nbfacet=3} : nombre de facettes latérales du tube.
        
        \item \opt{close=false} : booléen indiquant si la ligne polygonale doit être refermée.
        
        \item \opt{hollow=false} : booléen indiquant si les deux extrémités du tube doivent être ouvertes ou non. Lorsque l'option \opt{close} vaut \true, l'option \opt{hollow} prend automatiquement la valeur \true.
        
        \item \opt{addwall=0} : nombre qui vaut $0$ ou $1$. Lorsque cette option vaut $1$, la fonction renvoie, après le tube, une liste de facettes séparatrices (murs) entre chaque "tronçon" du tube, ce qui peut être utile avec la méthode \cmd{g:Dscene3d()} (uniquement).
        
        \item \opt{obj=\false} : avec la valeur \false cette fonction renvoie une liste de facettes, avec les valeur \true elle renvoie une table à trois champs : \code{\{vertices=\{points 3D\}, facets=\{\{index1,\ldots\},\ldots\}, normals=\{vecteurs 3D\}\}}.
    \end{itemize}
    
La fonction \cmd{ld.section2tube(section, L \fac{, options})} construit également un tube centré sur \argu{L} qui doit être une liste de points 3D, l'argument \argu{section} doit être une facette centrée sur le premier point de \argu{L}, elle représente une section du tube qui va être construit. L'argument \argu{options} est une table dont les champs définissent les options, qui sont (avec leur valeur  par défaut) :
    \begin{itemize}
        \item \opt{nbfacet=3} : nombre de facettes latérales du tube.
        
        \item \opt{close=false} : booléen indiquant si la ligne polygonale doit être refermée.
        
        \item \opt{hollow=false} : booléen indiquant si les deux extrémités du tube doivent être ouvertes ou non. Lorsque l'option \opt{close} vaut \true, l'option \opt{hollow} prend automatiquement la valeur \true.
        
        \item \opt{addwall=0} : nombre qui vaut $0$ ou $1$. Lorsque cette option vaut $1$, la fonction renvoie, après le tube, une liste de facettes séparatrices (murs) entre chaque "tronçon" du tube, ce qui peut être utile avec la méthode \cmd{g:Dscene3d()} (uniquement).
        
        \item \opt{obj=\false} : avec la valeur \false cette fonction renvoie une liste de facettes, avec les valeur \true elle renvoie une table à trois champs : \code{\{vertices=\{points 3D\}, facets=\{\{index1,\ldots\},\ldots\}, normals=\{vecteurs 3D\}\}}.        
    \end{itemize}
    
\begin{demo}{Exemple avec line2tube et section2tube}
\begin{luadraw}{name=line2tube_section2tube}
local ld = luadraw
local pt3d = ld.pt3d
local Origin, vecI, vecJ, vecK, M = pt3d.Origin, pt3d.vecI, pt3d.vecJ, pt3d.vecK, pt3d.M

local g = ld.graph3d:new{window={-5,6,-4.5,8}, viewdir={45,60}, margin={0,0,0,0}, size={10,10}}
local L1 = ld.map(pt3d.toPoint3d,ld.polyreg(0,3,6)) -- hexagone régulier dans le plan xOy, centre O de sommet M(3,0,0)
local L2 = ld.shift3d(ld.rotate3d(L1,90,{Origin,vecJ}),3*vecJ)
local L3 = ld.shift3d(ld.reverse(L1),6*vecK)
L3[6] = L3[5]-2*vecK -- modification du dernier point
local section = ld.shift3d({M(2,0,0.5),M(4,0,0.5),M(4,0,-0.5),M(2,0,-0.5)},6*vecK)
local T1 = ld.line2tube(L1,1,{nbfacet=8,close=true}) -- tube 1 refermé
local T2 = ld.line2tube(L2,1,{nbfacet=8}) -- tube 2 non refermé
local T3 = ld.section2tube(section, L3,{hollow=true})
g:Dmixfacet( T1, {color="Crimson",opacity=0.8}, T2, {color="SteelBlue"}, T3, {color="ForestGreen"} )
g:Show()
\end{luadraw}
\end{demo}

\subsubsection{rotcurve()}
La fonction \cmd{ld.rotcurve(p, t1, t2, axe, angle1, angle2 \fac{, options})} construit sous forme d'une liste de facettes, la surface balayée par la courbe paramétrée par \argu{p}$\colon t\mapsto p(t)\in \mathbf R^3$ sur l'intervalle défini par \argu{t1} et \argu{t2}, en la faisant tourner autour de \argu{axe} (qui est une table de la forme \{point 3D, vecteur 3D\} représentant une droite orientée de l'espace), d'un angle allant de \argu{angle1} (en degrés) à \argu{angle2}. L'argument \argu{options} est une table dont les champs définissent les options, qui sont (avec leur valeur  par défaut) :
    \begin{itemize}
        \item \opt{grid=\{25,25\}} : table constituée de deux nombres, le premier est le nombre de points calculés pour le paramètre $t$, et le second le nombre de points calculés pour le paramètre angulaire.

        \item \opt{addwall=0} : nombre qui vaut $0$ ou $1$ ou $2$. Lorsque cette option vaut $1$, la fonction renvoie, après la surface, une liste de facettes séparatrices (murs) entre chaque "couche" de facettes (une couche correspond à deux valeurs consécutives du paramètre $t$), et avec la valeur $2$ c'est une liste de facettes séparatrices (murs) entre chaque "tranche" de rotation (une couche correspond à deux valeurs consécutives du paramètre angulaire, ceci est intéressant lorsque la courbe est dans un même plan que l'axe de rotation). Cette option peut être utile avec la méthode \cmd{g:Dscene3d()} (uniquement).
        
        \item \opt{obj=\false} : avec la valeur \false cette fonction renvoie une liste de facettes, avec les valeur \true elle renvoie une table à trois champs : \code{\{vertices=\{points 3D\}, facets=\{\{index1,\ldots\},\ldots\}, normals=\{vecteurs 3D\}\}}.           
    \end{itemize} 
        
\begin{demo}{Exemple avec rotcurve}
\begin{luadraw}{name=rotcurve}
local ld = luadraw
local pt3d = ld.pt3d
local Origin, vecI, vecJ, vecK, M = pt3d.Origin, pt3d.vecI, pt3d.vecJ, pt3d.vecK, pt3d.M
local cos, sin, pi = math.cos, math.sin, math.pi
local g = ld.graph3d:new{viewdir={30,60},size={10,10}}
local p = function(t) return M(0,sin(t)+2,t) end -- courbe dans le plan yOz
local axe = {Origin,vecK}
local S = ld.rotcurve(p,pi,-pi,axe,0,360,{grid={25,35}})
local visible, hidden = g:Classifyfacet(S)
g:Dfacet(hidden, {mode=ld.mShadedOnly,color="cyan"})
g:Dline3d(axe,"red,line width=1.2pt")
g:Dfacet(visible, {mode=5,color="cyan"})
g:Dline3d(axe,"red,line width=1.2pt,dashed")
g:Dparametric3d(p,{t={-pi,pi},draw_options="red,line width=1.2pt"})
g:Show()
\end{luadraw}
\end{demo}

\paragraph{Remarque} : si l'orientation de la surface ne semble pas bonne, il suffit d'échanger les paramètres \argu{t1} et \argu{t2}, ou bien \argu{angle1} et \argu{angle2}.

\subsubsection{rotline()}

La fonction \cmd{ld.rotline(L, axe, angle1, angle2 \fac{, options})} construit sous forme d'une liste de facettes, la surface balayée par la liste de points 3D \argu{L} en la faisant tourner autour de \argu{axe} (qui est une table de la forme \{point 3D, vecteur 3D\} représentant une droite orientée de l'espace), d'un angle allant de \argu{angle1} (en degrés) à \argu{angle2}. L'argument \argu{options} est une table dont les champs définissent les options, qui sont (avec leur valeur  par défaut) :
    \begin{itemize}
        \item \opt{nbdots=25} : nombre de points calculés pour le paramètre angulaire.
        
        \item \opt{close=false} : booléen qui indique si \argu{L} doit être refermée.
        
        \item \opt{addwall=0} : nombre qui vaut $0$ ou $1$ ou $2$. Lorsque cette option vaut $1$, la fonction renvoie, après la liste des facettes composant la surface, une liste de facettes séparatrices (murs) entre chaque "couche" de facettes (une couche correspond à deux valeurs consécutives du paramètre $t$), et avec la valeur $2$ c'est une liste de facettes séparatrices (murs) entre chaque "tranche" de rotation (une couche correspond à deux valeurs consécutives du paramètre angulaire, ceci est intéressant lorsque la courbe est dans un même plan que l'axe de rotation). Cette option peut être utile avec la méthode \cmd{g:Dscene3d()} (uniquement).
        
        \item \opt{obj=\false} : avec la valeur \false cette fonction renvoie une liste de facettes, avec les valeur \true elle renvoie une table à trois champs : \code{\{vertices=\{points 3D\}, facets=\{\{index1,\ldots\},\ldots\}, normals=\{vecteurs 3D\}\}}.           
        \end{itemize} 
        
\begin{demo}{Exemple avec rotline}
\begin{luadraw}{name=rotline}
local ld = luadraw
local pt3d = ld.pt3d
local M = pt3d.M
local g = ld.graph3d:new{window={-4,4,-4,4},size={10,10}}
local L = {M(0,0,4),M(0,4,0),M(0,0,-4)} -- liste de points dans le plan yOz
local axe = {pt3d.Origin, pt3d.vecK}
local S = ld.rotline(L,axe,0,360,{nbdots=5}) -- le point 1 et le point 5 sont confondus
g:Dfacet(S,{color="Crimson",edgecolor="Gold",opacity=0.8})
g:Show()
\end{luadraw}
\end{demo}      


\subsection{Arêtes d'un solide}

Un objet de type "edge" est une table à deux champs, un champ nommé \emph{visible} qui contient une ligne polygonale 3D correspondant aux arêtes visibles, et un autre champ nommé \emph{hidden} qui contient une ligne polygonale 3D correspondant aux arêtes cachées.

\begin{itemize}
    \item La méthode \cmd{g:Edges(P)} où \argu{P} est un polyèdre, renvoie les arêtes de $P$ sous forme d'un objet de type "edge". Une arête de argu est visible lorsqu'elle appartient à au moins une face visible.
    
    \item La méthode \cmd{g:Intersection3d(P, plane)} où argu est un polyèdre ou bien une liste de facettes, renvoie sous forme d'objet de type "edge" l'intersection entre argu et le plan représenté par \argu{plane} (c'est une table de la forme \{A,u\} où $A$ est un point du plan et $u$ un vecteur normal, ce sont donc deux points 3D).
    
    \item La méthode \cmd{g:Dedges(edges, options)} permet de dessiner \argu{edges} qui doit être un objet de type "edge". L'argument \argu{options} est une table dont les champs définissent les options, qui sont (avec leur valeur  par défaut) :
    \begin{itemize}
        \item \opt{hidden=false} : booléen qui indique si les arêtes cachées doivent être dessinées).
        \item \opt{visible=true} : booléen qui indique si les arêtes visibles doivent être dessinées.
        \item \opt{clip=false} : booléen qui indique si les arêtes doivent être clippées par la fenêtre 3D.
        \item \opt{hiddenstyle=ld.Hiddenlinestyle} : chaîne de caractères définissant le style de ligne des arêtes cachées, par défaut cette option contient la valeur de la variable globale \varglob{ld.Hiddenlinestyle} (qui vaut \val{"dotted"} par défaut).
        \item \opt{hiddencolor=color} : chaîne de caractères définissant la couleur des arêtes cachées, par défaut cette option contient la même couleur que l'option \opt{color}.
        \item \opt{style=<style courant>} : chaîne de caractères définissant le style de ligne des arêtes visibles.
        \item \opt{color=<couleur courante>} : chaîne de caractères définissant la couleur des arêtes visibles.
        \item \opt{width=<épaisseur courante>} : nombre représentant l'épaisseur de trait des arêtes (en dixième de point).
    \end{itemize}

    \item \textbf{Complément} : 
        \begin{itemize}
            \item La fonction \cmd{ld.facetedges(F)} où \argu{F} est une liste de facettes ou bien un polyèdre, renvoie une liste de segments 3D représentant toutes les arêtes de \argu{F}. Le résultat n'est pas un objet de type "edge", et il se dessine avec la méthode \cmd{g:Dpolyline3d()}. Bien sûr, chaque arête n'apparaît qu'une seule fois dans la liste.
            \item La fonction \cmd{ld.facetvertices(F)} où \argu{F} est une liste de facettes ou bien un polyèdre, renvoie la liste de tous les sommets de \argu{F} (points 3D).
        \end{itemize}
\end{itemize}

\subsection{Méthodes et fonctions s'appliquant à des facettes ou polyèdres}

\begin{itemize}
    \item La méthode \cmd{g:Isvisible(F)} où \argu{F} désigne \textbf{une} facette (liste d'au moins 3 points 3D coplanaires et non alignés), renvoie \true si la facette \argu{F} est visible (vecteur normal dirigé vers l'observateur). Si $A$, $B$ et $C$ sont les trois premiers points de \argu{F}, le vecteur normal est calculé en faisant le produit vectoriel $\vec{AB}\wedge\vec{AC}$.
    
    \item La méthode \cmd{g:Classifyfacet(F)} où \argu{F} est une liste de facettes ou bien un polyèdre, renvoie \textbf{deux} listes de facettes, la première est la liste des facettes visibles, et la suivante, la liste des facettes non visibles.
    
    \item La méthode \cmd{g:Sortfacet(F \fac{, backcull})} où \argu{F} est une liste de facettes, renvoie cette liste de facettes triées de la plus éloignée à la plus proche de l'observateur. L'argument facultatif \argu{backcull} est un booléen qui vaut \false par défaut, lorsqu'il a la valeur \true, les facettes non visibles sont exclues du résultat (seules les facettes visibles sont alors renvoyées après avoir été triées). Le calcul de l'éloignement d'un facette se fait sur son centre de gravité. La technique dite du "peintre" consiste à afficher les facettes de la plus éloignée à la plus proche, donc dans l'ordre de la liste renvoyée par cette fonction (le résultat affiché n'est cependant pas toujours correct en fonction de la taille et de la forme des facettes).
    
    \item La méthode \cmd{g:Sortpolyfacet(P \fac{, backcull})} où \argu{P} est un polyèdre, renvoie la liste des facettes de \argu{P} (facettes avec points 3D) triées de la plus éloignée à la plus proche de l'observateur. L'argument facultatif \argu{backcull} est un booléen qui vaut \false par défaut, lorsqu'il a la valeur \true, les facettes non visibles sont exclues du résultat comme pour la méthode précédente. Ces deux méthodes de tris sont utilisées par les méthodes de dessin de polyèdres ou facettes (\cmd{g:Dpoly()}, \cmd{g:Dfacet()} et \cmd{g:Dmixfacet()}).
    
    \item La méthode \cmd{g:Outline(P)} où \argu{P} est un polyèdre, renvoie le "contour" de \argu{P} sous la forme d'une table à deux champs, un champ nommé \emph{visible} qui contient une ligne polygonale 3D représentant les "arêtes" (segments) appartenant à une seule facette, celle-ci étant visible, ou bien à deux facettes, une visible et une cachée; l'autre champ est nommé \emph{hidden} et contient une ligne polygonale 3D représentant les "arêtes" appartenant à une seule facette, celle-ci étant cachée.
    
    \item La fonction \cmd{ld.border(P)} où \argu{P} est un polyèdre ou une liste de facette, renvoie une ligne polygonale 3D qui correspond aux arêtes appartenant à une seule facette de \argu{P} (ces arêtes sont mises "bout à bout" pour former une ligne polygonale).
    
    \item La fonction \cmd{ld.getfacet(P \fac{, list})} où \argu{P} est un polyèdre, renvoie la liste des facettes de \argu{P} (avec points 3D) dont le numéro figure dans la table \argu{list}. Si l'argument \argu{list} n'est pas précisé, c'est la liste de toutes les facettes de \argu{P} qui est renvoyée (dans ce cas c'est la même chose que \cmd{ld.poly2facet(P)}).
    
    \item La fonction \cmd{ld.facet2plane(L)} où \argu{L} est soit une facette, soit une liste de facettes, renvoie soit le plan contenant la facette, soit la liste des plans contenant chacune des facettes de \argu{L}. Un plan est une table du type \{A,u\} où $A$ est un point du plan et $u$ un vecteur normal au plan (donc deux points 3D).
    
    \item La fonction \cmd{ld.reverse\_face\_orientation(F)} où \argu{F} et soit une facette, soit une liste de facette, soit un polyèdre, renvoie un résultat de même nature que \argu{F} mais dans lequel l'ordre sur les sommets de chaque facette a été inversé. Cela peut être utile lorsque l'orientation de l'espace à été modifiée.
    
\begin{demo}{Sphère inscrite dans un octaèdre avec projection du centre sur les faces}
\begin{luadraw}{name=sphere_octaedre}
local ld = luadraw
local pt3d = ld.pt3d
local Origin, vecI, vecJ, vecK, M = pt3d.Origin, pt3d.vecI, pt3d.vecJ, pt3d.vecK, pt3d.M

local poly = require "luadraw_polyhedrons"
local g = ld.graph3d:new{ window3d={-3,3,-3,3,-3,3}, size={10,10}}
local P = poly.octahedron(Origin,M(0,0,3)) -- polyèdre défini dans le module luadraw_polyhedrons
P = ld.rotate3d(P,-10,{Origin,vecK}) -- rotate3d sur un polyèdre renvoie un polyèdre
local V, H = g:Classifyfacet(P) -- V pour facettes visibles, H pour hidden
local S = ld.map(function(p) return {ld.proj3d(Origin,p),p[2]} end, ld.facet2plane(V) )
-- S contient la liste de : {projeté, vecteur normal} (projetés de Origin sur les faces visibles)
local R = pt3d.abs(S[1][1]) -- rayon de la sphère
g:Dboxaxes3d({grid=true, gridcolor="gray", fillcolor="LightGray"})
g:Dfacet(H, {color="blue",opacity=0.9}) -- dessin des facettes non visibles
g:Dsphere(Origin,R,{mode=ld.mBorder,color="orange"}) -- dessin de la sphère
g:Dballdots3d(Origin,"gray",0.75) -- centre de la sphère
for _,D in ipairs(S) do -- segments reliant l'origine aux projetés
    g:Dpolyline3d( {Origin,D[1]},"dashed,gray")
end
g:Dfacet(V,{opacity=0.4, color="LightBlue"}) -- facettes visibles de l'octaèdre
g:Dcrossdots3d(S,nil,0.75) -- dessin des projetés sur les faces
g:Dpolyline3d( {M(0,-3,3), M(0,0,3), M(-3,0,3)},"gray")
g:Show()
\end{luadraw}
\end{demo}
\end{itemize}

\subsection{Découper un solide : cutpoly et cutfacet}

\begin{itemize}
    \item La fonction \cmd{ld.cutpoly(P, plane \fac{, close})} permet de couper le polyèdre \argu{P} avec le plan \argu{plane} (table du type \{A,n\} où $A$ est un point du plan et $n$ un vecteur normal au plan). La fonction renvoie 3 choses : la partie située dans le demi-espace contenant le vecteur $n$ (sous forme d'un polyèdre), suivie de la partie située dans l'autre demi-espace (toujours sous forme d'un polyèdre), suivie de la section sous forme d'une facette orientée par $-n$. Lorsque l'argument facultatif \argu{close} vaut \true, la section est ajoutée aux deux polyèdres résultants, ce qui a pour effet de les refermer (\false par défaut).\par
    \textbf{Remarque} : lorsque le polyèdre \argu{P} n'est pas convexe, le résultat de la section n'est pas toujours correct.

\begin{demo}{Cube coupé par un plan (cutpoly), avec \emph{close}=false et avec \emph{close}=true}
\begin{luadraw}{name=cutpoly}
local ld = luadraw
local pt3d = ld.pt3d
local Origin, vecI, vecJ, vecK, M = pt3d.Origin, pt3d.vecI, pt3d.vecJ, pt3d.vecK, pt3d.M

local g = ld.graph3d:new{window3d={-3,3,-3,3,-3,3}, window={-4,4,-3,3},size={10,10}}
local P = ld.parallelep(M(-1,-1,-1),2*vecI,2*vecJ,2*vecK)
local A, B, C = M(0,-1,1), M(0,1,1), M(1,-1,0)
local plane = {A, pt3d.prod(B-A,C-A)}
local P1 = ld.cutpoly(P,plane)
local P2 = ld.cutpoly(P,plane,true)
g:Lineoptions(nil,"Gold",8)
g:Dpoly( ld.shift3d(P1,-2*vecJ), {color="Crimson",mode=ld.mShadedHidden} )
g:Dpoly( ld.shift3d(P2,2*vecJ), {color="Crimson",mode=ld.mShadedHidden} )
g:Dlabel3d(
    "close=false", M(2,-2,-1), {dir={vecJ,vecK}},
    "close=true", M(2,2,-1), {} )
g:Show()
\end{luadraw}
\end{demo}

     \item La fonction \cmd{ld.cutfacet(F, plane \fac{, close})}, où \argu{F} est une facette, une liste de facettes, ou un polyèdre, fait la même chose que la fonction précédente sauf que cette fonction renvoie des listes de facettes et non pas des polyèdres. Cette fonction a été utilisée dans l'exemple des courbes de niveau à la figure \ref{courbeniv}.
\end{itemize}

\subsection{Clipper des facettes avec un polyèdre convexe : clip3d}

La fonction \cmd{ld.clip3d(S, P \fac{, exterior})} clippe le solide \argu{S} (liste de facettes ou bien polyèdre) avec le solide convexe \argu{P} (liste de facettes ou bien polyèdre) et renvoie la liste de facettes qui en résulte. L'argument facultatif \argu{exterior} est un booléen qui vaut \false par défaut, dans ce cas c'est la partie de \argu{S} qui est intérieure à \argu{P} qui est renvoyée, sinon c'est la partie de \argu{S} extérieure à \argu{P} qui est renvoyée.

\textbf{Remarque} : le résultat n'est pas toujours satisfaisant pour la partie extérieure.

\paragraph{Cas particulier} : clipper une liste de facettes $S$ (ou bien un polyèdre) avec la fenêtre 3D courante peut se faire avec cette fonction de la manière suivante :

\begin{center}
\textbf{S = ld.clip3d(S, g:Box3d())}
\end{center}

En effet, la méthode \cmd{g:Box3d()} renvoie la fenêtre 3D courante sous forme d'un parallélépipède.

\begin{demo}[clip3d]{Exemple avec clip3d : construction d'un dé à partir d'un cube et d'une sphère}
\begin{luadraw}{name=clip3d}
local ld = luadraw
local pt3d = ld.pt3d
local Origin, vecI, vecJ, vecK, M = pt3d.Origin, pt3d.vecI, pt3d.vecJ, pt3d.vecK, pt3d.M

local g = ld.graph3d:new{window={-3,3,-3,3},size={10,10}, viewdir="central"}
local S = ld.sphere(Origin,3)
local C = ld.parallelep(M(-2,-2,-2),4*vecI,4*vecJ,4*vecK)
local C1 = ld.clip3d(S,C) -- sphère clippée par le cube
local C2 = ld.clip3d(C,S) -- cube clippé par la sphère
local V = g:Classifyfacet(C2) -- facettes visibles de C2
g:Dfacet( ld.concat(C1,C2), {color="Beige",mode=ld.mShadedOnly,backcull=true} ) -- que les faces visibles
g:Dpolyline3d(V,true,"line width=0.8pt") -- contour des faces visibles de C2
local A, B, C, D = M(2,-2,-2), M(2,2,2), M(-2,2,-2), M(0,0,2) -- dessin des points noirs
g:Filloptions("full","black")
g:Dcircle3d( D,0.25,vecK); g:Dcircle3d( (2*A+B)/3,0.25,vecI)
g:Dcircle3d( (A+2*B)/3,0.25,vecI); g:Dcircle3d( (3*B+C)/4,0.25,vecJ)
g:Dcircle3d( (B+C)/2,0.25,vecJ); g:Dcircle3d( (B+3*C)/4,0.25,vecJ)
g:Show()
\end{luadraw}
\end{demo}

\subsection{Clipper un plan avec un polyèdre convexe : clipplane}

La fonction \cmd{ld.clipplane(plane, P)}, où l'argument \argu{plane} est une table de la forme \emph{\{A,n\}} représentant le plan passant par $A$ (point 3D) et de vecteur normal $n$ (point 3D non nul), et \argu{P} est un polyèdre convexe, renvoie la section du polyèdre par le plan, si elle existe, sous forme d'une facette (liste de points 3D) orientée par $n$.
%
\input{body-fr/luadraw-doc3d-section-5-La-methode-Dscene3d.tex}%
\section{Constructions géométriques}

Dans cette section sont regroupées les fonctions construisant des figures géométriques sans méthode graphique dédiée.

\subsection{Cercle circonscrit, cercle inscrit : circumcircle3d(), incircle3d()}

\begin{itemize}
    \item La fonction \cmd{ld.circumcircle3d(A, B, C)}, où \argu{A}, \argu{B} et \argu{C} sont trois points 3D non alignés, renvoie le cercle circonscrit au triangle formé par ces trois points, sous la forme d'une séquence: $O,R,n$, où $O$ est le centre du cercle, $R$ son rayon, et $n$ un vecteur normal au plan du cercle.
    \item La fonction \cmd{ld.incircle3d(A, B, C)}, où \argu{A}, \argu{B} et \argu{C} sont trois points 3D non alignés, renvoie le cercle inscrit dans le triangle formé par ces trois points, sous la forme d'une séquence: $I,R,n$, où $I$ est le centre du cercle, $R$ son rayon, et $n$ un vecteur normal au plan du cercle.    
\end{itemize}

\subsection{Enveloppe convexe : cvx\_hull3d()}

La fonction \cmd{ld.cvx\_hull3d(L)} où \argu{L} est une liste de points 3D \textbf{distincts}, calcule et renvoie l'enveloppe convexe de \argu{L} sous la forme d'une liste de facettes.

\begin{demo}{Utilisation de cvx\_hull3d()}
\begin{luadraw}{name=cvx_hull3d}
local ld = luadraw
local pt3d = ld.pt3d
local Origin, vecI, vecJ, vecK, M = pt3d.Origin, pt3d.vecI, pt3d.vecJ, pt3d.vecK, pt3d.M

local g = ld.graph3d:new{window={-2,4,-6,1},bbox=false,size={10,10}}
local L = {Origin, 4*vecI, M(4,4,0), 4*vecJ}
ld.insert(L, ld.shift3d(L,-3*vecK))
ld.insert(L, {M(2,1,2), M(2,3,2)})
local V = ld.cvx_hull3d(L)
local P = ld.facet2poly(V)
g:Dpoly(P , {color="cyan",mode=ld.mShadedHidden})
g:Show()
\end{luadraw}
\end{demo}

\paragraph{Cas particulier}: lorsque tous les points de \argu{L} sont dans un même plan, on peut utiliser la fonction \cmd{ld.cvx\_hull3dcoplanar(L, n)} où \argu{n} est un vecteur orthogonal au plan. Cette fonction renvoie une facette (liste de points 3D).

\subsection{Plans : plane(), planeEq(), orthoframe(), plane2ABC()}

Un plan de l'espace est une table de la forme $\{A,n\}$ où $A$ est un point du plan (point 3D) et $n$ un vecteur normal au plan (point 3D non nul).
\begin{itemize}
    \item La fonction \cmd{ld.plane(A, B, C)} envoie le plan passant par les trois points 3D \argu{A}, \argu{B} et \argu{C} (s'ils sont non alignés, sinon le résultat est \nil).
    
    \item La fonction \cmd{ld.planeEq(a, b, c, d)} envoie le plan dont une équation cartésienne est $ax+by+cz+d=0$ (si les coefficients \argu{a}, \argu{a} et \argu{a} ne sont pas tous nuls, sinon le résultat est\nil).
    
    \item La fonction \cmd{ld.plane2ABC(P)} où \argu{P}$=\{A,n\}$ désigne un plan, renvoie une séquence de trois points 3D $A,B,C$, appartenant au plan, et tels que $(A,\vec{AB},\vec{AC})$ soit un repère orthonormal direct de ce plan.
    
    \item La fonction \cmd{ld.orthoframe(P)} où \argu{P}$=\{A,n\}$ désigne un plan, renvoie une séquence de trois points 3D $A,u,v$, tels que $(A,u,v)$ soit un repère orthonormal direct de ce plan.
\end{itemize}

\begin{demo}{Faces d'un cube trouées avec un hexagone régulier}
\begin{luadraw}{name=plans}
local ld = luadraw
local pt3d = ld.pt3d
local Origin, vecI, vecJ, vecK, M = pt3d.Origin, pt3d.vecI, pt3d.vecJ, pt3d.vecK, pt3d.M

local g = ld.graph3d:new{window={-3,3,-3.25,3.25}, margin={0,0,0,0},
    viewdir={"central",20,60}, bg="LightGray", size={10,10}}
ld.Hiddenlines = true; ld.Hiddenlinestyle = "dashed"
local p = ld.polyreg(0,1,6)
local P = ld.parallelep(M(-2,-2,-2),4*vecI,4*vecJ,4*vecK)
local V = g:Sortpolyfacet(P)
local list = {}
g:Filloptions("full","Crimson",1,true); -- true pour le mode evenodd
g:Lineoptions("solid","Gold",8)
for _, F in ipairs(V) do
    local P1 = ld.plane(pt3d.isobar3d(F),F[1],F[2]) -- plan de la facette F
    local A, u, v = ld.orthoframe(P1) -- repère orthonormé sur la facette avec centre de gravité comme origine
    local p1 = ld.map(function(z) return A+z.re*u+z.im*v end,p) -- hexagone reproduit sur la facette
    table.insert(p1,2,"m")
    local color = "Crimson"
    if not g:Isvisible(F) then color = "Crimson!60!black" end
    g:Dpath3d( ld.concat(F,{"l"},p1,{"l","cl"}),"fill="..color ) -- dessin de la facette "trouée" avec l'hexagone
end
g:Show()
\end{luadraw}
\end{demo}

\subsection{Sphère circonscrite, Sphère inscrite : circumsphere(), insphere()}

\begin{itemize}
    \item La fonction \cmd{ld.circumsphere(A, B, C, D)}, où \argu{A}, \argu{B}, \argu{C} et \argu{D} sont quatre points 3D non coplanaires, renvoie la sphère circonscrite au tétraèdre formé par ces quatre points, sous la forme d'une séquence: $I,R$, où $I$ est le centre de la sphère, et $R$ son rayon.
    
    \item La fonction \cmd{ld.insphere(A, B, C, D)}, où \argu{A}, \argu{B}, \argu{C} et \argu{D} sont quatre points 3D non coplanaires, renvoie la sphère inscrite dans le tétraèdre formé par ces quatre points, sous la forme d'une séquence: $I,R$, où $I$ est le centre de la sphère, et $R$ son rayon.
\end{itemize}

\subsection{Tétraèdre à longueurs fixées : tetra\_len()}

La fonction \cmd{ld.tetra\_len(ab, ac, ad, bc, bd, cd)} calcule les sommets $A,B,C,D$ d'un tétraèdre dont les longueurs des arêtes sont données, c'est à dire tels que $AB=$\argu{ab}, $AC=$\argu{ac}, $AD=$\argu{ad}, $BC=$\argu{bc}, $BD=$\argu{bd} et $CD=$\argu{cd}. La fonction renvoie la séquence de quatre points $A,B,C,D$. Le sommet $A$ est toujours le point $M(0,0,0)$ (\varglob{pt3d.Origin}) et le sommet $B$ est toujours le point \code{ab*pt3d.vecI} et le sommet $C$ dans le plan $xy$. Le tétraèdre en tant que polyèdre peut ensuite être construit avec la fonction \cmd{ld.tetra(A, B-A, C-A, D-A)}.

\begin{demo}{Un tétraèdre avec la longueur des arêtes fixée}
\begin{luadraw}{name=tetra_len}
local ld = luadraw
local g = ld.graph3d:new{window={-4,4,-4,4},margin={0,0,0,0},viewdir={25,65},size={10,10}}
ld.Hiddenlines = true; ld.Hiddenlinestyle = "dashed"
require 'luadraw_spherical'
local R = 4
local A,B,C,D = ld.tetra_len(R,R,R,R,R,R)
local T = ld.tetra(A,B-A,C-A,D-A)
g:Define_sphere({radius=R})
g:DSpolyline( ld.facetedges(T), {color="DarkGreen"})
g:DSbigcircle( {B,C},{color="Blue"} )
g:DSbigcircle( {B,D},{color="Blue"} )
g:DSbigcircle( {C,D},{color="Blue"}  )
g:DSlabel("$R$",(2*A+C)/3,{pos="S"})
g:Dspherical()
g:Ddots3d({A,B,C,D})
g:Dlabel3d("$A$",A,{pos="S"},"$B$",B,{pos="SW"},"$C$",C,{},"$D$",D,{pos="N"} )
g:Show()
\end{luadraw}
\end{demo}

\subsection{Triangles : sss\_triangle3d(), sas\_triangle3d(), asa\_triangle3d()}

Ces fonctions sont la version 3D des fonctions  sss\_triangle(), sas\_triangle(), asa\_triangle() déjà décrites.
\begin{itemize}
    \item La fonction \cmd{ld.sss\_triangle3d(ab, bc, ca)} où \argu{ab}, \argu{bc} et \argu{ca} sont trois longueurs, calcule et renvoie une liste de trois points 3D $\{A,B,C\}$ formant les sommets d'un triangle direct dans le plan $xOy$ dont les longueurs des côtés sont les arguments, c'est à dire $AB=$\argu{ab}, $BC=$\argu{bc} et $CA=$\argu{ca}, lorsque cela est possible. Le sommet $A$ est toujours le point $M(0,0,0)$ (\varglob{pt3d.Origin}) et le sommet $B$ est toujours le point \code{ab*pt3d.vecI}. Ce triangle peut être dessiné avec la méthode \cmd{g:Dpolyline3d()}.
    
    \item La fonction \cmd{ld.sas\_triangle3d(ab, alpha, ca)} où \argu{ab} et \argu{ca} sont deux longueurs, \argu{alpha} un angle en degrés, calcule et renvoie une liste de trois points 3D $\{A,B,C\}$ formant les sommets d'un triangle dans le plan $xOy$ tel que $AB=$\argu{ab}, $CA=$\argu{ca}, et tel que l'angle $(\vec{AB},\vec{AC})$ a pour mesure \argu{alpha}, lorsque cela est possible. Le sommet $A$ est toujours le point $M(0,0,0)$ (\varglob{pt3d.Origin}) et le sommet $B$ est toujours le point \code{ab*pt3d.vecI}. Ce triangle peut être dessiné avec la méthode \cmd{g:Dpolyline3d()}.
    
    \item La fonction \cmd{ld.asa\_triangle3d(alpha, ab, beta)} où \argu{ab} est une longueur, \argu{alpha} et \argu{beta} deux angles en degrés, calcule et renvoie une liste de trois points 3D $\{A,B,C\}$ formant les sommets d'un triangle dans le plan $xOy$ tel que $AB=$\argu{ab}, tel que l'angle $(\vec{AB},\vec{AC})$ a pour mesure \argu{alpha}, et tel que l'angle $(\vec{BA},\vec{BC})$ a pour mesure \argu{beta}, lorsque cela est possible. Le sommet $A$ est toujours le point $M(0,0,0)$ (\varglob{pt3d.Origin}) et le sommet $B$ est toujours le point \code{ab*pt3d.vecI}. Ce triangle peut être dessiné avec la méthode \cmd{g:Dpolyline3d()}.
\end{itemize}
%
\input{body-fr/luadraw-doc3d-section-7-Transformations-calcul-matriciel-fonctions-mathematiques.tex}%
\input{body-fr/luadraw-doc3d-section-8-Exemples-plus-pousses.tex}%

\include{body-fr/luadraw-doc-Annexes.tex}


\end{document}
